\documentclass[11pt, a4paper]{article}

%---------------------------------------------------------
% Lingua E Codifica
%---------------------------------------------------------
\usepackage[italian]{babel}
\usepackage[T1]{fontenc}
\usepackage{lmodern}

%---------------------------------------------------------
% Layout Di Pagina E Tipografia
%---------------------------------------------------------
\usepackage[margin=2cm, top=1.5cm, bottom=1.5cm]{geometry}
\usepackage{parskip}
\usepackage{microtype}
\usepackage{csquotes}

%---------------------------------------------------------
% Gestione Dei Colori
%---------------------------------------------------------
\usepackage[table]{xcolor}
\definecolor{darkgreen}{RGB}{0,100,0}
\usepackage{soul}
%---------------------------------------------------------
% Matematica E Simboli
%---------------------------------------------------------
\usepackage{amsmath, amssymb, amsfonts}
\usepackage{mathtools}
\usepackage{siunitx}

% BB OTTIMIZZAZIONE MATRICI LUNGHE PER LA STAMPA
\setlength{\arraycolsep}{10pt}
\usepackage{bm}
\usepackage{arydshln}

\DeclareMathOperator*{\argmax}{arg\,max}
\DeclareMathOperator*{\argmin}{arg\,min} % Utile se si usa anche l'argmin

%---------------------------------------------------------
% Grafica, Immagini E Tabelle
%---------------------------------------------------------
\usepackage{graphicx}
\graphicspath{{./Images/Convolutional Neural Networks/}}
\usepackage{float}
\usepackage{booktabs}
\usepackage{array}
\usepackage{multirow}
\usepackage{tabularx}
\usepackage{listings}
\usepackage{tikz}
\usetikzlibrary{arrows.meta, patterns, angles, quotes, calc}
\usepackage{pgfplots}
\pgfplotsset{compat=1.18}


%---------------------------------------------------------
% Link, Riferimenti Ed Elenchi
%---------------------------------------------------------
\usepackage[unicode]{hyperref}
\hypersetup{
    colorlinks=true,
    linkcolor=blue,
    urlcolor=blue,
    citecolor=blue,
    pdftitle={Capitolo 11 - Convolutional Neural Networks},
    pdfauthor={Francesco Lo Verde},
    pdfpagemode=FullScreen,
}
\usepackage{bookmark}
\usepackage{enumitem}
\setlist{noitemsep}

%---------------------------------------------------------
% Personalizzazione Titoli E Altri Pacchetti
%---------------------------------------------------------
\usepackage{titlesec}
\titleformat{\section}{\normalfont\LARGE\bfseries}{\thesection}{1em}{}
\titleformat{\subsection}{\normalfont\Large\bfseries}{\thesubsection}{1em}{}
\titleformat{\subsubsection}{\normalfont\large\bfseries}{\thesubsubsection}{1em}{}

\usepackage{comment}
\newtheorem{theorem}{Definizione}

% Pacchetto per gli algoritmi
\usepackage[ruled,vlined,linesnumbered]{algorithm2e}

\SetKwInput{Dati}{Input}
\SetKwInput{Risultato}{Output}
\SetKw{Inizializza}{Inizializza}
\SetKwIF{Se}{AltrimentiSe}{Altrimenti}{se}{allora}{altrimenti se}{altrimenti}{finese}
\SetKwProg{Fn}{Funzione}{:}{fine}
\SetKwFor{Per}{per}{fai}{fine per}
\SetKwRepeat{Ripeti}{ripeti}{finché}
\SetKwComment{Commento}{// }{}

\author{Francesco Lo Verde}
\begin{document}
\section{Introduzione Al Deep Learning E Reti Feedforward}

Le reti neurali artificiali rappresentano la spina dorsale del moderno apprendimento automatico. Per comprendere appieno il loro funzionamento, è essenziale analizzare come l'informazione fluisca al loro interno e come la loro struttura influenzi la capacità di approssimare funzioni complesse.

% [IMMAGINE RIF FILE NOMIMATO "TUTTO"] Confronto strutturale tra una rete neurale semplice (Simple Neural Network) con un solo strato nascosto e una rete neurale profonda (Deep Learning Neural Network) con quattro strati nascosti — Slide 1 / Pagina 1 del PDF

\subsection{La Struttura Di Una Rete Feedforward}

Le \textbf{reti neurali artificiali feedforward} (note anche come \emph{feedforward neural networks} o \textbf{percettroni multistrato}, abbreviati in \textbf{MLP}) rappresentano i modelli di \emph{deep learning} per eccellenza \cite{goodfellow2016deep}. L'obiettivo primario di una rete feedforward è quello di approssimare una qualche funzione obiettivo incognita, indicata con $f^*(\bm{x})$. Ad esempio, nel caso specifico di un classificatore, la funzione ideale $y = f^*(\bm{x})$ si occupa di mappare un vettore di ingresso $\bm{x}$ a una determinata categoria o etichetta $y$.

Una rete feedforward definisce a sua volta una mappatura matematica parametrizzata del tipo:
\[
\bm{y} = f(\bm{x}; \bm{\theta})
\]
Il processo di addestramento si focalizza sull'apprendimento del valore ottimale dei parametri $\bm{\theta}$ che permetta di ottenere la migliore approssimazione possibile della funzione target $f^*(\bm{x})$.

Il termine \textbf{feedforward} (ovvero a propagazione in avanti) deriva direttamente dal modo in cui l'informazione fluisce all'interno del modello: essa parte dall'ingresso $\bm{x}$, attraversa i calcoli intermedi necessari per definire la funzione $f$, e giunge infine all'uscita $\bm{y}$. In questa tipologia strutturale non sono presenti connessioni di retroazione (\emph{feedback}), vale a dire collegamenti in cui le uscite del modello vengono reintrodotte come ingressi al modello stesso. Qualora la struttura venga estesa per includere tali connessioni di retroazione, si entra nel dominio delle \textbf{reti neurali ricorrenti} (\emph{Recurrent Neural Networks}, o RNN).

Dal punto di vista matematico, le reti feedforward possono essere interpretate come una composizione di molteplici funzioni elementari strutturate a catena. Consideriamo, a titolo esemplificativo, una catena formata da tre funzioni collegate in serie:
\[
f(\bm{x}) = f^{(1)}\left(f^{(2)}\left(f^{(3)}(\bm{x})\right)\right)
\]
In questa formulazione, la funzione $f^{(1)}$ viene associata al primo strato (\emph{layer}) della rete, $f^{(2)}$ rappresenta il secondo strato, e così via procedendo lungo la catena. Il numero complessivo di funzioni (ovvero di strati) che compongono il modello definisce la \textbf{profondità} (\emph{depth}) della rete neurale; da questa caratteristica deriva la denominazione stessa di \textbf{apprendimento profondo} (\emph{deep learning}). L'ultimo strato della catena prende il nome di \textbf{strato di uscita} (\emph{output layer}).

Durante la fase di calcolo e ottimizzazione, se la vera funzione da stimare è $f^*(\bm{x})$, l'intero processo di addestramento mira a trovare una funzione $f(\bm{x})$ che si avvicini il più possibile a essa. Nello specifico, dato un campione di addestramento $\bm{x}$ associato a un'etichetta reale $y = f^*(\bm{x})$, l'algoritmo richiede che lo strato di uscita generi un valore numerico estremamente prossimo a $y$.

Gli strati intermedi della rete, invece, non hanno un valore di uscita desiderato specificato a priori dai dati. Proprio per questo motivo essi vengono definiti \textbf{strati nascosti} (\emph{hidden layers}). Il compito della rete consiste nell'utilizzare autonomamente questi strati intermedi per elaborare l'informazione in modo da produrre l'output finale corretto.

Ogni strato nascosto è costituito da un insieme di elementi computazionali elementari detti \textbf{neuroni} o \textbf{unità}; il numero di tali unità determina la \textbf{larghezza} (\emph{width}) dello strato e, di riflesso, del modello stesso. Piuttosto che concepire uno strato nascosto esclusivamente come una generica funzione vettoriale, risulta più intuitivo immaginarlo come un insieme di unità operanti in parallelo. Ciascuna unità riceve in ingresso le attivazioni provenienti da numerose altre unità dello strato precedente e produce il proprio valore di uscita applicando una specifica \textbf{funzione di attivazione} non lineare alla combinazione lineare dei suoi ingressi.

\newpage

\section{Perché Il Deep Learning: Il Ruolo Della Rappresentazione Non Lineare}

I modelli lineari tradizionali, come la regressione lineare o la regressione logistica, possiedono proprietà indubbiamente interessanti: possono essere addestrati in modo estremamente efficiente e affidabile, sia tramite formule in forma chiusa sia sfruttando le tecniche classiche dell'ottimizzazione convessa. Tuttavia, essi presentano una limitazione intrinseca invalicabile: la loro capacità espressiva è strettamente limitata alle sole funzioni lineari. Di conseguenza, un modello lineare non è in grado di comprendere o catturare l'interazione complessa che può sussistere tra due o più variabili di ingresso.

\subsection{I Limiti Dei Modelli Lineari}

Per superare questo limite e consentire ai modelli lineari di rappresentare funzioni fortemente non lineari di un input $\bm{x}$, si può scegliere di non applicare il modello direttamente a $\bm{x}$, bensì a una sua versione trasformata $\bm{\varphi}(\bm{x})$, dove $\bm{\varphi}$ rappresenta una trasformazione non lineare.

In maniera del tutto equivalente, è possibile ricorrere al celebre \textbf{trucco del kernel} (\emph{kernel trick}) per ottenere un algoritmo di apprendimento non lineare basato sull'applicazione implicita della mappatura $\bm{\varphi}$. Possiamo interpretare $\bm{\varphi}$ come una funzione che fornisce un insieme di caratteristiche (\emph{features}) utili a descrivere $\bm{x}$, oppure come una vera e propria nuova rappresentazione spaziale di $\bm{x}$.

La questione fondamentale risiede nel modo in cui scegliere questa trasformazione $\bm{\varphi}$:
\begin{itemize}[noitemsep]
    \item \textbf{Usare una funzione $\bm{\varphi}$ molto generica}: approccio tipico dei metodi a kernel robusti, che tuttavia spesso non generalizzano in modo efficiente su problemi ad altissima dimensionalità.
    \item \textbf{Progettare manualmente $\bm{\varphi}$}: pratica comune nel machine learning classico (\emph{feature engineering}), laboriosa e strettamente dipendente dal singolo dominio applicativo.
    \item \textbf{Apprendere la funzione $\bm{\varphi}$}: questo rappresenta il nucleo concettuale e l'obiettivo fondamentale del \emph{deep learning}. Il modello apprende autonomamente una rappresentazione dei dati che rende il problema linearmente risolvibile negli strati successivi.
\end{itemize}

\subsection{L'Esempio Del Problema XOR}

Per comprendere concretamente la necessità di introdurre non linearità e strati nascosti, analizziamo il classico problema della funzione logica \textbf{XOR} (OR esclusivo). Si consideri un piccolo insieme di addestramento definito dalla matrice degli ingressi:
\[
\bm{X} = \begin{bmatrix} 0 & 0 \\ 0 & 1 \\ 1 & 0 \\ 1 & 1 \end{bmatrix}
\]
L'obiettivo consiste unicamente nel far sì che il modello si adatti perfettamente a questo specifico training set, senza considerare per il momento problematiche legate alla generalizzazione su un test set esterno. Formuliamo il problema come un compito di regressione che punta a minimizzare l'errore quadratico medio (\emph{Mean Squared Error}, MSE) sul dataset di addestramento:
\[
J(\bm{\theta}) = \frac{1}{4} \sum_{\bm{x} \in \mathcal{X}} \left( f^*(\bm{x}) - f(\bm{x}; \bm{\theta}) \right)^2
\]
Proviamo inizialmente a risolvere questo problema apprendendo la funzione XOR per mezzo di un modello puramente lineare, definito dall'equazione:
\[
f(\bm{x}; \bm{w}, b) = \bm{x}^\top \bm{w} + b
\]

% [IMMAGINE RIF FILE NOMIMATO "TUTTO"] Rappresentazione grafica dei punti dello XOR nello spazio originale x. I punti (0,0) e (1,1) hanno etichetta 0, mentre i punti (0,1) e (1,0) hanno etichetta 1 — Slide 6 / Pagina 6 del PDF

Un modello lineare applicato direttamente allo spazio originario degli ingressi non può in alcun modo implementare la funzione logica XOR. Analizzando il comportamento dei punti, si osserva che:
\begin{itemize}[noitemsep]
    \item Quando l'ingresso $x_1 = 0$, l'output del modello deve necessariamente aumentare all'aumentare di $x_2$.
    \item Quando l'ingresso $x_1 = 1$, l'output del modello deve invece diminuire all'aumentare di $x_2$.
\end{itemize}
Tuttavia, un modello lineare è vincolato ad applicare un coefficiente fisso $w_2$ alla variabile $x_2$. Di conseguenza, il modello non ha modo di utilizzare il valore assunto da $x_1$ per modificare dinamicamente il coefficiente associato a $x_2$, rendendo il problema matematicamente impossibile da risolvere in modo lineare.

\subsection{Risoluzione Dello XOR Tramite Una Rete Feedforward Con Attivazione ReLU}

Per superare questa limitazione, introduciamo una semplice rete neurale feedforward dotata di un singolo strato nascosto contenente due unità.

% [IMMAGINE RIF FILE NOMIMATO "TUTTO"] Diagramma architetturale di una rete feedforward con un input vettoriale x, uno strato nascosto h composto da due unità (h1, h2) e un output scalare y — Slide 7 / Pagina 7 del PDF

Questa rete calcola un vettore di attivazioni nascoste $\bm{h}$ per mezzo di una funzione vettoriale $\bm{f}^{(1)}(\bm{x}; \bm{W}, \bm{c})$. I valori calcolati all'interno di questo strato nascosto vengono successivamente impiegati come input per il secondo strato, che rappresenta lo strato di uscita della rete. Quest'ultimo è a tutti gli effetti un classico modello di regressione lineare, applicato tuttavia alle rappresentazioni intermediate $\bm{h}$ piuttosto che agli ingressi originali $\bm{x}$.

La struttura complessiva prevede quindi due funzioni concatenate in serie:
\[
\bm{h} = \bm{f}^{(1)}(\bm{x}; \bm{W}, \bm{c}) \quad \text{e} \quad y = f^{(2)}(\bm{h}; \bm{w}, b)
\]
Il modello completo assume dunque la forma complessiva:
\[
f(\bm{x}; \bm{W}, \bm{c}, \bm{w}, b) = f^{(2)}\left(\bm{f}^{(1)}(\bm{x})\right)
\]
Quale forma matematica dovrebbe assumere la funzione del primo strato $\bm{f}^{(1)}$? Poiché i modelli lineari sono semplici ed efficienti, si potrebbe essere tentati di definire lineare anche $\bm{f}^{(1)}$. Sfortunatamente, se $\bm{f}^{(1)}$ fosse una trasformazione lineare, l'intera rete feedforward rimarrebbe una funzione lineare del suo input originale.

Per dimostrarlo matematicamente, ignoriamo momentaneamente i termini di bias (intercetta). Se poniamo $\bm{f}^{(1)}(\bm{x}) = \bm{W}^\top \bm{x}$ e $f^{(2)}(\bm{h}) = \bm{h}^\top \bm{w}$, la composizione delle funzioni produce:
\[
f(\bm{x}) = \left( \bm{W}^\top \bm{x} \right)^\top \bm{w} = \bm{x}^\top \bm{W} \bm{w}
\]
Possiamo ridefinire questa relazione come un'unica funzione lineare nella forma:
\[
f(\bm{x}) = \bm{x}^\top \bm{w}' \quad \text{dove} \quad \bm{w}' = \bm{W}\bm{w}
\]
Questo dimostra che la composizione di strati puramente lineari collassa inevitabilmente in un unico operatore lineare, preservando l'incapacità intrinseca di risolvere problemi non lineari come lo XOR. Diventa quindi indispensabile selezionare una funzione di attivazione non lineare. Scegliamo di utilizzare l'unità lineare rettificata, nota come \textbf{ReLU} (\emph{Rectified Linear Unit}):
\[
g(z) = \max\{0, z\}
\]

% [IMMAGINE RIF FILE NOMIMATO "TUTTO"] Grafico della funzione di attivazione non lineare ReLU g(z) = max{0, z} in funzione dell'input z — Slide 8 / Pagina 8 del PDF

Siamo ora in grado di specificare la formulazione analitica completa della nostra rete neurale:
\[
f(\bm{x}; \bm{W}, \bm{c}, \bm{w}, b) = \bm{w}^\top \max\left\{0, \bm{W}^\top \bm{x} + \bm{c}\right\} + b
\]
Mostriamo come questa formulazione permetta di risolvere il problema dello XOR. Definiamo i seguenti parametri per il modello:
\[
\bm{W} = \begin{bmatrix} 1 & 1 \\ 1 & 1 \end{bmatrix}, \quad \bm{c} = \begin{bmatrix} 0 \\ -1 \end{bmatrix}, \quad \bm{w} = \begin{bmatrix} 1 \\ -2 \end{bmatrix}, \quad b = 0
\]
Applichiamo passo dopo passo i calcoli per ciascuno dei quattro possibili vettori di ingresso $\bm{x} \in \mathcal{X}$:

\begin{enumerate}
    \item Per l'ingresso $\bm{x} = \begin{bmatrix} 0 & 0 \end{bmatrix}^\top$:
    \begin{align*}
    \bm{W}^\top \bm{x} + \bm{c} &= \begin{bmatrix} 1 & 1 \\ 1 & 1 \end{bmatrix} \begin{bmatrix} 0 \\ 0 \end{bmatrix} + \begin{bmatrix} 0 \\ -1 \end{bmatrix} = \begin{bmatrix} 0 \\ -1 \end{bmatrix} \\
    \bm{h} &= \max\left\{0, \begin{bmatrix} 0 \\ -1 \end{bmatrix}\right\} = \begin{bmatrix} 0 \\ 0 \end{bmatrix} \\
    y &= \bm{w}^\top \bm{h} + b = \begin{bmatrix} 1 & -2 \end{bmatrix} \begin{bmatrix} 0 \\ 0 \end{bmatrix} + 0 = \boxed{0}
    \end{align*}

    \item Per l'ingresso $\bm{x} = \begin{bmatrix} 0 & 1 \end{bmatrix}^\top$:
    \begin{align*}
    \bm{W}^\top \bm{x} + \bm{c} &= \begin{bmatrix} 1 & 1 \\ 1 & 1 \end{bmatrix} \begin{bmatrix} 0 \\ 1 \end{bmatrix} + \begin{bmatrix} 0 \\ -1 \end{bmatrix} = \begin{bmatrix} 1 \\ 1 \end{bmatrix} + \begin{bmatrix} 0 \\ -1 \end{bmatrix} = \begin{bmatrix} 1 \\ 0 \end{bmatrix} \\
    \bm{h} &= \max\left\{0, \begin{bmatrix} 1 \\ 0 \end{bmatrix}\right\} = \begin{bmatrix} 1 \\ 0 \end{bmatrix} \\
    y &= \bm{w}^\top \bm{h} + b = \begin{bmatrix} 1 & -2 \end{bmatrix} \begin{bmatrix} 1 \\ 0 \end{bmatrix} + 0 = \boxed{1}
    \end{align*}

    \item Per l'ingresso $\bm{x} = \begin{bmatrix} 1 & 0 \end{bmatrix}^\top$:
    \begin{align*}
    \bm{W}^\top \bm{x} + \bm{c} &= \begin{bmatrix} 1 & 1 \\ 1 & 1 \end{bmatrix} \begin{bmatrix} 1 \\ 0 \end{bmatrix} + \begin{bmatrix} 0 \\ -1 \end{bmatrix} = \begin{bmatrix} 1 \\ 1 \end{bmatrix} + \begin{bmatrix} 0 \\ -1 \end{bmatrix} = \begin{bmatrix} 1 \\ 0 \end{bmatrix} \\
    \bm{h} &= \max\left\{0, \begin{bmatrix} 1 \\ 0 \end{bmatrix}\right\} = \begin{bmatrix} 1 \\ 0 \end{bmatrix} \\
    y &= \bm{w}^\top \bm{h} + b = \begin{bmatrix} 1 & -2 \end{bmatrix} \begin{bmatrix} 1 \\ 0 \end{bmatrix} + 0 = \boxed{1}
    \end{align*}

    \item Per l'ingresso $\bm{x} = \begin{bmatrix} 1 & 1 \end{bmatrix}^\top$:
    \begin{align*}
    \bm{W}^\top \bm{x} + \bm{c} &= \begin{bmatrix} 1 & 1 \\ 1 & 1 \end{bmatrix} \begin{bmatrix} 1 \\ 1 \end{bmatrix} + \begin{bmatrix} 0 \\ -1 \end{bmatrix} = \begin{bmatrix} 2 \\ 2 \end{bmatrix} + \begin{bmatrix} 0 \\ -1 \end{bmatrix} = \begin{bmatrix} 2 \\ 1 \end{bmatrix} \\
    \bm{h} &= \max\left\{0, \begin{bmatrix} 2 \\ 1 \end{bmatrix}\right\} = \begin{bmatrix} 2 \\ 1 \end{bmatrix} \\
    y &= \bm{w}^\top \bm{h} + b = \begin{bmatrix} 1 & -2 \end{bmatrix} \begin{bmatrix} 2 \\ 1 \end{bmatrix} + 0 = 2 - 2 = \boxed{0}
    \end{align*}
\end{enumerate}

Grazie alla non linearità introdotta dalla funzione di attivazione ReLU, la rete proietta i dati in uno spazio latente trasformato $\bm{h}$ in cui le classi sono perfettamente separabili mediante un iperpiano lineare.

% [IMMAGINE RIF FILE NOMIMATO "TUTTO"] Rappresentazione grafica dello spazio trasformato h (Learned h space) in cui i punti del problema XOR sono diventati linearmente separabili — Slide 9 / Pagina 9 del PDF

\newpage

\section{Perché Andare Più In Profondità: Reti Deep Contro Reti Shallow}

L'analisi strutturale dei modelli basati su reti neurali impone una riflessione cruciale riguardante la loro geometria complessiva: per quale motivo risulta vantaggioso sviluppare architetture profonde composte da numerosi strati sovrapposti, anziché limitarsi a modelli con un singolo strato nascosto estremamente esteso?

\subsection{Il Teorema Di Approssimazione Universale}

Il fondamento teorico delle reti neurali è sancito dal celebre \textbf{Teorema di Approssimazione Universale} (formulato da Hornik et al. nel 1989 e Cybenko nel 1989). Tale teorema assicura che:
\begin{quote}
Una rete neurale di tipo feedforward dotata di uno strato di uscita lineare e di almeno uno strato nascosto contenente una qualsiasi funzione di attivazione non lineare (come ad esempio la funzione sigmoidea logistica) è in grado di approssimare con un livello di errore arbitrariamente piccolo una qualunque funzione continua definita su sottoinsiemi compatti dello spazio reale, a patto che lo strato nascosto disponga di un numero sufficientemente elevato di unità.
\end{quote}

Questo enunciato sembrerebbe suggerire che un singolo strato nascosto (architettura \emph{shallow}) sia teoricamente sufficiente per risolvere qualsiasi problema di approssimazione. Di conseguenza, sorge spontaneo domandarsi per quale motivo sia necessario ricorrere a reti con un numero di strati superiore.

La risposta risiede nei limiti pratici legati all'addestramento e alla complessità parametrica del modello:
\begin{itemize}
    \item \textbf{Mancanza di garanzie sull'apprendimento}: Il teorema sancisce l'esistenza matematica di una configurazione idonea di parametri, ma non fornisce alcuna garanzia sul fatto che l'algoritmo di ottimizzazione sia effettivamente in grado di individuarla durante l'addestramento. L'apprendimento può infatti fallire per due principali motivazioni:
    \begin{itemize}[noitemsep]
        \item L'algoritmo di ottimizzazione potrebbe non riuscire a convergere verso i valori ottimali dei parametri corrispondenti alla funzione desiderata.
        \item L'algoritmo di addestramento potrebbe selezionare una funzione errata a causa di fenomeni di \textbf{overfitting} (sovradattamento).
    \end{itemize}
    \item \textbf{Crescita esponenziale delle dimensioni}: Qualora la funzione da approssimare risulti complessa o presenti specifiche regolarità, l'approccio a singolo strato potrebbe richiedere un numero di unità nascoste che cresce in modo \textbf{esponenziale} rispetto alla dimensionalità dell'input.
    \item \textbf{Assenza di assunzioni a priori (priors)}: Uno strato singolo contenente un numero enorme di neuroni non impone alcuna struttura o conoscenza a priori sul modello, rendendo il processo di ottimizzazione numerica estremamente più arduo e inefficiente.
\end{itemize}

\subsection{I Vantaggi Delle Architetture Profonde}

Tutte le criticità sopra delineate vengono notevolmente mitigate adottando strutture profonde, ovvero aumentando il numero di strati e riducendo parallelamente il numero di neuroni interni a ciascuno strato (\emph{going deeper}):
\begin{itemize}[noitemsep]
    \item \textbf{Codifica di assunzioni a priori sulla struttura dei dati}: La scelta di un'architettura profonda equivale a integrare nel modello l'ipotesi forte che la funzione complessa che si desidera apprendere possa essere espressa come la composizione gerarchica di molteplici funzioni più semplici.
    \item \textbf{Efficienza parametrica}: L'impiego di modelli profondi consente di ridurre drasticamente il numero complessivo di unità necessarie per descrivere accuratamente la medesima funzione complessa, abbattendo contemporaneamente la varianza e l'errore di generalizzazione.
    \item \textbf{Ottimizzazione semplificata}: Una struttura profonda è tipicamente molto più semplice da addestrare rispetto a una rete caratterizzata da un singolo strato nascosto dotato di una quantità sproporzionata di neuroni paralleli.
\end{itemize}

\newpage

\section{Il Deep Learning Nella Visione Artificiale}

Le considerazioni svolte finora evidenziano l'efficacia del \emph{deep learning} per l'approssimazione di funzioni non lineari complesse. Ciononostante, è noto che compiti standard come la classificazione di immagini possono essere affrontati con discreto successo anche mediante algoritmi di machine learning classici più semplici, tra cui:
\begin{itemize}[noitemsep]
    \item L'algoritmo dei $K$-vicini più prossimi (\textbf{KNN})
    \item Le macchine a vettori di supporto (\textbf{SVM})
    \item Le reti neurali semplici dotate di un singolo strato nascosto
\end{itemize}

Diventa quindi lecito chiedersi quali siano i reali benefici introdotti da una rete profonda nel contesto specifico dei compiti di visione artificiale (\emph{Computer Vision}, CV), e in che modo essa riesca a migliorare sensibilmente le prestazioni dei modelli.

Per chiarire i vantaggi derivanti dalle architetture profonde in questo specifico ambito applicativo, si rende necessario analizzare preventivamente le modalità di funzionamento dei tradizionali algoritmi di machine learning, mettendone in luce i limiti strutturali. Successivamente, sarà fondamentale introdurre una specifica e cruciale modifica architetturale espressamente concepita per rendere le reti profonde idonee ed efficienti per l'elaborazione di dati bidimensionali complessi come le immagini.

% AA Slides 11.1
\section{Introduzione Al Deep Learning E Alle Reti Neurali Convoluzionali}

% [IMMAGINE DEL FILE DI SLIDE CON NOME "1-28"] Slide di copertina del corso "Convolutional Neural Networks" (CMSC 733, Fall 2015) di Angjoo Kanazawa — Slide 1 / Pagina 1 del PDF

Il presente documento raccoglie gli appunti relativi alle Reti Neurali Convoluzionali (\emph{Convolutional Neural Networks} o \textbf{ConvNets}), con l'obiettivo fondamentale di comprendere la loro struttura, il loro funzionamento e l'intuizione teorica e geometrica che ne sottende l'architettura. Il percorso didattico si articola in tre macro-aree principali:
\begin{enumerate}[noitemsep]
    \item Una breve motivazione teorica all'adozione del paradigma del Deep Learning.
    \item La trattazione formale e intuitiva di cosa siano le ConvNets.
    \item L'applicazione delle ConvNets a problemi di localizzazione e riconoscimento visivo, in particolare per il compito di Rilevamento di Oggetti (\emph{Object Detection}).
\end{enumerate}

\subsection{Definizione E Obiettivi Del Deep Learning}

L'apprendimento profondo, o \textbf{Deep Learning}, può essere formalmente definito come la composizione di molteplici trasformazioni non lineari applicate ai dati di input. L'obiettivo primario di questo paradigma consiste nell'apprendere rappresentazioni utili dei dati direttamente a partire dai dati grezzi, senza la necessità di un intervento umano nella progettazione dei descrittori. Tali rappresentazioni intermedie prendono il nome di caratteristiche o \textbf{features}.

Il Deep Learning si presenta in una grande varietà di implementazioni e approcci, che possono essere ricondotti a due macro-categorie fondamentali:
\begin{itemize}[noitemsep]
    \item Apprendimento non supervisionato (\emph{unsupervised learning}).
    \item Apprendimento supervisionato (\emph{supervised learning}).
\end{itemize}
Le reti neurali convoluzionali si collocano primariamente all'interno del paradigma dell'apprendimento profondo di tipo supervisionato.

\subsection{Riepilogo Dell'Apprendimento Supervisionato}

Nell'apprendimento supervisionato disponiamo di un insieme di dati noto come dataset di addestramento (\emph{training dataset}), formalizzato come:
\[ \mathcal{D} = \left\{ \left(\bm{x}^i, y^i\right) \right\}_{i=1}^P \]
dove ciascun termine della formulazione corrisponde a:
\begin{itemize}[noitemsep]
    \item $\bm{x}^i$: l'i-esimo esempio di input per l'addestramento.
    \item $y^i$: l'etichetta target associata all'i-esimo input.
    \item $P$: il numero complessivo di esempi di addestramento disponibili nel dataset.
\end{itemize}

% [IMMAGINE DEL FILE DI SLIDE CON NOME "1-28"] Rappresentazione a scatola nera del modello di apprendimento supervisionato con input x e output y — Slide 4 / Pagina 4 del PDF

Il modello appreso si comporta come una scatola nera che mappa l'input $\bm{x}$ nell'output predetto $y$. A seconda della natura matematica del target $y$, si distinguono diverse tipologie di compiti:
\begin{itemize}
    \item \textbf{Classificazione} (\emph{Classification}): l'output è un'etichetta discreta (ad esempio, data l'immagine di un cucciolo in input, il modello deve produrre l'etichetta testuale \enquote{dog}).
    \item \textbf{Rimozione del rumore} (\emph{Denoising}): l'input è un'immagine corrotta da rumore e il sistema deve ricostruire l'immagine nitida originale. Poiché l'output è continuo per ogni pixel, questo compito si configura come un problema di regressione (\emph{regression}).
    \item \textbf{Riconoscimento Ottico dei Caratteri} (\emph{OCR}): data un'immagine contenente una sequenza di cifre (ad esempio, \enquote{2345}), il modello deve mappare l'input visivo nella stringa di caratteri corrispondente, operazione che rientra nella classe della predizione strutturata (\emph{structured prediction}).
\end{itemize}

% [IMMAGINE DEL FILE DI SLIDE CON NOME "1-28" Esempi visivi di compiti supervisionati: classificazione di immagini, denoising (regressione) e OCR (structured prediction) — Slide 5 / Pagina 5 del PDF

Nell'ambito del Deep Learning supervisionato, questa scatola nera viene spacchettata e sostituita da una sequenza ordinata di molteplici livelli di elaborazione intermedi, i quali trasformano progressivamente l'input attraverso una gerarchia di rappresentazioni.

% [IMMAGINE DEL FILE DI SLIDE CON NOME "1-28" Rappresentazione del Deep Learning supervisionato in cui la scatola nera è divisa in quattro blocchi verticali in serie — Slide 6 / Pagina 6 del PDF

\newpage

\section{Limitazioni Dei Metodi Tradizionali E Necessità Di Apprendere Le Caratteristiche}

Per comprendere appieno l'importanza del Deep Learning, è fondamentale analizzare i limiti intrinseci dell'approccio tradizionale al riconoscimento visivo. Nelle metodologie classiche, la scatola nera di elaborazione è strutturata in una serie di stadi sequenziali ben definiti:

% [IMMAGINE DEL FILE DI SLIDE CON NOME "1-28""] Pipeline di riconoscimento tradizionale: Preprocessing, Feature Extraction (HOG, SIFT), Post-processing (Feature selection, MKL) ed elaborazione tramite un Classificatore — Slide 8-9 / Pagina 8-9 del PDF

Il flusso di lavoro tradizionale prevede:
\begin{enumerate}
    \item \textbf{Pre-elaborazione} (\emph{Preprocessing}): operazioni preliminari sui dati grezzi (es. normalizzazione o ridimensionamento).
    \item \textbf{Estrazione delle caratteristiche} (\emph{Feature Extraction}): algoritmo deputato a convertire i pixel in vettori di caratteristiche stabili e discriminative. Esempi tipici includono HOG (\emph{Histogram of Oriented Gradients}) e SIFT (\emph{Scale-Invariant Feature Transform}).
    \item \textbf{Post-elaborazione} (\emph{Post-processing}): tecniche di selezione delle caratteristiche (\emph{Feature selection}) o metodi di apprendimento a kernel multipli (\textbf{MKL}, \emph{Multiple Kernel Learning}).
    \item \textbf{Classificatore} (\emph{Classifier}): l'algoritmo finale di decisione, come ad esempio macchine a vettori di supporto (\textbf{SVM}) o algoritmi di \emph{boosting}.
\end{enumerate}

La criticità fondamentale di questo schema risiede nel fatto che le fasi di pre-elaborazione, estrazione delle caratteristiche e post-elaborazione sono interamente progettate a mano (\emph{Hand Engineered}) dall'esperto di dominio. Questa dipendenza presenta notevoli svantaggi:
\begin{itemize}[noitemsep]
    \item Lo stadio di estrazione delle caratteristiche è il fattore più critico per determinare l'accuratezza finale dell'intero sistema.
    \item La progettazione e il raffinamento manuale di queste caratteristiche richiedono uno sforzo ingegneristico enorme in termini di tempo e risorse.
    \item Non esiste un criterio matematico universale per stabilire quale sia la caratteristica ottimale per un dato problema visivo.
    \item Continuare a progettare manualmente e con fatica caratteristiche sempre più complesse non rappresenta una soluzione scalabile nel lungo periodo.
\end{itemize}

La soluzione a queste limitazioni consiste nel superare la separazione rigida tra estrazione delle caratteristiche e classificazione, andando ad apprendere la rappresentazione delle caratteristiche direttamente dai dati.

\subsection{Apprendimento Congiunto Di Caratteristiche E Classificatore}

Il Deep Learning si propone di apprendere simultaneamente e in un unico processo le caratteristiche e il classificatore associato, realizzando un sistema di riconoscimento da estremo a estremo (\emph{end-to-end}). Da un punto di vista formale, questo equivale a definire una mappa non lineare complessa che converte direttamente i pixel grezzi dell'immagine di input nelle etichette finali di classe.

Per costruire un sistema così fortemente non lineare, la strategia vincente consiste nell'evitare la progettazione di un'unica funzione enormemente complessa, preferendo invece la combinazione e la sovrapposizione di blocchi costruttivi estremamente semplici.

\newpage

\section{Costruzione Di Funzioni Complesse E Struttura Multilivello}

Per generare una funzione matematica complessa a partire da elementi elementari, si possono ipotizzare due diverse proposte architetturali:

% [IMMAGINE DEL FILE DI SLIDE CON NOME "1-28"] Due proposte di architettura: somma in parallelo (Proposta 1) e composizione sequenziale in serie (Proposta 2) di funzioni non lineari — Slide 12 / Pagina 12 del PDF

\begin{itemize}
    \item \textbf{Proposta 1 (Somma in parallelo)}: l'input viene elaborato in parallelo da diversi blocchi non lineari semplici, e i singoli risultati vengono sommati tra loro.
    \item \textbf{Proposta 2 (Composizione in serie)}: i blocchi non lineari semplici vengono disposti in cascata (in serie), in modo che l'output di ciascun blocco diventi l'input del successivo.
\end{itemize}

Il Deep Learning adotta fermamente la seconda proposta. La composizione di funzioni è infatti il nucleo concettuale dei metodi di apprendimento profondo. Per comprenderne l'efficacia, consideriamo un insieme limitato di funzioni matematiche elementari e non lineari, quali ad esempio:
\[ \mathcal{F} = \left\{ \sin(x), \cos(x), \log(x), x^3, \exp(x) \right\} \]

% [IMMAGINE DEL FILE DI SLIDE CON NOME "1-28"14"] Diagramma raffigurante la transizione da funzioni semplici in un insieme a una funzione composta complessa — Slide 13-14 / Pagina 13-14 del PDF

Attraverso l'applicazione sequenziale di queste funzioni, è possibile descrivere una funzione complessa e fortemente non lineare come:
\[ g(x) = \log\left(\cos\left(\exp\left(\sin^3(x)\right)\right)\right) \]
In un'architettura di rete profonda, ciascuna di queste funzioni semplici possiede dei parametri interni liberi che saranno ottimizzati durante il processo di addestramento.

\subsection{Intuizione Geometrica Delle Reti Neurali Profonde}

% [IMMAGINE DEL FILE DI SLIDE CON NOME "1-28"] Schema gerarchico di una rete neurale profonda a 4 livelli che elabora l'immagine di un'auto per produrre una distribuzione di probabilità — Slide 15-17 / Pagina 15-17 del PDF

L'intuizione alla base delle reti profonde risiede nella decomposizione del problema in livelli di astrazione gerarchici. Considerando l'elaborazione di un'immagine (ad esempio, un'automobile), il flusso attraversa una serie di livelli sequenziali (indicati come \emph{Layer 1}, \emph{Layer 2}, \emph{Layer 3} e \emph{Layer 4}):
\begin{itemize}[noitemsep]
    \item Ciascun livello (rappresentato come un blocco nero nel diagramma) esegue una specifica trasformazione dotata di parametri addestrabili.
    \item La composizione sequenziale di tali livelli produce un sistema complessivo altamente non lineare.
    \item L'ultimo livello (\emph{Layer 4}) genera in output una distribuzione di probabilità sulle diverse categorie ammissibili del problema.
\end{itemize}

\newpage

\subsection{Formulazione Matematica Di Un Singolo Livello}

Un singolo livello elementare di una rete neurale consiste nell'applicazione di una trasformazione lineare (prodotto tra matrici e vettori con aggiunta di un bias) seguita da una funzione di attivazione non lineare. Matematicamente, l'operazione si esprime come:
\begin{align*}
\bm{z} &= \bm{W}\bm{x} + \bm{b} \\
\bm{a} &= f(\bm{z})
\end{align*}
dove:
\begin{itemize}[noitemsep]
    \item $\bm{W}$ rappresenta la matrice dei pesi (\emph{weights}), ossia i parametri da apprendere.
    \item $\bm{x}$ è il vettore di input al livello corrente (corrispondente all'output del livello precedente).
    \item $\bm{b}$ è il vettore di bias.
    \item $f$ è una funzione non lineare semplice, detta funzione di attivazione. Una scelta standard e di grande efficacia è la ReLU (\emph{Rectified Linear Unit}), definita come:
    \[ f(x) = \max(0, x) \]
\end{itemize}

% [IMMAGINE DEL FILE DI SLIDE CON NOME "1-28"] Rappresentazioni grafiche di un singolo livello neurale: schema a livello di singolo neurone e schema in forma vettoriale compatta — Slide 19 / Pagina 19 del PDF

Da un punto di vista grafico, l'output di un singolo neurone può essere descritto dalla relazione:
\[ h = f\left(\bm{w}^T\bm{x}\right) \]
dove il valore scalare $h$ prende il nome di neurone, unità nascosta (\emph{hidden unit}) o caratteristica (\emph{feature}). La slide illustra due schemi di visualizzazione:
\begin{enumerate}
    \item Una scomposizione in cui le singole componenti dell'input $x_1, x_2, x_3, x_4$ vengono pesate individualmente da pesi $w_1, w_2, w_3, w_4$, sommate e passate all'operatore non lineare $f$.
    \item Uno schema vettoriale in cui l'intero vettore di input $\bm{x}$ viene trasformato tramite la matrice $\bm{W}$ per produrre il vettore di caratteristiche nascoste $\bm{h}$.
\end{enumerate}

\newpage

\section{Addestramento Congiunto E Flusso Di Ottimizzazione}

L'addestramento congiunto (\emph{joint training}) di tutti i moduli che compongono la rete profonda richiede la definizione di un'architettura complessiva per il calcolo dell'errore:

% [IMMAGINE DEL FILE DI SLIDE CON NOME "1-28"] Schema dell'architettura di addestramento congiunto con blocchi in serie e calcolo del blocco di errore finale — Slide 20 / Pagina 20 del PDF

Una rete neurale multistrato contenente più di due livelli prende il nome di rete profonda (\emph{deep net}). Per configurare il modello, l'utente deve definire a priori i seguenti iperparametri:
\begin{itemize}[noitemsep]
    \item Il numero complessivo di livelli (\emph{layers}).
    \item Il numero di unità nascoste per ciascun livello.
    \item La tipologia di livelli (es. convoluzionali, completamente connessi).
    \item La funzione di perdita (\emph{loss function}), utilizzata per quantificare lo scostamento tra la predizione del modello $\hat{y}$ e l'etichetta reale $y$.
\end{itemize}

\subsection{Le Fasi Di Addestramento Della Rete Neurale}

L'addestramento della rete avviene in modo iterativo su piccoli gruppi di dati di addestramento, chiamati mini-batch. Ciascun passo di ottimizzazione si compone di tre fasi sequenziali:

% [IMMAGINE DEL FILE DI SLIDE CON NOME "1-28"] Animazione delle fasi di addestramento: Forward Propagation (F-PROP), Backward Propagation (B-PROP) e aggiornamento dei parametri — Slide 21-27 / Pagina 21-27 del PDF

\begin{enumerate}
    \item \textbf{Fase A: Propagazione in avanti (F-PROP / Forward Propagation)} \\
    I dati del mini-batch corrente vengono immessi nella rete e fluiscono sequenzialmente attraverso i livelli (dal livello 1 al livello 3) per calcolare la predizione finale e il valore della funzione di perdita $\mathcal{L}$.
    \item \textbf{Fase B: Retropropagazione (B-PROP / Backward Propagation)} \\
    Si calcola il gradiente della funzione di perdita rispetto a tutti i parametri liberi (pesi e bias) della rete. Il calcolo avviene a ritroso, propagando l'errore dal livello finale verso i livelli iniziali (dal livello 3 al livello 1).
    \item \textbf{Fase C: Aggiornamento dei parametri} \\
    I gradienti calcolati vengono utilizzati per aggiornare i pesi $\bm{W}$ della rete attraverso una regola di discesa del gradiente. La formulazione matematica dell'aggiornamento è:
    \[ \boxed{\bm{W} \leftarrow \bm{W} - \eta \frac{\mathcal{L}}{d\bm{W}}} \]
    dove $\eta$ rappresenta il tasso di apprendimento (\emph{learning rate}) e $\frac{\mathcal{L}}{d\bm{W}}$ è il gradiente della perdita rispetto alla matrice dei pesi $\bm{W}$.
\end{enumerate}

\newpage

\subsection{La Sfida Dei Dati Immagine Nei Livelli Completamente Connessi}

Quando si lavora con dati strutturati come le immagini, l'applicazione diretta di livelli completamente connessi (\emph{fully connected}) mostra limiti evidenti in termini di scalabilità e numero di parametri.

% [IMMAGINE DEL FILE DI SLIDE CON NOME "1-28"] Schema di connessione completa applicata a un'immagine di input, in cui ogni pixel è connesso a un'unità del livello successivo — Slide 28 / Pagina 28 del PDF

In un livello completamente connesso, ogni singolo pixel dell'immagine bidimensionale di input viene trattato come una caratteristica indipendente e connesso a ciascun neurone del livello successivo tramite un peso specifico appartenente alla matrice dei pesi $\bm{W}$. Questo approccio ignora completamente la topologia spaziale dell'immagine e porta a un'esplosione del numero di parametri da apprendere, rendendo necessaria l'introduzione di operatori con connettività locale e condivisione dei pesi, elementi fondanti delle reti neurali convoluzionali.

\section{Dalla Connessione Densa Alla Connessione Locale}

Nell'ambito dell'elaborazione di immagini digitali tramite modelli di apprendimento profondo, la transizione dalle architetture completamente connesse (\textbf{fully connected}) a quelle convoluzionali rappresenta una svolta fondamentale dettata da vincoli fisici e computazionali.

Se consideriamo un'immagine ad alta risoluzione come input di una rete neurale classica, sorge immediatamente il problema della proliferazione incontrollata dei parametri. Supponiamo di avere un'immagine in ingresso di dimensioni $H \times W = 1000 \times 1000$ pixel (pari a un input di $10^6$ valori). Nel caso in cui si decida di utilizzare uno strato nascosto completamente connesso composto da $H_u = 10^6$ (un milione) di unità, il numero di pesi della matrice di connessione (trascurando i termini di bias) è modellabile come:

\[
P_{\text{FC}} \approx (H \times W) \times H_u = 10^6 \times 10^6 = 10^{12} \text{ parametri}
\]

Nelle slide viene evidenziato come, anche in scenari intermedi con un numero inferiore di connessioni globali, si arrivi rapidamente alla quota di ben \textcolor{red}{$1\text{B}$ (un miliardo) di parametri}, una dimensione che rende l'addestramento impraticabile per motivi di memoria, computazione e forte rischio di sovradimensionamento (\emph{overfitting}).

% [IMMAGINE DEL FILE DI SLIDE CON NOME "29-55"] Rappresentazione concettuale di una connessione densa (fully connected) applicata a un'immagine ad alta risoluzione con conseguente esplosione parametrica. — Slide 65 / Pagina 1 del PDF

Questa inefficienza strutturale si spiega attraverso due considerazioni chiave:
\begin{itemize}
    \item \textbf{La correlazione spaziale è locale} (\emph{Spatial correlation is local}): in un'immagine, l'informazione semantica rilevante è racchiusa nelle relazioni tra pixel vicini. Un pixel nell'angolo in alto a sinistra ha una correlazione statistica pressoché nulla con un pixel situato nell'angolo in basso a destra. Non vi è quindi alcun motivo fisico per connettere ogni singolo pixel a tutte le unità dello strato successivo.
    \item \textbf{Allocazione ottimale delle risorse}: è preferibile concentrare la capacità computazionale e i parametri del modello sulla cattura di pattern locali piuttosto che disperderli in connessioni globali ridondanti.
\end{itemize}

Per superare questa limitazione, si introduce il principio della \textbf{connessione a regioni locali} (\emph{local receptive fields}). Invece di collegare ciascuna unità nascosta all'intera immagine, ogni unità viene connessa esclusivamente a una piccola porzione spaziale contigua dell'input.

Considerando nuovamente lo scenario precedente con un'immagine di $1000 \times 1000$ pixel e uno strato nascosto di $1\text{M}$ di unità, se limitiamo l'area di connessione di ciascuna unità a una finestra locale (denominata anche dimensione del filtro o del kernel) pari a $K_h \times K_w = 10 \times 10$ pixel, il numero di parametri dello strato $P_{\text{local}}$ si riduce drasticamente:

\[
P_{\text{local}} = H_u \times (K_h \times K_w) = 10^6 \times (10 \times 10) = 10\text{M} \text{ parametri}
\]

% [IMMAGINE DEL FILE DI SLIDE CON NOME "29-55"] Connessioni limitate a regioni spaziali locali (local receptive fields) sull'immagine di Einstein. — Slide 67 / Pagina 2 del PDF

In questo modo, riducendo la connessione a regioni locali, il numero di parametri scende da un miliardo a soli $10$ milioni.

\newpage

\subsection{Condivisione Dei Pesi E Convoluzione (Weight Sharing)}

La riduzione a $10\text{M}$ di parametri rappresenta un miglioramento sensibile, ma per reti profonde con molti strati questo numero risulta ancora eccessivamente elevato. Per ottimizzare ulteriormente il modello, si introduce la tecnica della \textbf{condivisione dei pesi} (\emph{weight sharing}).

L'intuizione alla base di questa scelta è di natura prettamente fisica:
\[
\text{\emph{Le caratteristiche di interesse (come i bordi o i contorni) possono apparire in qualsiasi punto dell'immagine.}}
\]
Se un determinato filtro è efficace nel rilevare un bordo orientato nell'angolo superiore sinistro dell'immagine, lo stesso filtro sarà altrettanto efficace nel rilevare lo stesso tipo di bordo nell'angolo inferiore destro. Di conseguenza, ha senso \textbf{condividere gli stessi parametri del filtro} (il kernel) in tutte le diverse posizioni spaziali dell'immagine. Questa operazione corrisponde matematicamente alla \textbf{convoluzione con kernel appresi}.

% [IMMAGINE DEL FILE DI SLIDE CON NOME "29-55"] Condivisione dei parametri (weight sharing) dello stesso kernel su diverse posizioni spaziali dell'immagine. — Slide 69 / Pagina 3 del PDF

Nelle reti neurali convoluzionali (\textbf{CNN}), per poter catturare una varietà sufficiente di pattern differenti nello stesso strato, non ci si limita ad apprendere un solo filtro, ma si apprende un insieme di \textbf{molteplici filtri}.

Prendiamo come esempio un'architettura convoluzionale con le seguenti specifiche:
\begin{itemize}
    \item Immagine in ingresso: $1000 \times 1000$ pixel.
    \item Numero di filtri indipendenti: $N_f = 100$.
    \item Dimensione di ciascun filtro: $10 \times 10$ pixel.
\end{itemize}

Grazie al meccanismo di condivisione dei pesi, il numero complessivo di parametri $P_{\text{conv}}$ da apprendere dipende unicamente dalle dimensioni e dal numero dei filtri stessi, risultando completamente indipendente dalle dimensioni dell'immagine e dal numero di unità spaziali dello strato successivo:

\[
\boxed{P_{\text{conv}} = N_f \times (K_h \times K_w) = 100 \times (10 \times 10) = 10\,000 = 10\text{K} \text{ parametri}}
\]

% [IMMAGINE DEL FILE DI SLIDE CON NOME "29-55"] Rappresentazione schematica di una rete convoluzionale che apprende molteplici filtri condivisi. — Slide 70 / Pagina 4 del PDF

Questo principio di conservazione dei parametri consente di ridurre la complessità del modello da $10\text{M}$ a soli $10\text{K}$ parametri, garantendo al contempo un'elevata efficienza computazionale e una forte regolarizzazione implicita. Questa formulazione teorica risale al lavoro pionieristico di LeCun et al., \emph{"Gradient-based learning applied to document recognition"}, IEEE 1998.

\newpage

\subsection{Dettagli Sui Canali Multipli E Mappe Di Caratteristiche}

Le immagini reali sono generalmente caratterizzate da molteplici canali di colore (ad esempio i canali RGB per il rosso, verde e blu). Quando l'input presenta $C$ canali, la convoluzione si estende alla terza dimensione:
\begin{itemize}
    \item Se l'input ha $C=3$ canali (R, G, B), vengono applicati $3$ filtri separati di dimensione $K \times K$, uno per ciascun canale.
    \item Questi filtri formano un unico \textbf{kernel 3D} di dimensioni complessive $K \times K \times C$.
    \item I contributi convolutivi calcolati sui singoli canali vengono sommati punto a punto per generare un singolo valore di output per ciascuna posizione spaziale.
\end{itemize}

L'output derivante dalla convoluzione dell'input con un singolo filtro tridimensionale prende il nome di \textbf{mappa di caratteristiche} (\emph{feature map}).

% [IMMAGINE DEL FILE DI SLIDE CON NOME "29-55"] Applicazione di un kernel 3D a mappe di caratteristiche in ingresso con 3 canali (RGB) per produrre una mappa di caratteristiche in uscita. — Slide 71 / Pagina 5 del PDF

Questo meccanismo equivale all'applicazione di una finestra scorrevole (\emph{sliding window}). Ad esempio, l'output fornito da un filtro parziale nel modello classico DPM (\emph{Deformable Parts Model}) può essere interpretato a tutti gli effetti come una mappa di caratteristiche.

Per poter rappresentare l'informazione visiva in modo completo, si utilizzano molteplici filtri in parallelo. Ciascun filtro ha il compito di rilevare specifiche caratteristiche (come texture, bordi o variazioni cromatiche) nell'output dello strato precedente. Per catturare questa molteplicità di aspetti, la rete apprende un set di filtri differenti, generando così un volume di mappe di caratteristiche in uscita.

\[
\text{\textbf{Nota bene:} il numero di mappe di caratteristiche in uscita è solitamente maggiore del numero di mappe in ingresso.}
\]

% [IMMAGINE DEL FILE DI SLIDE CON NOME "29-55"] Generazione di molteplici mappe di caratteristiche in uscita mediante l'applicazione di più filtri 3D. — Slide 71 / Pagina 6 del PDF

L'operazione di filtraggio convoluzionale gode di due proprietà matematiche e strutturali fondamentali:
\begin{enumerate}
    \item \textbf{Equivarianza alla traslazione} (\emph{Translation equivariance}): se l'immagine di input subisce una traslazione spaziale, la mappa di caratteristiche risultante subirà la medesima traslazione. Formalmente, definendo l'operatore di traslazione come $T_g$ e l'operazione di convoluzione come $\ast$, si ha:
    \[
    (T_g \bm{x}) \ast \bm{w} = T_g (\bm{x} \ast \bm{w})
    \]
    \item \textbf{Pesi del filtro vincolati} (\emph{Tied filter weights}): lo stesso identico valore di peso viene utilizzato in ogni posizione spaziale dell'immagine, garantendo il drastico contenimento del numero di parametri del modello.
\end{enumerate}

% [IMMAGINE DEL FILE DI SLIDE CON NOME "29-55"] Esempio reale di filtraggio convoluzionale (equivarianza alla traslazione) applicato a un'immagine urbana con un filtro per bordi orientati. — Slide 72 / Pagina 7 del PDF

\newpage

\section{Costruzione Dell'Invarianza Alla Traslazione Tramite Pooling Spaziale}

Mentre la convoluzione garantisce l'equivarianza alla traslazione, in molti compiti di percezione visiva è necessario ottenere una proprietà differente: l'\textbf{invarianza alla traslazione} (\emph{translation invariance}).

Per comprendere la differenza, supponiamo che un determinato filtro sia specializzato come rilevatore di occhi (\emph{"eye" detector}). La domanda fondamentale che sorge è la seguente:
\[
\text{\emph{Come possiamo rendere il rilevamento robusto rispetto all'esatta posizione spaziale dell'occhio?}}
\]

% [IMMAGINE DEL FILE DI SLIDE CON NOME "29-55"] Rappresentazione concettuale della necessità di invarianza spaziale per un rilevatore di elementi locali (es. occhio). — Slide 72 / Pagina 8 del PDF

La risposta a questo problema risiede nell'introduzione del \textbf{pooling spaziale} (\emph{spatial pooling}).

Aggregando o eseguendo un'operazione di "pooling" (come ad esempio il calcolo del massimo, \emph{max pooling}, o della media, \emph{average pooling}) sulle risposte del filtro in posizioni limitrofe differenti, si ottiene una forte robustezza rispetto all'esatta localizzazione spaziale delle caratteristiche rilevate. Se la caratteristica si sposta leggermente all'interno della finestra di pooling, l'output dell'operazione rimarrà invariato o varierà in modo minimo.

In aggiunta a questa proprietà, l'operazione di pooling effettua un sotto-campionamento (\emph{subsampling}) dell'immagine. Questo processo riduce la risoluzione spaziale delle mappe di caratteristiche, consentendo agli strati successivi della rete di analizzare regioni spaziali progressivamente più ampie, aumentando così il campo ricettivo efficace (\emph{receptive field}) dei filtri più profondi.

% [IMMAGINE DEL FILE DI SLIDE CON NOME "29-55"] Meccanismo di pooling spaziale per combinare le risposte di filtri locali adiacenti in un'unica unità più robusta. — Slide 73 / Pagina 9 del PDF

\newpage

\subsection{Sintesi Di Un Tipico Strato Convoluzionale}

L'insieme coordinato di queste operazioni definisce la struttura di un singolo \textbf{strato convoluzionale tipico}. Il flusso sequenziale di elaborazione è così strutturato:

\begin{align*}
\text{\textbf{Immagine di Input}} &\longrightarrow \text{\textbf{Convoluzione (Appresa)}} \\
&\longrightarrow \text{\textbf{Non-linearità}} \\
&\longrightarrow \text{\textbf{Pooling Spaziale (Opzionale)}} \\
&\longrightarrow \text{\textbf{Normalizzazione (Opzionale)}} \\
&\longrightarrow \text{\textbf{Mappe di Caratteristiche di Output}}
\end{align*}

% [IMMAGINE DEL FILE DI SLIDE CON NOME "29-55"] Diagramma di flusso delle operazioni sequenziali all'interno di un tipico strato convoluzionale. — Slide 74 / Pagina 10 del PDF

Questi blocchi funzionali possono essere impilati sequenzialmente per formare architetture profonde, le quali vengono addestrate end-to-end con le medesime tecniche di ottimizzazione (retropropagazione del gradiente e discesa stocastica del gradiente) impiegate per le reti neurali densamente connesse. L'ultimo strato di una tipica rete convoluzionale è solitamente costituito da uno strato completamente connesso (\emph{fully connected}), avente una dimensione di uscita pari al numero di classi del problema di classificazione.

\newpage

\subsection{Confronto Con L'Algoritmo SIFT}

È estremamente istruttivo confrontare la struttura di uno strato convoluzionale profondo con i descrittori locali classici sviluppati manualmente dagli esperti, in particolare con l'algoritmo \textbf{SIFT} (\emph{Scale-Invariant Feature Transform}, Lowe, IJCV 2004).

La pipeline SIFT si articola nelle seguenti fasi:
\begin{enumerate}
    \item \textbf{Applicazione di filtri orientati} (\emph{Apply oriented filters}): calcolo dei gradienti orientati sull'immagine di input.
    \item \textbf{Pooling spaziale} (\emph{Spatial pool}): aggregazione locale delle risposte dei gradienti tramite somma (\emph{Sum}).
    \item \textbf{Normalizzazione} (\emph{Normalize}): normalizzazione del vettore di caratteristiche finale a lunghezza unitaria per garantire robustezza alle variazioni di illuminazione.
\end{enumerate}

% [IMMAGINE DEL FILE DI SLIDE CON NOME "29-55"] Diagramma di flusso della pipeline classica SIFT (Lowe, IJCV 2004). — Slide 75 / Pagina 11 del PDF

Esiste una corrispondenza diretta e biunivoca tra le fasi dell'algoritmo SIFT e le operazioni che compongono uno strato di una moderna rete neurale convoluzionale:

% [IMMAGINE DEL FILE DI SLIDE CON NOME "29-55"] Corrispondenza visiva tra i blocchi funzionali dell'algoritmo SIFT e uno strato convoluzionale di una CNN. — Slide 76 / Pagina 12 del PDF

\begin{itemize}
    \item La fase SIFT \emph{"Apply oriented filters"} corrisponde alla \textbf{Convoluzione}. La differenza fondamentale risiede nel fatto che in SIFT i filtri sono ingegnerizzati manualmente e fissi, mentre nella CNN i pesi dei filtri vengono interamente appresi dai dati tramite ottimizzazione.
    \item La fase SIFT \emph{"Spatial pool (Sum)"} corrisponde al \textbf{Pooling Spaziale} della CNN.
    \item La fase SIFT \emph{"Normalize to unit length"} corrisponde allo strato di \textbf{Normalizzazione} della CNN.
\end{itemize}

\newpage

\section{Apprendimento Gerarchico E Composizionale}

Il successo delle reti convoluzionali profonde si basa sul concetto di \textbf{composizionalità} delle caratteristiche (\emph{composition idea}):
\[
\text{\emph{Ogni strato impara un rilevatore di caratteristiche combinando l'output dello strato precedente.}}
\]
Questa struttura gerarchica comporta una conseguenza di grande rilievo: man mano che si sale verso gli strati più profondi dell'architettura, la rete apprende rappresentazioni caratterizzate da un livello di astrazione via via crescente.

Per analizzare concretamente questo fenomeno, si fa riferimento all'architettura sviluppata da Alex Krizhevsky, Ilya Sutskever e Geoffrey Hinton, nota come \textbf{AlexNet}, vincitrice della competizione internazionale ILSVRC-2012 con un errore top-5 del \textcolor{red}{$18.2\%$}.

Le caratteristiche principali di questa architettura storica sono:
\begin{itemize}
    \item Un totale di \textbf{8 strati} con parametri apprendibili.
    \item Addestramento effettuato sul dataset \textbf{ImageNet}, composto da $1000$ categorie di oggetti, $1.2\text{M}$ di immagini di addestramento e $150\text{K}$ immagini di test.
    \item Una sequenza strutturale così definita:
    \begin{align*}
    \text{Input} &\longrightarrow \text{Layer 1: Conv + Pool} \longrightarrow \text{Layer 2: Conv + Pool} \longrightarrow \text{Layer 3: Conv} \\
    &\longrightarrow \text{Layer 4: Conv} \longrightarrow \text{Layer 5: Conv + Pool} \longrightarrow \text{Layer 6: Full} \\
    &\longrightarrow \text{Layer 7: Full} \longrightarrow \text{Softmax Output}
    \end{align*}
\end{itemize}

% [IMMAGINE DEL FILE DI SLIDE CON NOME "29-55"] Architettura a blocchi a 8 strati di AlexNet (Winner ILSVRC-2012). — Slide 77 / Pagina 14 del PDF

Dal punto di vista implementativo, come illustrato nel diagramma originale, la rete era originariamente suddivisa su due GPU parallele per far fronte ai severi limiti di memoria dell'epoca. L'input dell'architettura è un volume di dimensioni $224 \times 224 \times 3$, elaborato nel primo strato da filtri di dimensione $11 \times 11 \times 3$ applicati con uno stride spaziale pari a 4 pixel.

% [IMMAGINE DEL FILE DI SLIDE CON NOME "29-55"] Diagramma dettagliato della scomposizione sui due canali GPU dell'architettura AlexNet dal paper originale. — Slide 78 / Pagina 15 del PDF

\newpage

\subsection{Visualizzazione Delle Caratteristiche Apprese}

Esaminando visivamente i filtri appresi e i corrispondenti patch d'immagine che massimizzano l'attivazione delle unità ai vari livelli della rete, si riscontra sperimentalmente la validità dell'ipotesi gerarchica:

\begin{enumerate}
    \item \textbf{Filtri del Primo Strato (Layer 1)}:
    Visualizzando gli 81 filtri di dimensione $11 \times 11 \times 3$ del primo strato, si nota come essi catturino esclusivamente caratteristiche a basso livello (\emph{low-level features}), quali bordi orientati (\emph{oriented edges}) e variazioni locali di colore (\emph{blobs}). Questi filtri sono del tutto analoghi a quelli utilizzati dall'algoritmo SIFT per il calcolo dei gradienti spaziali.

    % [IMMAGINE DEL FILE DI SLIDE CON NOME "29-55"] Visualizzazione degli 81 filtri di dimensione 11x11x3 appresi nel primo strato di AlexNet. — Slide 79 / Pagina 16 del PDF

    % [IMMAGINE DEL FILE DI SLIDE CON NOME "29-55"] I top 9 patch di immagini reali che attivano maggiormente ciascun filtro del primo strato. — Slide 80 / Pagina 17 del PDF

    \item \textbf{Filtri del Secondo Strato (Layer 2)}:
    Combinando le risposte dei filtri a basso livello del primo strato, le unità del secondo strato iniziano a rispondere a strutture dotate di una maggiore complessità geometrica, arrivando a rilevare motivi superficiali ripetitivi e \textbf{consistenze materiche} (\emph{textures}).

    % [IMMAGINE DEL FILE DI SLIDE CON NOME "29-55"] Visualizzazione dei patch reali che attivano i filtri del secondo strato, evidenziando motivi e texture astratte. — Slide 81 / Pagina 18 e 19 del PDF

    \item \textbf{Filtri del Terzo Strato (Layer 3)}:
    In questo livello si assiste alla transizione verso la rappresentazione di parti fisiche di oggetti o strutture complesse (come ruote di biciclette, reticoli, motivi a nido d'ape, porzioni di volti o caratteri tipografici).

    % [IMMAGINE DEL FILE DI SLIDE CON NOME "29-55"] Visualizzazione dei patch di massima attivazione per il terzo strato, che mostrano parti di oggetti e pattern complessi. — Slide 82 / Pagina 20 e 21 del PDF

    \item \textbf{Filtri del Quarto Strato (Layer 4)}:
    Le caratteristiche apprese mostrano una spiccata astrazione semantica, attivandosi selettivamente in presenza di pattern complessi e parti anatomiche complesse (ad esempio, interi musi di cani, zampe di volatili o strutture geometriche circolari).

    % [IMMAGINE DEL FILE DI SLIDE CON NOME "29-55"] Patch di massima attivazione per lo strato 4, indicanti la comparsa di classi semantiche e strutture complesse (es. musi di cane). — Slide 83 / Pagina 22 del PDF

    \item \textbf{Filtri del Quinto Strato (Layer 5)}:
    I filtri dello strato convoluzionale più profondo rispondono in modo specifico a interi oggetti appartenenti a categorie semantiche ben definite (quali tastiere di computer, cani nella loro interezza, fiori o biciclette).

    % [IMMAGINE DEL FILE DI SLIDE CON NOME "29-55"] Patch di massima attivazione per lo strato 5, con rilevazione di interi oggetti o categorie semantiche complesse. — Slide 84 / Pagina 23 del PDF
\end{enumerate}

\newpage

\section{Uso Delle ConvNet Come Estrattori Generici Di Caratteristiche}

L'evidenza empirica mostra che una rete convoluzionale profonda adeguatamente addestrata su un dataset di ampie dimensioni (come ImageNet) costituisce un eccellente \textbf{estrattore di caratteristiche generiche}.

È possibile utilizzare una rete pre-addestrata interrompendo l'elaborazione a un determinato strato d'interesse (\emph{"chopping" the network}) e impiegando le attivazioni di tale strato come rappresentazione vettoriale delle caratteristiche (\emph{feature representation}). Questa rappresentazione può essere usata per addestrare un classificatore esterno, come una \textbf{SVM} (\emph{Support Vector Machine}), su un dataset visivo completamente differente.

Per quantificare le prestazioni di questa metodologia, le slide riportano uno studio comparativo condotto sui dataset di riferimento \textbf{Caltech-101} (con 30 campioni per classe) e \textbf{Caltech-256} (con 60 campioni per classe). Le caratteristiche estratte dai diversi strati di AlexNet sono state utilizzate per addestrare classificatori SVM lineari, o confrontate con le uscite dirette dello strato Softmax:

\begin{table}[H]
\centering
\caption{Accuratezza di classificazione (\%) su dataset Caltech-101 e Caltech-256 al variare dello strato di estrazione delle caratteristiche (dal Layer 1 al Layer 7).}
\label{tab:feature_extraction_results}
\vspace{0.2cm}
\begin{tabular}{lcc}
\toprule
\textbf{Metodo di Classificazione (Strato)} & \textbf{Caltech-101 (30/classe)} & \textbf{Caltech-256 (60/classe)} \\
\midrule
SVM (1)     & $44.8 \pm 0.7$ & $24.6 \pm 0.4$ \\
SVM (2)     & $66.2 \pm 0.5$ & $39.6 \pm 0.3$ \\
SVM (3)     & $72.3 \pm 0.4$ & $46.0 \pm 0.3$ \\
SVM (4)     & $76.6 \pm 0.4$ & $51.3 \pm 0.1$ \\
SVM (5)     & $86.2 \pm 0.8$ & $65.6 \pm 0.3$ \\
SVM (7)     & $85.5 \pm 0.4$ & $71.7 \pm 0.2$ \\
\midrule
Softmax (5) & $82.9 \pm 0.4$ & $65.7 \pm 0.5$ \\
Softmax (7) & $85.4 \pm 0.4$ & $72.6 \pm 0.1$ \\
\bottomrule
\end{tabular}
\end{table}

Dall'analisi quantitativa si evincono chiaramente due conclusioni fondamentali:
\begin{itemize}
    \item L'accuratezza cresce costantemente all'aumentare della profondità dello strato da cui vengono estratte le caratteristiche, confermando l'importanza dell'astrazione gerarchica.
    \item L'estrazione dal Layer 7 permette di raggiungere le prestazioni più elevate, pari all' $85.5\%$ su Caltech-101 e al $71.7\%$ ($72.6\%$ con Softmax) su Caltech-256.
\end{itemize}

Oltre all'utilizzo della rete come estrattore di caratteristiche statiche, è possibile migliorare ulteriormente l'accuratezza applicando la tecnica del \textbf{fine-tuning} (messa a punto fine). Questa procedura prevede di inizializzare i pesi della rete con quelli del modello pre-addestrato su ImageNet, per poi aggiornarli parzialmente o totalmente tramite retropropagazione sul nuovo dataset di destinazione.

\newpage

\subsection{Rilevamento Di Oggetti Tramite ConvNet E Modello R-CNN}

Sfruttando l'elevato potere discriminativo delle caratteristiche estratte dalle reti convoluzionali, è possibile estendere il loro utilizzo dal compito di classificazione globale a quello di rilevamento e localizzazione di oggetti (\textbf{object detection}).

Un metodo lineare per formulare il problema consiste nei seguenti passaggi:
\begin{enumerate}
    \item Ritagliare l'immagine di input in corrispondenza di box delimitatori candidati (\emph{bounding boxes}).
    \item Estrarre le caratteristiche convoluzionali per ciascun ritaglio.
    \item Addestrare una SVM lineare per valutare l'appartenenza a ciascuna categoria di oggetto.
    \item Identificare le regioni ad alta probabilità, calcolare il punteggio di confidenza associato ed eseguire l'algoritmo di \textbf{soppressione dei non-massimi} (\emph{non-maximum suppression}, NMS) per rimuovere le sovrapposizioni ridondanti dello stesso oggetto.
\end{enumerate}

Questo approccio è stato formalizzato e ottimizzato da Ross Girshick nel modello denominato \textbf{R-CNN: Regions with CNN features}.

% [IMMAGINE DEL FILE DI SLIDE CON NOME "29-55"] Diagramma di flusso del funzionamento del modello R-CNN per l'object detection. — Slide 87 / Pagina 26 del PDF

La pipeline di funzionamento di R-CNN si articola in quattro fasi sequenziali:
\begin{enumerate}
    \item \textbf{Immagine di input}: inserimento dell'immagine da analizzare.
    \item \textbf{Estrazione delle proposte di regione}: generazione di circa $2000$ regioni candidate per immagine (\emph{region proposals}) che potrebbero contenere oggetti di interesse.
    \item \textbf{Calcolo delle caratteristiche}: elaborazione di ciascuna regione tramite una rete convoluzionale profonda per estrarne il corrispondente vettore di caratteristiche.
    \item \textbf{Classificazione delle regioni}: classificazione di ciascuna proposta tramite SVM lineari specifiche per classe (ad esempio valutando se l'oggetto rilevato sia un aeroplano, una persona, uno schermo televisivo, ecc.).
\end{enumerate}

R-CNN ha rappresentato un'importante innovazione nel campo della percezione visiva, posizionandosi come il miglior metodo sul benchmark \textbf{PASCAL 2012 Detection} e registrando un incremento prestazionale del \textcolor{red}{$30\%$} rispetto a tutte le metodologie concorrenti dell'epoca.

\newpage

\subsection{Librerie Di Riferimento Per Reti Convoluzionali}

In ambito accademico e industriale, lo sviluppo storico delle reti convoluzionali si è consolidato attraverso l'uso di librerie software dedicate e ottimizzate per il calcolo su GPU:

\begin{itemize}
    \item \textbf{Cuda-convnet}: sviluppata direttamente da Alex Krizhevsky (all'epoca Google), nota per essere stata una delle primissime implementazioni altamente efficienti scritte in C++/CUDA per l'addestramento accelerato su schede grafiche NVIDIA.
    \item \textbf{Caffe}: sviluppata da Yangqing Jia presso l'Università di Berkeley (successivamente Google). Ha riscosso una notevole diffusione grazie alla sua velocità, alla separazione pulita tra definizione dell'architettura (tramite file di configurazione in formato testuale prototxt) e codice sorgente, e alla disponibilità del \emph{"Model Zoo"} contenente i pesi pre-addestrati di AlexNet.
    \item \textbf{Torch7}: libreria basata sul linguaggio di programmazione Lua e ampiamente supportata da importanti istituzioni accademiche e industriali, tra cui Idiap, NYU (sotto la direzione di Yann LeCun) e NEC.
\end{itemize}

% AA Slides 11.2

\section{Informazioni Amministrative E Panoramica Dei Framework}

Prima di addentrarsi nell'analisi dettagliata delle architetture delle reti convoluzionali, vengono riportati alcuni avvisi amministrativi relativi al corso:
\begin{itemize}
    \item La seconda serie di esercizi (\textbf{A2}) deve essere consegnata entro giovedì 4 maggio.
    \item L'esame parziale (\textbf{Midterm}) si terrà in aula martedì 9 maggio e coprirà tutti gli argomenti trattati fino alla lezione di giovedì 4 maggio inclusa.
    \item La sessione dei poster (\textbf{Poster session}) è programmata per martedì 6 giugno, dalle ore 12:00 alle ore 15:00.
\end{itemize}

\subsection{Evoluzione Dei Framework Di Deep Learning}

I framework per il deep learning hanno subito un'importante evoluzione tecnologica volta a semplificare la prototipazione e a migliorare le prestazioni computazionali. Di seguito vengono tracciate le principali linee di sviluppo e di transizione dei software più diffusi:

% [IMMAGINE DEL FILE DI SLIDE CON NOME 11.2 "1 - 26"] Panoramica dei framework di deep learning e delle loro evoluzioni storiche — Slide 3 / Pagina 3 del PDF

\begin{itemize}
    \item \textbf{Caffe} (sviluppato originariamente presso la UC Berkeley) è evoluto in \textbf{Caffe2} (sotto l'egida di Facebook).
    \item \textbf{Torch} (nato dalla collaborazione tra NYU e Facebook) è stato rifondato nella versione Python come \textbf{PyTorch} (Facebook).
    \item \textbf{Theano} (sviluppato alla U Montreal) ha storicamente spianato la strada per la creazione di \textbf{TensorFlow} (Google).
    \item Tra gli altri framework significativi si annoverano \textbf{Paddle} (Baidu), \textbf{CNTK} (Microsoft) e \textbf{MXNet} (Amazon). Quest'ultimo, sviluppato originariamente da istituzioni quali l'Università di Washington, la CMU, il MIT e la Hong Kong University, rappresenta il framework di riferimento principale all'interno dell'infrastruttura AWS.
\end{itemize}

Questi strumenti di sviluppo assolvono a tre scopi fondamentali:
\begin{enumerate}
    \item Permettere di costruire agevolmente grafi computazionali di grandi dimensioni.
    \item Consentire il calcolo automatico ed efficiente dei gradienti all'interno di tali grafi computazionali (differenziazione automatica).
    \item Eseguire tutte le operazioni in modo altamente ottimizzato su GPU (avvolgendo librerie di basso livello come cuDNN e cuBLAS).
\end{enumerate}

\subsection{Specificazione Dei Grafi Computazionali}

Nei framework moderni come PyTorch, i livelli vengono strutturati in modo modulare definendo in maniera esplicita sia il passo in avanti (\emph{forward pass}) sia il passo all'indietro (\emph{backward pass}). L'esempio seguente mostra l'implementazione personalizzata di una funzione di attivazione \textbf{ReLU}:

\begin{verbatim}
class ReLU(torch.autograd.Function):
    def forward(self, x):
        self.save_for_backward(x)
        return x.clamp(min=0)

    def backward(self, grad_y):
        x, = self.saved_tensors
        grad_input = grad_y.clone()
        grad_input[x < 0] = 0
        return grad_input
\end{verbatim}

La modularità del framework consente poi di comporre l'intera architettura del modello come una sequenza ordinata di livelli, gestendo l'addestramento e l'aggiornamento dei parametri tramite un ciclo iterativo. Di seguito viene riportato un esempio di definizione e addestramento di una rete neurale a due livelli fully-connected usando l'interfaccia sequenziale:

\begin{verbatim}
import torch
from torch.autograd import Variable

N, D_in, H, D_out = 64, 1000, 100, 10
x = Variable(torch.randn(N, D_in))
y = Variable(torch.randn(N, D_out), requires_grad=False)

model = torch.nn.Sequential(
    torch.nn.Linear(D_in, H),
    torch.nn.ReLU(),
    torch.nn.Linear(H, D_out)
)
loss_fn = torch.nn.MSELoss(size_average=False)

learning_rate = 1e-4
for t in range(500):
    y_pred = model(x)
    loss = loss_fn(y_pred, y)

    model.zero_grad()
    loss.backward()

    for param in model.parameters():
        param.data -= learning_rate * param.grad.data
\end{verbatim}

\newpage

\section{Architetture Classiche Per Reti Neurali Convoluzionali}

La progettazione di architetture per le reti neurali convoluzionali (\textbf{CNN}) ha guidato lo stato dell'arte nella computer vision negli ultimi decenni. Vengono qui analizzati i modelli principali attraverso analisi di casi studio dettagliati, partendo dai precursori storici fino ai modelli profondi più complessi:
\begin{itemize}
    \item \textbf{Casi di studio principali}: AlexNet, VGG, GoogLeNet, ResNet.
    \item \textbf{Architetture e varianti correlate}: NiN (Network in Network), Wide ResNet, ResNeXT, Stochastic Depth, DenseNet, FractalNet, SqueezeNet.
\end{itemize}

\subsection{Ripasso Storico: LeNet-5}

Introdotta nel 1998 da Yann LeCun e colleghi, l'architettura \textbf{LeNet-5} rappresenta uno dei primi successi nell'applicazione delle reti convoluzionali per compiti di riconoscimento di cifre scritte a mano.

% [IMMAGINE DEL FILE DI SLIDE CON NOME 11.2 "1 - 26"] Schema dell'architettura classica LeNet-5 con livelli alternati di convoluzione e subsampling — Slide 8 / Pagina 8 del PDF

L'architettura segue una struttura sequenziale basata sull'alternanza di convoluzioni e pooling:
\[
\text{[CONV - POOL - CONV - POOL - FC - FC]}
\]
Le caratteristiche salienti dei livelli di LeNet-5 includono:
\begin{itemize}
    \item I filtri convoluzionali (\textbf{CONV}) hanno una dimensione spaziale pari a $5 \times 5$ e vengono applicati con uno stride pari a $1$.
    \item I livelli di sottocampionamento (\textbf{Subsampling / Pooling}) hanno filtri di dimensione $2 \times 2$ applicati con stride pari a $2$.
\end{itemize}

\newpage

\section{Studio Di Caso: AlexNet}

Sviluppata da Krizhevsky, Sutskever e Hinton nel 2012, \textbf{AlexNet} rappresenta l'architettura che ha segnato la rinascita e l'affermazione su larga scala del deep learning vincendo la competizione ImageNet.

% [IMMAGINE DEL FILE DI SLIDE CON NOME 11.2 "1 - 26"] Diagramma strutturale di AlexNet con la ripartizione dei canali su due GPU — Slide 9 / Pagina 9 del PDF

L'architettura sequenziale e semplificata di AlexNet è costituita dai seguenti passaggi successivi:
\begin{enumerate}
    \item \textbf{INPUT}: Immagini di dimensione $[227 \times 227 \times 3]$.
    \item \textbf{CONV1}: Livello convoluzionale con filtri di grandi dimensioni.
    \item \textbf{MAX POOL1}: Livello di max pooling spaziale.
    \item \textbf{NORM1}: Livello di normalizzazione dei canali.
    \item \textbf{CONV2}: Livello convoluzionale.
    \item \textbf{MAX POOL2}: Livello di max pooling.
    \item \textbf{NORM2}: Livello di normalizzazione.
    \item \textbf{CONV3}, \textbf{CONV4}, \textbf{CONV5}: Tre livelli convoluzionali consecutivi senza pooling intermedio.
    \item \textbf{MAX POOL3}: Livello di pooling finale prima dei moduli fully-connected.
    \item \textbf{FC6}, \textbf{FC7}, \textbf{FC8}: Livelli densamente connessi, di cui l'ultimo restituisce i punteggi delle classi.
\end{enumerate}

\subsection{Analisi Dettagliata Dei Livelli Di AlexNet}

Di seguito si esegue lo studio matematico della propagazione delle dimensioni e del calcolo dei parametri per i primi livelli dell'architettura.

\subsubsection{Livello Convoluzionale CONV1}

Il volume di input ha dimensioni spaziali e di canale pari a:
\[
W_{\text{in}} = 227, \quad H_{\text{in}} = 227, \quad C_{\text{in}} = 3
\]
Su questo input viene applicato un insieme di $F_{\text{num}} = 96$ filtri convoluzionali di dimensione $F \times F = 11 \times 11$, impiegando uno stride $S = 4$ e un padding $P = 0$.

Per calcolare la dimensione spaziale dell'output $W_{\text{out}}$ (e simmetricamente l'altezza $H_{\text{out}}$), si applica la formula generale:
\[
W_{\text{out}} = \frac{W_{\text{in}} - F + 2P}{S} + 1
\]
Sostituendo i valori numerici si ottiene:
\[
W_{\text{out}} = \frac{227 - 11 + 2(0)}{4} + 1 = \frac{216}{4} + 1 = 54 + 1 = 55
\]
Poiché vengono utilizzati $96$ filtri, la profondità del volume risultante è pari a $96$. L'output di questo livello ha pertanto dimensioni complessive pari a:
\[
\boxed{[55 \times 55 \times 96]}
\]

Per determinare il numero totale di parametri addestrabili all'interno di questo livello convoluzionale, si considera che ciascuno dei $96$ filtri ha una dimensione pari a $11 \times 11 \times 3$ (la profondità del filtro deve eguagliare il numero di canali in ingresso). Il calcolo dei parametri dei pesi della rete (escludendo i termini di bias) è dato da:
\[
\text{Parametri} = (11 \times 11 \times 3) \times 96 = 363 \times 96 = 34.848 \approx \mathbf{35\text{K}}
\]

\subsubsection{Livello Di Pooling POOL1}

L'input di questo livello corrisponde all'output del precedente livello convoluzionale, avente dimensioni $[55 \times 55 \times 96]$.
Vengono applicati filtri di max pooling di dimensione $F \times F = 3 \times 3$ con uno stride $S = 2$ e nessun padding ($P = 0$).

La dimensione spaziale di output viene calcolata come:
\[
W_{\text{out}} = \frac{55 - 3 + 2(0)}{2} + 1 = \frac{52}{2} + 1 = 26 + 1 = 27
\]
Poiché l'operazione di pooling viene eseguita in maniera indipendente su ciascun canale e non altera la profondità del volume, il volume di output finale è pari a:
\[
\boxed{[27 \times 27 \times 96]}
\]

Un livello di pooling applica una funzione matematica non parametrica (la selezione del valore massimo in ogni regione locale di dimensione $3 \times 3$). Pertanto, non vi sono pesi o bias da ottimizzare:
\[
\text{Parametri} = \mathbf{0}
\]

\subsection{Sintesi Dell'Intera Architettura AlexNet}

Sulla base delle regole geometriche applicate, le dimensioni spaziali e i canali si propagano lungo tutta la rete secondo la seguente sequenza di livelli:
\begin{itemize}
    \item \textbf{INPUT}: $[227 \times 227 \times 3]$
    \item \textbf{CONV1}: $96$ filtri di dimensione $11 \times 11$ con stride $4$, pad $0 \rightarrow [55 \times 55 \times 96]$
    \item \textbf{MAX POOL1}: filtri di dimensione $3 \times 3$ con stride $2 \rightarrow [27 \times 27 \times 96]$
    \item \textbf{NORM1}: Normalizzazione locale dei canali $\rightarrow [27 \times 27 \times 96]$
    \item \textbf{CONV2}: $256$ filtri di dimensione $5 \times 5$ con stride $1$, pad $2 \rightarrow [27 \times 27 \times 256]$
    \item \textbf{MAX POOL2}: filtri di dimensione $3 \times 3$ con stride $2 \rightarrow [13 \times 13 \times 256]$
    \item \textbf{NORM2}: Normalizzazione locale dei canali $\rightarrow [13 \times 13 \times 256]$
    \item \textbf{CONV3}: $384$ filtri di dimensione $3 \times 3$ con stride $1$, pad $1 \rightarrow [13 \times 13 \times 384]$
    \item \textbf{CONV4}: $384$ filtri di dimensione $3 \times 3$ con stride $1$, pad $1 \rightarrow [13 \times 13 \times 384]$
    \item \textbf{CONV5}: $256$ filtri di dimensione $3 \times 3$ con stride $1$, pad $1 \rightarrow [13 \times 13 \times 256]$
    \item \textbf{MAX POOL3}: filtri di dimensione $3 \times 3$ con stride $2 \rightarrow [6 \times 6 \times 256]$
    \item \textbf{FC6}: Strato fully-connected con $4096$ neuroni
    \item \textbf{FC7}: Strato fully-connected con $4096$ neuroni
    \item \textbf{FC8}: Strato fully-connected con $1000$ neuroni (corrispondenti alle classi di ImageNet)
\end{itemize}

\subsection{Dettagli Implementativi E Addestramento Multi-GPU}

\subsubsection{Limiti Hardware Storici}

Durante lo sviluppo di AlexNet nel 2012, l'addestramento è stato condotto su GPU \textbf{NVIDIA GTX 580}. Questa scheda grafica disponeva di soli \SI{3}{GB} di memoria fisica. Per ovviare a questo limite stringente di memoria e poter addestrare una rete di tali dimensioni, gli autori hanno dovuto suddividere il grafo computazionale della rete su \textbf{due GPU separate}, ripartendo a metà i neuroni (le mappe delle caratteristiche) all'interno di ciascun livello.

% [IMMAGINE DEL FILE DI SLIDE CON NOME 11.2 "1 - 26"] Evidenziazione sul diagramma di AlexNet della ripartizione del volume del primo livello convoluzionale su due flussi paralleli — Slide 19 / Pagina 19 del PDF

Il volume risultante dal primo livello convoluzionale $[55 \times 55 \times 96]$ viene diviso simmetricamente:
\[
[55 \times 55 \times 48] \times 2 \text{ GPU}
\]

Il flusso di comunicazione tra le due GPU segue uno schema preciso a seconda del livello:
\begin{itemize}
    \item \textbf{Livelli a connessione locale (stessa GPU)}: Per i livelli \textbf{CONV1}, \textbf{CONV2}, \textbf{CONV4} e \textbf{CONV5}, le connessioni e le computazioni avvengono esclusivamente all'interno delle mappe delle caratteristiche residenti sulla singola GPU fisica. Non avviene alcuna comunicazione di rete inter-GPU.
    \item \textbf{Livelli a connessione globale (comunicazione cross-GPU)}: Per i livelli \textbf{CONV3}, \textbf{FC6}, \textbf{FC7} e \textbf{FC8}, le connessioni si estendono a tutte le mappe delle caratteristiche del livello precedente. Ciò richiede uno scambio continuo di dati e sincronizzazione tra le memorie delle due schede video.
\end{itemize}

\subsubsection{Iperparametri Di Addestramento E Scelte Tecnologiche}

I dettagli sperimentali e le innovazioni introdotte con AlexNet includono:
\begin{itemize}
    \item \textbf{Attivazione ReLU}: Rappresenta il primo utilizzo di rilievo della funzione di attivazione lineare rettificata, in sostituzione delle tradizionali funzioni sigmoide o tangente iperbolica.
    \item \textbf{Livelli di normalizzazione (Norm)}: Uso di livelli di normalizzazione locale della risposta (LRN), una tecnica che non trova più larga diffusione nelle architetture odierne.
    \item \textbf{Regolarizzazione tramite Dropout}: Applicazione di una probabilità di spegnimento dei neuroni pari a $0.5$ per contrastare l'overfitting nei livelli densamente connessi.
    \item \textbf{Data Augmentation massivo}: Incremento artificiale del dataset tramite traslazioni, specchiature e variazioni di colore.
    \item \textbf{Dettagli dell'ottimizzatore}:
    \begin{itemize}
        \item Dimensione del batch (\emph{batch size}): $128$.
        \item Ottimizzazione tramite \textbf{SGD con Momentum} impostato a $0.9$.
        \item Tasso di apprendimento (\emph{learning rate}) iniziale pari a $10^{-2}$. Riduzione manuale di un fattore $10$ ogni volta che l'accuratezza sul set di validazione raggiunge un plateau.
        \item Regolarizzazione L2 tramite decadimento dei pesi (\emph{weight decay}) pari a $5 \times 10^{-4}$.
    \end{itemize}
    \item \textbf{Ensemble di modelli}: L'utilizzo di un insieme composto da 7 reti convoluzionali differenti ha permesso di ridurre il tasso di errore in top-5 dal $18.2\%$ al $15.4\%$.
\end{itemize}

\newpage

\section{Evoluzione Dei Vincitori ILSVRC: ZFNet E VGGNet}

L'ImageNet Large Scale Visual Recognition Challenge (\textbf{ILSVRC}) costituisce il punto di riferimento per misurare l'evoluzione delle prestazioni dei modelli di classificazione delle immagini.

% [IMMAGINE DEL FILE DI SLIDE CON NOME 11.2 "1 - 26"] Grafico storico delle percentuali di errore del vincitore dell'ImageNet Challenge dal 2010 al 2015, con l'evidenziazione del trend verso reti più profonde — Slide 22 / Pagina 22 del PDF

Il grafico evidenzia la netta discontinuità introdotta dalle CNN a partire dal 2012:
\begin{itemize}
    \item \textbf{2010} (Shallow, metodi classici): Errore pari a $28.2\%$.
    \item \textbf{2011} (Shallow, metodi classici): Errore pari a $25.8\%$.
    \item \textbf{2012} (\textbf{AlexNet}): Primo modello basato su CNN a vincere la competizione. Rete a 8 livelli con errore del $16.4\%$.
    \item \textbf{2013} (\textbf{ZFNet}): Rete a 8 livelli con iperparametri migliorati, che riduce l'errore all' $11.7\%$.
    \item \textbf{2014} (\textbf{VGG} / \textbf{GoogLeNet}): Transizione verso reti significativamente più profonde. VGG (19 livelli) raggiunge il $7.3\%$ di errore, mentre GoogLeNet (22 livelli) vince con il $6.7\%$.
    \item \textbf{2015} (\textbf{ResNet}): Rete estremamente profonda composta da 152 livelli, che abbatte l'errore al $3.57\%$.
\end{itemize}

\subsection{ZFNet: Ottimizzazione Degli Iperparametri}

Sviluppata da Zeiler e Fergus nel 2013, \textbf{ZFNet} mantiene la medesima struttura complessiva a 8 livelli di AlexNet, ma ne ottimizza la configurazione per preservare meglio i dettagli spaziali ed estrarre caratteristiche più discriminanti.

% [IMMAGINE DEL FILE DI SLIDE CON NOME 11.2 "1 - 26"] Diagramma dei blocchi architetturali di ZFNet con l'evidenziazione delle modifiche ai filtri convoluzionali — Slide 24 / Pagina 24 del PDF

Le modifiche sostanziali apportate rispetto ad AlexNet sono le seguenti:
\begin{itemize}
    \item \textbf{Modifica al livello CONV1}: La dimensione del filtro e lo stride vengono ridotti per limitare la perdita di informazione spaziale nei primi passi di elaborazione:
    \[
    \text{AlexNet: } (11 \times 11, \text{ stride } 4) \quad \longrightarrow \quad \text{ZFNet: } \mathbf{(7 \times 7, \text{ stride } 2)}
    \]
    \item \textbf{Aumento del numero di filtri nei livelli intermedi (CONV3, CONV4, CONV5)}: Si aumenta la capacità della rete espandendo il numero di canali estratti nei livelli centrali:
    \[
    \text{AlexNet: } [384, 384, 256] \quad \longrightarrow \quad \text{ZFNet: } \mathbf{[512, 1024, 512]}
    \]
\end{itemize}
Queste modifiche hanno consentito di abbassare il tasso di errore top-5 su ImageNet dal $16.4\%$ di AlexNet all' \textbf{$11.7\%$}.

\subsection{VGGNet: Architetture Profonde Con Filtri Piccoli}

Sviluppata da Simonyan e Zisserman nel 2014, l'architettura \textbf{VGGNet} dimostra che la profondità della rete costituisce un fattore cruciale per le prestazioni della rappresentazione visiva.

% [IMMAGINE DEL FILE DI SLIDE CON NOME 11.2 "1 - 26"] Confronto grafico affiancato delle tre architetture sequenziali AlexNet, VGG16 e VGG19 — Slide 26 / Pagina 26 del PDF

Le principali caratteristiche distintive di VGG rispetto ai suoi predecessori includono:
\begin{itemize}
    \item \textbf{Aumento sostanziale della profondità}: Si passa dagli 8 livelli di AlexNet a versioni nettamente più profonde, note come \textbf{VGG16} (16 livelli con pesi addestrabili) e \textbf{VGG19} (19 livelli).
    \item \textbf{Uso esclusivo di filtri molto piccoli}: Tutti i livelli convoluzionali adottano esclusivamente filtri di dimensione $3 \times 3$, applicati sempre con stride $1$ e padding di bordo pari a $1$. Questa scelta permette di mantenere invariata la risoluzione spaziale del volume a ogni operazione di convoluzione.
    \item \textbf{Standardizzazione del Pooling}: I livelli di riduzione spaziale sono implementati tramite max pooling con filtri di dimensione fissa $2 \times 2$ e stride pari a $2$.
    \item \textbf{Progressione delle prestazioni}: L'approccio sistematico e standardizzato di VGGNet ha permesso di scendere dal $11.7\%$ di errore top-5 di ZFNet (ILSVRC'13) al \textbf{$7.3\%$} nell'edizione ILSVRC'14.
\end{itemize}

\section{Studio Di Caso: VGGNet}

L'architettura \textbf{VGGNet}, introdotta da Simonyan e Zisserman nel 2014, rappresenta una pietra miliare nello sviluppo di reti neurali convoluzionali profonde. Il principio fondamentale alla base del suo design consiste nella semplificazione strutturale ottenuta tramite l'uso metodico di filtri di dimensioni ridotte in reti significativamente più profonde rispetto ai modelli precedenti (come AlexNet).

% [IMMAGINE DEL FILE DI SLIDE CON NOME "27-52"] Confronto strutturale tra AlexNet, VGG16 e VGG19 — Slide 27 / Pagina 1 del PDF

\subsection{L'Efficacia Dei Filtri Convoluzionali Piccoli ($3 \times 3$)}

Una delle scelte di design più rilevanti di VGGNet riguarda la sostituzione dei filtri convoluzionali di grandi dimensioni (quali quelli $7 \times 7$ o $11 \times 11$ impiegati in AlexNet) con una sequenza di filtri più piccoli da $3 \times 3$.

Un quesito fondamentale nello studio di questa architettura riguarda il motivo per cui l'uso di filtri più piccoli risulti vantaggioso. Per comprenderlo, è necessario analizzare il concetto di \textbf{campo ricettivo effettivo} (\emph{effective receptive field}).

Si consideri una pila (\emph{stack}) di tre strati convoluzionali consecutivi con filtri di dimensione $3 \times 3$ applicati con uno stride pari a $1$. Ciascuno strato opera come segue:
\begin{enumerate}
    \item Il primo strato convoluzionale da $3 \times 3$ ha un campo ricettivo effettivo di dimensione $3 \times 3$ rispetto all'input originario.
    \item Il secondo strato convoluzionale da $3 \times 3$, operando sull'output del primo, estende la sua influenza spaziale sull'input originario fino a una dimensione di $5 \times 5$.
    \item Il terzo strato convoluzionale da $3 \times 3$, elaborando l'output del secondo, estende ulteriormente la propria finestra visiva sull'input iniziale, raggiungendo una dimensione complessiva pari a:
    \[
    \boxed{7 \times 7}
    \]
\end{enumerate}

Di conseguenza, una sequenza ordinata di tre livelli convoluzionali da $3 \times 3$ (con stride $1$) possiede esattamente lo stesso campo ricettivo effettivo di un singolo livello convoluzionale da $7 \times 7$. Questa equivalenza spaziale comporta tuttavia due vantaggi computazionali cruciali:

\begin{enumerate}
    \item \textbf{Maggiore profondità e non-linearità}: L'adozione di tre strati convoluzionali in cascata anziché uno solo permette di inserire tre funzioni di attivazione non lineari (ad esempio la ReLU) lungo il percorso computazionale, anziché una singola attivazione. Ciò aumenta significativamente la capacità rappresentativa e la capacità della rete di apprendere funzioni altamente non lineari.
    \item \textbf{Drastica riduzione del numero di parametri}: Si assuma che ciascuno strato gestisca un numero di canali di input e di output pari a $C$.
    \begin{itemize}
        \item Un singolo livello convoluzionale con filtri $7 \times 7$ richiede un numero di parametri (escludendo i bias) pari a:
        \[
        7^2 \times C^2 = 49 C^2
        \]
        \item Tre livelli convoluzionali in cascata con filtri $3 \times 3$ richiedono complessivamente:
        \[
        3 \times (3^2 \times C^2) = 27 C^2
        \]
    \end{itemize}
    Il confronto evidenzia che l'approccio con filtri più piccoli consente di ridurre il numero totale dei parametri di circa il $45\%$, limitando il rischio di \emph{overfitting} pur mantenendo la medesima capacità di estrazione spaziale dei pattern.
\end{enumerate}

\subsection{Analisi Quantitativa Di Memoria E Parametri In VGG16}

Per comprendere l'impatto di queste scelte progettuali sui requisiti hardware, si analizza dettagliatamente la propagazione dei volumi, l'occupazione di memoria e il numero di parametri (senza considerare i termini di bias) per ciascun livello del modello di riferimento \textbf{VGG16}:

% [IMMAGINE DEL FILE DI SLIDE CON NOME "27-52"] Struttura a blocchi di VGG16 con analisi di memoria e parametri — Slide 31 / Pagina 5 del PDF

\begin{enumerate}[noitemsep]
    \item \textbf{INPUT}: $[224 \times 224 \times 3]$
    \begin{itemize}
        \item Memoria: $224 \times 224 \times 3 \approx 150\text{K}$ valori floating-point.
        \item Parametri: $0$.
    \end{itemize}
    \item \textbf{CONV3-64} (\emph{conv1-1}): $[224 \times 224 \times 64]$
    \begin{itemize}
        \item Memoria: $224 \times 224 \times 64 \approx 3.2\text{M}$ valori.
        \item Parametri: $(3 \times 3 \times 3) \times 64 = 1.728$.
    \end{itemize}
    \item \textbf{CONV3-64} (\emph{conv1-2}): $[224 \times 224 \times 64]$
    \begin{itemize}
        \item Memoria: $224 \times 224 \times 64 \approx 3.2\text{M}$ valori.
        \item Parametri: $(3 \times 3 \times 64) \times 64 = 36.864$.
    \end{itemize}
    \item \textbf{POOL2} (\emph{pool1}): $[112 \times 112 \times 64]$
    \begin{itemize}
        \item Memoria: $112 \times 112 \times 64 \approx 800\text{K}$ valori.
        \item Parametri: $0$.
    \end{itemize}
    \item \textbf{CONV3-128} (\emph{conv2-1}): $[112 \times 112 \times 128]$
    \begin{itemize}
        \item Memoria: $112 \times 112 \times 128 \approx 1.6\text{M}$ valori.
        \item Parametri: $(3 \times 3 \times 64) \times 128 = 73.728$.
    \end{itemize}
    \item \textbf{CONV3-128} (\emph{conv2-2}): $[112 \times 112 \times 128]$
    \begin{itemize}
        \item Memoria: $112 \times 112 \times 128 \approx 1.6\text{M}$ valori.
        \item Parametri: $(3 \times 3 \times 128) \times 128 = 147.456$.
    \end{itemize}
    \item \textbf{POOL2} (\emph{pool2}): $[56 \times 56 \times 128]$
    \begin{itemize}
        \item Memoria: $56 \times 56 \times 128 \approx 400\text{K}$ valori.
        \item Parametri: $0$.
    \end{itemize}
    \item \textbf{CONV3-256} (\emph{conv3-1}): $[56 \times 56 \times 256]$
    \begin{itemize}
        \item Memoria: $56 \times 56 \times 256 \approx 800\text{K}$ valori.
        \item Parametri: $(3 \times 3 \times 128) \times 256 = 294.912$.
    \end{itemize}
    \item \textbf{CONV3-256} (\emph{conv3-2}): $[56 \times 56 \times 256]$
    \begin{itemize}
        \item Memoria: $56 \times 56 \times 256 \approx 800\text{K}$ valori.
        \item Parametri: $(3 \times 3 \times 256) \times 256 = 589.824$.
    \end{itemize}
    \item \textbf{CONV3-256} (\emph{conv3-3}): $[56 \times 56 \times 256]$
    \begin{itemize}
        \item Memoria: $56 \times 56 \times 256 \approx 800\text{K}$ valori.
        \item Parametri: $(3 \times 3 \times 256) \times 256 = 589.824$.
    \end{itemize}
    \item \textbf{POOL2} (\emph{pool3}): $[28 \times 28 \times 256]$
    \begin{itemize}
        \item Memoria: $28 \times 28 \times 256 \approx 200\text{K}$ valori.
        \item Parametri: $0$.
    \end{itemize}
    \item \textbf{CONV3-512} (\emph{conv4-1}): $[28 \times 28 \times 512]$
    \begin{itemize}
        \item Memoria: $28 \times 28 \times 512 \approx 400\text{K}$ valori.
        \item Parametri: $(3 \times 3 \times 256) \times 512 = 1.179.648$.
    \end{itemize}
    \item \textbf{CONV3-512} (\emph{conv4-2}): $[28 \times 28 \times 512]$
    \begin{itemize}
        \item Memoria: $28 \times 28 \times 512 \approx 400\text{K}$ valori.
        \item Parametri: $(3 \times 3 \times 512) \times 512 = 2.359.296$.
    \end{itemize}
    \item \textbf{CONV3-512} (\emph{conv4-3}): $[28 \times 28 \times 512]$
    \begin{itemize}
        \item Memoria: $28 \times 28 \times 512 \approx 400\text{K}$ valori.
        \item Parametri: $(3 \times 3 \times 512) \times 512 = 2.359.296$.
    \end{itemize}
    \item \textbf{POOL2} (\emph{pool4}): $[14 \times 14 \times 512]$
    \begin{itemize}
        \item Memoria: $14 \times 14 \times 512 \approx 100\text{K}$ valori.
        \item Parametri: $0$.
    \end{itemize}
    \item \textbf{CONV3-512} (\emph{conv5-1}): $[14 \times 14 \times 512]$
    \begin{itemize}
        \item Memoria: $14 \times 14 \times 512 \approx 100\text{K}$ valori.
        \item Parametri: $(3 \times 3 \times 512) \times 512 = 2.359.296$.
    \end{itemize}
    \item \textbf{CONV3-512} (\emph{conv5-2}): $[14 \times 14 \times 512]$
    \begin{itemize}
        \item Memoria: $14 \times 14 \times 512 \approx 100\text{K}$ valori.
        \item Parametri: $(3 \times 3 \times 512) \times 512 = 2.359.296$.
    \end{itemize}
    \item \textbf{CONV3-512} (\emph{conv5-3}): $[14 \times 14 \times 512]$
    \begin{itemize}
        \item Memoria: $14 \times 14 \times 512 \approx 100\text{K}$ valori.
        \item Parametri: $(3 \times 3 \times 512) \times 512 = 2.359.296$.
    \end{itemize}
    \item \textbf{POOL2} (\emph{pool5}): $[7 \times 7 \times 512]$
    \begin{itemize}
        \item Memoria: $7 \times 7 \times 512 \approx 25\text{K}$ valori.
        \item Parametri: $0$.
    \end{itemize}
    \item \textbf{FC} (\emph{fc6}): $[1 \times 1 \times 4096]$
    \begin{itemize}
        \item Memoria: $4096$ valori.
        \item Parametri: $7 \times 7 \times 512 \times 4096 = 102.760.448$.
    \end{itemize}
    \item \textbf{FC} (\emph{fc7}): $[1 \times 1 \times 4096]$
    \begin{itemize}
        \item Memoria: $4096$ valori.
        \item Parametri: $4096 \times 4096 = 16.777.216$.
    \end{itemize}
    \item \textbf{FC} (\emph{fc8}): $[1 \times 1 \times 1000]$
    \begin{itemize}
        \item Memoria: $1000$ valori.
        \item Parametri: $4096 \times 1000 = 4.096.000$.
    \end{itemize}
\end{enumerate}

% [IMMAGINE DEL FILE DI SLIDE CON NOME "27-52"] Nomenclatura comune dei blocchi convoluzionali e fully-connected in VGG16 — Slide 34 / Pagina 8 del PDF

L'analisi quantitativa complessiva della rete rileva i seguenti requisiti aggregati:
\begin{align*}
\text{Memoria Totale} &\approx 24\text{M di valori floating-point} \\
\text{Parametri Totali} &\approx 138\text{M di parametri}
\end{align*}

Se si considera una codifica standard a single-precision (4 byte per valore reale), la memoria occupata durante il solo passo di propagazione in avanti (\emph{forward pass}) risulta essere:
\[
24\text{M} \times 4 \text{ bytes} \approx \mathbf{96\text{MB / immagine}}
\]
Durante la fase di addestramento, il calcolo dei gradienti per la retropropagazione raddoppia approssimativamente questa richiesta di memoria, portandola a circa $\approx \mathbf{192\text{MB / immagine}}$.

Dall'analisi della distribuzione spaziale di tali requisiti, emergono due considerazioni asimmetriche fondamentali:
\begin{itemize}
    \item \textbf{Distribuzione della Memoria}: La stragrande maggioranza della memoria di attivazione viene consumata nei \textbf{primi livelli convoluzionali} (\emph{early CONV}), in cui i volumi preservano dimensioni spaziali elevate (ad esempio, conv1-1 e conv1-2 richiedono ciascuno circa $3.2\text{M}$ di valori).
    \item \textbf{Distribuzione dei Parametri}: La quasi totalità dei parametri della rete si concentra negli \textbf{ultimi strati fully-connected} (\emph{late FC}). In particolare, il passaggio dal volume convoluzionale finale $[7 \times 7 \times 512]$ al primo livello denso fc6 comporta da solo l'introduzione di oltre $102\text{M}$ di parametri (ossia circa il $74\%$ dell'intera rete).
\end{itemize}

\subsection{Prestazioni E Dettagli Di Addestramento Di VGGNet}

L'architettura VGGNet si caratterizza per una serie di risultati competitivi e metodologie sperimentali:
\begin{itemize}
    \item \textbf{Risultati ILSVRC'14}: Il modello ha ottenuto il secondo posto assoluto nel task di classificazione e il primo posto nel task di localizzazione spaziale degli oggetti.
    \item \textbf{Protocollo di Addestramento}: La procedura di ottimizzazione segue fedelmente quella descritta per AlexNet (Krizhevsky 2012).
    \item \textbf{Eliminazione della Normalizzazione}: A differenza di AlexNet, in VGGNet si omette l'uso dello strato di Normalizzazione della Risposta Locale (LRN), in quanto non si sono riscontrati benefici prestazionali empirici a fronte di un costo computazionale aggiuntivo.
    \item \textbf{Scelta del Modello (VGG16 vs VGG19)}: Sebbene l'architettura VGG19 offra un'accuratezza marginalmente superiore, essa comporta un notevole incremento del consumo di memoria a causa dei tre livelli convoluzionali aggiuntivi. Per questa ragione, nella pratica applicativa si preferisce spesso VGG16.
    \item \textbf{Uso di Ensemble}: L'integrazione probabilistica di un insieme di più modelli convoluzionali VGG consente di raggiungere i livelli di errore minimi registrati nella competizione.
    \item \textbf{Generalizzazione dei Feature Vector}: Le caratteristiche numeriche estratte dallo strato fc7 dimostrano eccellenti proprietà di trasferibilità di conoscenza, generalizzando efficacemente su una vasta gamma di task computazionali esterni.
\end{itemize}

\newpage

\section{Studio Di Caso: GoogLeNet}

Presentata da Szegedy et al. nel 2014, l'architettura \textbf{GoogLeNet} (nota anche come Inception v1) introduce un cambiamento di paradigma focalizzato sull'efficienza computazionale del grafo per consentire la scalabilità in profondità senza far esplodere la complessità di calcolo.

% [IMMAGINE DEL FILE DI SLIDE CON NOME "27-52"] Struttura macroscopica di GoogLeNet e schema logico del modulo Inception — Slide 37 / Pagina 11 del PDF

La filosofia di progettazione si fonda sulle seguenti specifiche:
\begin{itemize}
    \item \textbf{Profondità elevata}: L'architettura si articola su ben $22$ strati computazionali organizzati in blocchi strutturati.
    \item \textbf{Modulo Inception}: Introduzione di una topologia locale efficiente che racchiude operazioni parallele all'interno di un unico macro-blocco.
    \item \textbf{Assenza di livelli Fully-Connected (FC)}: Per contenere l'esplosione dei parametri evidenziata nel modello VGG, GoogLeNet elimina i tradizionali strati finali densamente connessi, sostituendoli con un livello di Global Average Pooling.
    \item \textbf{Efficienza nei Parametri}: La rete conta solo \textbf{5 milioni di parametri}, un quantitativo circa 12 volte inferiore a quello richiesto da AlexNet, pur garantendo una profondità quasi tripla.
    \item \textbf{Vittoria a ILSVRC'14}: GoogLeNet ha conquistato il primo posto nel task di classificazione con un tasso di errore top-5 pari al \textbf{$6.7\%$}.
\end{itemize}

Il nucleo teorico di GoogLeNet risiede nel concetto di \textbf{Modulo Inception}: l'obiettivo è progettare una topologia locale efficace (definita come una sub-rete o \emph{network within a network}) e successivamente impilare questi moduli l'uno sull'altro in modo sequenziale.

\subsection{Analisi Del Modulo Inception Naive E I Suoi Limiti}

La versione preliminare e non ottimizzata del modulo Inception (chiamata \textbf{Naive Inception Module}) opera applicando in parallelo diverse trasformazioni locali sul volume di input proveniente dallo strato precedente:

% [IMMAGINE DEL FILE DI SLIDE CON NOME "27-52"] Rappresentazione concettuale del modulo Naive Inception con rami paralleli — Slide 39 / Pagina 13 del PDF

\begin{itemize}
    \item Convoluzioni con campi ricettivi multipli per estrarre informazioni a diverse scale spaziali: $1 \times 1$, $3 \times 3$, $5 \times 5$.
    \item Un'operazione di pooling locale parallela: Max Pooling $3 \times 3$.
\end{itemize}
Tutti i flussi in uscita da queste quattro operazioni parallele vengono infine concatenati lungo la dimensione del canale (profondità).

\subsubsection{Esempio Numerico Di Criticità Computazionale}

Al fine di quantificare la complessità computazionale di questa architettura naive, si consideri uno scenario esemplificativo in cui l'input fornito al modulo Inception sia un volume di dimensioni:
\[
W_{\text{in}} = 28, \quad H_{\text{in}} = 28, \quad C_{\text{in}} = 256
\]
Si applichino i seguenti operatori paralleli, configurati in modo da preservare le dimensioni spaziali dell'input tramite opportuni parametri di padding e stride:
\begin{itemize}
    \item \textbf{Ramo 1}: Convoluzione $1 \times 1$ con $128$ filtri.
    \[
    \text{Output 1} = [28 \times 28 \times 128]
    \]
    \item \textbf{Ramo 2}: Convoluzione $3 \times 3$ con $192$ filtri (con $P = 1$, $S = 1$).
    \[
    \text{Output 2} = [28 \times 28 \times 192]
    \]
    \item \textbf{Ramo 3}: Convoluzione $5 \times 5$ con $96$ filtri (con $P = 2$, $S = 1$).
    \[
    \text{Output 3} = [28 \times 28 \times 96]
    \]
    \item \textbf{Ramo 4}: Max Pooling $3 \times 3$ con stride $1$ e padding $1$. Poiché l'operazione di pooling preserva la profondità originale del canale:
    \[
    \text{Output 4} = [28 \times 28 \times 256]
    \]
\end{itemize}

Al termine dei rami paralleli, il volume di output finale si ottiene tramite la concatenazione dei canali di ciascun ramo:
\begin{align*}
W_{\text{out}} &= 28 \\
H_{\text{out}} &= 28 \\
C_{\text{out}} &= 128 + 192 + 96 + 256 = 672
\end{align*}
Il volume risultante ha pertanto dimensioni spaziali e di canale pari a:
\[
\boxed{[28 \times 28 \times 672]}
\]
Il numero totale di valori floating-point memorizzati in questo singolo livello ammonta a $28 \times 28 \times 672 \approx 527\text{K}$.

Per valutare il costo di calcolo di questa configurazione naive, si determina il numero complessivo di operazioni convoluzionali (moltiplicazioni e accumuli, o \emph{Ops}) richieste. La formula generale per il calcolo delle operazioni in uno strato convoluzionale è definita come:
\[
\text{Ops} = \underbrace{(W_{\text{out}} \times H_{\text{out}} \times C_{\text{out}})}_{\text{Numero di neuroni in uscita}} \times \underbrace{(F \times F \times C_{\text{in}})}_{\text{Dimensione spaziale e profondità del filtro}}
\]

Applicando la formula ai singoli rami della struttura parallela si ottiene:
\begin{itemize}
    \item \textbf{Costo della convoluzione $1 \times 1$}:
    \[
    \text{Ops}_{[1 \times 1]} = (28 \times 28 \times 128) \times (1 \times 1 \times 256) \approx \mathbf{25.7\text{M ops}}
    \]
    \item \textbf{Costo della convoluzione $3 \times 3$}:
    \[
    \text{Ops}_{[3 \times 3]} = (28 \times 28 \times 192) \times (3 \times 3 \times 256) \approx \mathbf{346.8\text{M ops}}
    \]
    \item \textbf{Costo della convoluzione $5 \times 5$}:
    \[
    \text{Ops}_{[5 \times 5]} = (28 \times 28 \times 96) \times (5 \times 5 \times 256) \approx \mathbf{481.7\text{M ops}}
    \]
\end{itemize}

Sommando i contributi dei singoli rami, si ottiene il costo di calcolo complessivo del modulo:
\[
\text{Ops}_{\text{Totale}} = 25.7\text{M} + 346.8\text{M} + 481.7\text{M} \approx \mathbf{854\text{M ops}}
\]

Questo valore evidenzia l'insostenibilità computazionale del modulo Inception Naive. Oltre all'elevato numero di operazioni richieste, un ulteriore limite strutturale risiede nel fatto che lo strato di pooling preserva la profondità del canale d'ingresso ($256$). Di conseguenza, la profondità del volume concatenato finale può solo aumentare di livello in livello, portando rapidamente a un collo di bottiglia computazionale e di memoria.

\subsection{La Soluzione: I Livelli Bottleneck Con Convoluzioni $1 \times 1$}

Per risolvere le criticità legate alla complessità computazionale e all'esplosione dei canali, GoogLeNet introduce i cosiddetti strati \textbf{bottleneck} basati sull'uso di convoluzioni $1 \times 1$.

\subsubsection{Funzionamento Delle Convoluzioni $1 \times 1$}

Si consideri un volume tridimensionale di ingresso di dimensioni $[56 \times 56 \times 64]$. Si applichi a questo volume un livello convoluzionale $1 \times 1$ configurato con $32$ filtri.

% [IMMAGINE DEL FILE DI SLIDE CON NOME "27-52"] Effetto dimensionale di una convoluzione 1x1 su un volume tridimensionale — Slide 51 / Pagina 25 del PDF

Ciascuno dei $32$ filtri ha dimensioni effettive pari a $1 \times 1 \times 64$ e calcola un prodotto scalare a 64 dimensioni su ogni singola cella spaziale del volume di input. La trasformazione lineare risultante genera un volume spaziale avente dimensioni pari a:
\[
\boxed{[56 \times 56 \times 32]}
\]

Questa operazione offre due importanti vantaggi di tipo geometrico e rappresentativo:
\begin{itemize}
    \item \textbf{Conservazione della risoluzione spaziale}: Le dimensioni di larghezza e altezza ($56 \times 56$) rimangono inalterate.
    \item \textbf{Riduzione controllata della profondità}: Consente di comprimere la profondità del volume (in questo caso dimezzando il numero di canali da 64 a 32).
    \item \textbf{Proiezione lineare e combinazione}: Opera una proiezione lineare dello spazio dei canali in una dimensione inferiore, consentendo alla rete di apprendere combinazioni ottimali e non lineari (se seguite da funzioni di attivazione) delle mappe delle caratteristiche estratte in precedenza.
\end{itemize}

\section{Caso Di Studio: GoogLeNet}

L'architettura \textbf{GoogLeNet}, introdotta da Szegedy et al. nel 2014, rappresenta uno dei principali punti di svolta nell'evoluzione delle reti neurali convoluzionali profonde, ponendo una forte enfasi sull'efficienza computazionale oltre che sulla capacità di rappresentazione del modello.

\subsection{L'Inception Module E Il Ruolo Dei Bottleneck}

La pietra angolare di GoogLeNet è il modulo \textbf{Inception}, il quale si basa sull'idea di progettare una topologia locale efficace (una sorta di \emph{"network within a network"}) da ripetere e impilare per formare l'architettura complessiva.

\subsubsection{Modulo Inception Naïve}

Nella sua formulazione più intuitiva, definita \textbf{Inception Naïve}, il modulo riceve l'output del livello precedente (\emph{Previous Layer}) e vi applica in parallelo diverse operazioni convoluzionali e di pooling:
\begin{itemize}
    \item Convoluzioni $1\times1$;
    \item Convoluzioni $3\times3$;
    \item Convoluzioni $5\times5$;
    \item Max pooling $3\times3$ (eseguito con stride unitario e opportuno padding per preservare la risoluzione spaziale).
\end{itemize}
Tutti i feature map prodotti in parallelo da queste ramificazioni vengono infine uniti lungo la dimensione dei canali tramite un'operazione di concatenazione dei filtri (\textbf{Filter concatenation}). Ciò permette alla rete di elaborare le informazioni a diverse scale spaziali all'interno dello stesso livello.

Tuttavia, questo approccio presenta un grave limite computazionale: le convoluzioni $3\times3$ e $5\times5$ applicate direttamente su canali di input molto profondi richiedono una quantità insostenibile di operazioni. Inoltre, lo strato di pooling conserva la profondità dell'input, provocando un aumento incontrollato dei canali a ogni modulo successivo.

% [IMMAGINE DEL FILE DI SLIDE CON NOME "53-78"] Schema a blocchi del modulo Inception Naïve rispetto al modulo Inception con riduzione dimensionale — Slide 53 / Pagina 1 del PDF

\subsubsection{Modulo Inception Con Riduzione Dimensionale (Bottleneck)}

Per risolvere il problema della complessità computazionale, viene introdotta la versione con \textbf{riduzione dimensionale}. Essa prevede l'inserimento di convoluzioni $1\times1$ con funzione di \textbf{"bottleneck"} (collo di bottiglia) posizionate strategicamente:
\begin{itemize}
    \item Prima delle convoluzioni $3\times3$ e $5\times5$, per proiettare i canali di input in uno spazio di dimensione inferiore prima di eseguire i calcoli con i filtri spaziali più ampi;
    \item Dopo il max pooling $3\times3$, in modo da ridurre la profondità del canale della mappa risultante prima che questa confluisca nella concatenazione finale.
\end{itemize}

% [IMMAGINE DEL FILE DI SLIDE CON NOME "53-78"] Evidenziazione in rosso dei blocchi di convoluzione 1x1 impiegati come bottleneck nel modulo Inception — Slide 54 / Pagina 2 del PDF

\subsection{Analisi Delle Operazioni Di Convoluzione (Conv Ops)}

Per comprendere quantitativamente il risparmio computazionale introdotto dai livelli di bottleneck, analizziamo un esempio numerico dettagliato in cui l'input del modulo ha dimensioni spaziali pari a $28\times28$ e una profondità di $256$ canali (ovvero un tensore $28\times28\times256$).

La formula per calcolare il numero di operazioni (moltiplicazioni e accumulazioni) richieste da uno strato convoluzionale con filtri di dimensione $K \times K$, input di profondità $C_{\text{in}}$ e mappa di output di dimensioni $H \times W \times C_{\text{out}}$ è:
\[ \text{Ops} = H \times W \times C_{\text{out}} \times K \times K \times C_{\text{in}} \]

Utilizzando la medesima struttura parallela della versione naïve e integrando i bottleneck da $64$ filtri, le operazioni necessarie per ciascun blocco del modulo Inception con riduzione dimensionale si sviluppano come segue:

\begin{align*}
\text{Ops}_{[1\times1 \text{ conv, } 64 \text{ prima di } 3\times3]} &= 28 \times 28 \times 64 \times 1 \times 1 \times 256 \approx 12.85\text{M} \\
\text{Ops}_{[1\times1 \text{ conv, } 64 \text{ prima di } 5\times5]} &= 28 \times 28 \times 64 \times 1 \times 1 \times 256 \approx 12.85\text{M} \\
\text{Ops}_{[1\times1 \text{ conv, } 128 \text{ diretto}]} &= 28 \times 28 \times 128 \times 1 \times 1 \times 256 \approx 25.69\text{M} \\
\text{Ops}_{[3\times3 \text{ conv, } 192]} &= 28 \times 28 \times 192 \times 3 \times 3 \times 64 \approx 86.70\text{M} \\
\text{Ops}_{[5\times5 \text{ conv, } 96]} &= 28 \times 28 \times 96 \times 5 \times 5 \times 64 \approx 120.42\text{M} \\
\text{Ops}_{[1\times1 \text{ conv, } 64 \text{ dopo pool}]} &= 28 \times 28 \times 64 \times 1 \times 1 \times 256 \approx 12.85\text{M}
\end{align*}

Sommando tutti i contributi sopra descritti si ricava il costo computazionale totale dell'Inception module ottimizzato:
\[ \boxed{\text{Totale Bottleneck} = 358\text{M ops}} \]

In assenza di bottleneck (modulo Naïve), l'applicazione diretta delle convoluzioni $3\times3$ e $5\times5$ sui $256$ canali di input comporterebbe un numero di operazioni notevolmente superiore:
\begin{align*}
\text{Ops}_{[1\times1 \text{ conv, } 128 \text{ diretto}]} &= 28 \times 28 \times 128 \times 1 \times 1 \times 256 \approx 25.69\text{M} \\
\text{Ops}_{[3\times3 \text{ conv, } 192]} &= 28 \times 28 \times 192 \times 3 \times 3 \times 256 \approx 346.82\text{M} \\
\text{Ops}_{[5\times5 \text{ conv, } 96]} &= 28 \times 28 \times 96 \times 5 \times 5 \times 256 \approx 481.69\text{M}
\end{align*}
Il costo complessivo della versione Naïve ammonterebbe a:
\[ \boxed{\text{Totale Naïve} = 854\text{M ops}} \]

L'integrazione dei bottleneck consente pertanto di ridurre le operazioni totali da $854\text{M}$ a $358\text{M}$ (un risparmio computazionale di circa il $58\%$). Inoltre, limitando lo spessore dell'output del ramo di pooling tramite convoluzione $1\times1$, la profondità finale del tensore concatenato si riduce, prevenendone la crescita esplosiva. Nel nostro caso, la concatenazione produce un tensore di dimensioni:
\[ 28 \times 28 \times (128 + 192 + 96 + 64) = 28 \times 28 \times 480 \]

% [IMMAGINE DEL FILE DI SLIDE CON NOME "53-78"] Esempio numerico con calcolo dettagliato delle operazioni di convoluzione nel modulo con riduzione dimensionale — Slide 55 / Pagina 3 del PDF

\subsection{Architettura GoogLeNet Completa}

L'architettura complessiva di GoogLeNet si ottiene impilando una serie di moduli Inception con riduzione della dimensionalità l'uno sull'altro.

% [IMMAGINE DEL FILE DI SLIDE CON NOME "53-78"] Moduli Inception con riduzione dimensionale impilati uno sull'altro — Slide 56 / Pagina 4 del PDF

La rete si compone di tre macro-sezioni principali:

\begin{enumerate}
    \item \textbf{Stem Network (Rete Di Supporto Iniziale)}: Posta all'ingresso del modello, elabora l'immagine originale tramite un'alternanza di livelli convoluzionali classici e pooling. La struttura segue lo schema: \textbf{Conv-Pool-2x Conv-Pool}.

    % [IMMAGINE DEL FILE DI SLIDE CON NOME "53-78"] Diagramma dell'architettura completa di GoogLeNet con evidenziata la Stem Network iniziale — Slide 57 / Pagina 5 del PDF

    \item \textbf{Stacked Inception Modules}: Costituisce il corpo centrale del modello ed è composto da una sequenza di moduli Inception collegati in cascata.

    % [IMMAGINE DEL FILE DI SLIDE CON NOME "53-78"] Diagramma dell'architettura di GoogLeNet con evidenziati i moduli Inception impilati nel corpo centrale — Slide 58 / Pagina 6 del PDF

    \item \textbf{Classifier Output}: Situato al termine della rete, genera la classificazione finale. Una peculiarità fondamentale di questa sezione è la \textbf{rimozione dei costosi strati Fully Connected (FC)} tipici delle reti precedenti, sostituiti da un livello di \emph{Global Average Pooling} che riduce drasticamente il numero totale di parametri.

    % [IMMAGINE DEL FILE DI SLIDE CON NOME "53-78"] Diagramma dell'architettura di GoogLeNet con evidenziato l'output del classificatore finale — Slide 59 / Pagina 7 del PDF
    % [IMMAGINE DEL FILE DI SLIDE CON NOME "53-78"] Diagramma dell'architettura di GoogLeNet che mostra la rimozione dei costosi livelli Fully Connected a favore dell'efficienza — Slide 60 / Pagina 8 del PDF
\end{enumerate}

\subsubsection{Classificatori Ausiliari}

Dato il notevole sviluppo in profondità di GoogLeNet, sorge il rischio che il gradiente si attenui progressivamente durante la backpropagation, impedendo l'efficace addestramento dei livelli iniziali (problema del gradiente svanente o \emph{vanishing gradient}).

Per superare questo limite, sono stati introdotti due **classificatori ausiliari** (\emph{Auxiliary classification outputs}) nelle porzioni intermedie della rete. Ciascuno di essi è strutturato come segue:
\[ \text{AvgPool} \rightarrow 1\times1 \text{ Conv} \rightarrow \text{FC} \rightarrow \text{FC} \rightarrow \text{Softmax} \]
Durante la fase di addestramento, questi rami calcolano delle loss intermedie che vengono sommate alla loss finale, iniettando gradiente aggiuntivo direttamente nei livelli inferiori per stabilizzare l'ottimizzazione. In fase di inferenza, i rami ausiliari vengono rimossi.

% [IMMAGINE DEL FILE DI SLIDE CON NOME "53-78"] Diagramma di GoogLeNet con evidenziazione dei rami di classificazione ausiliari per l'iniezione del gradiente — Slide 61 / Pagina 9 del PDF

In totale, GoogLeNet conta \textbf{22 livelli dotati di pesi} (includendo ciascuna delle ramificazioni parallele interne ai moduli Inception).

% [IMMAGINE DEL FILE DI SLIDE CON NOME "53-78"] Diagramma di GoogLeNet con indicazione del numero totale di strati dotati di pesi (22 livelli) — Slide 62 / Pagina 10 del PDF

\subsection{Riepilogo Delle Proprietà Di GoogLeNet}

GoogLeNet riesce a coniugare una notevole profondità d'analisi con una spiccata efficienza computazionale:
\begin{itemize}
    \item $22$ livelli con parametri addestrabili;
    \item Utilizzo di moduli Inception altamente ottimizzati;
    \item Rimozione quasi totale dei livelli Fully Connected;
    \item Riduzione dei parametri pari a \textbf{$12\times$ in meno rispetto ad AlexNet};
    \item Vincitore della categoria di classificazione di \textbf{ILSVRC'14} con un errore top-5 del \textbf{$6.7\%$}.
\end{itemize}

% [IMMAGINE DEL FILE DI SLIDE CON NOME "53-78"] Tabella riepilogativa delle caratteristiche chiave e dei risultati ottenuti da GoogLeNet — Slide 63 / Pagina 11 del PDF

\newpage

\section{Caso Di Studio: ResNet}

L'analisi dell'evoluzione dei modelli vincitori della competizione ImageNet (ILSVRC) evidenzia un chiaro trend noto come \emph{"Revolution of Depth"}. L'errore top-5 decresce costantemente all'aumentare del numero di strati dell'architettura: si passa dagli 8 livelli di AlexNet ($16.4\%$) e i 19 di VGG ($7.3\%$), fino ad arrivare ai ben \textbf{152 livelli} di **ResNet** nel 2015, che abbatte l'errore al \textbf{$3.57\%$}.

% [IMMAGINE DEL FILE DI SLIDE CON NOME "53-78"] Grafico "Revolution of Depth" che mostra la diminuzione dell'errore su ImageNet all'aumentare della profondità delle reti — Slide 64 / Pagina 12 del PDF

L'architettura **ResNet** (He et al., 2015) rende possibile l'addestramento di reti di questa portata introducendo le **connessioni residuali** (\emph{residual connections}). Grazie a questa intuizione, il modello ha dominato tutte le principali competizioni di classificazione e rilevamento di ILSVRC'15 e COCO'15.

% [IMMAGINE DEL FILE DI SLIDE CON NOME "53-78"] Schema a blocchi elementare di un blocco residuale e vista d'insieme del modello ResNet — Slide 65 / Pagina 13 del PDF

\subsection{Il Problema Della Degradazione}

Cosa accade se proviamo ad addestrare reti convoluzionali tradizionali (prive di connessioni residuali, definite reti \emph{"plain"}) aumentando semplicemente il numero di strati?

% [IMMAGINE DEL FILE DI SLIDE CON NOME "53-78"] Domanda introduttiva sugli effetti dello stacking continuo di strati in reti "plain" — Slide 66 / Pagina 14 del PDF

Empiricamente si osserva che, superata una certa soglia di profondità, le prestazioni della rete iniziano a peggiorare rapidamente. Dal confronto tra una rete plain a 20 livelli e una a 56 livelli emergono i seguenti comportamenti:
\begin{itemize}
    \item La rete a 56 livelli presenta un \textbf{errore di addestramento (training error)} sistematicamente superiore rispetto alla rete a 20 livelli;
    \item Di conseguenza, anche l'errore di test (\emph{test error}) risulta peggiore nella rete più profonda.
\end{itemize}

Poiché l'errore cresce già sul set di addestramento, si deduce che questo fenomeno \textbf{non è causato da overfitting} (il quale si manifesterebbe con un errore di training basso e un errore di test alto). Si tratta invece di un problema di **degradazione** legato alla difficoltà di ottimizzazione: all'aumentare dei livelli, la superficie di costo diventa complessa e i tradizionali algoritmi basati su gradiente faticano a convergere verso soluzioni ottimali.

% [IMMAGINE DEL FILE DI SLIDE CON NOME "53-78"] Grafici degli errori di addestramento (sinistra) e di test (destra) che evidenziano le scarse prestazioni della rete a 56 strati rispetto a quella a 20 strati — Slide 67 / Pagina 15 del PDF
% [IMMAGINE DEL FILE DI SLIDE CON NOME "53-78"] Spiegazione del problema di degradazione e precisazione sul fatto che non sia legato a overfitting — Slide 68 / Pagina 16 del PDF

\subsection{L'Apprendimento Residuale (Residual Learning)}

% [IMMAGINE DEL FILE DI SLIDE CON NOME "53-78"] Formulazione dell'ipotesi sul problema di ottimizzazione nelle reti profonde — Slide 69 / Pagina 17 del PDF

Dal punto di vista logico, una rete più profonda dovrebbe essere in grado di ottenere prestazioni \emph{almeno pari} a quelle di una rete più superficiale. Una soluzione ideale "per costruzione" consisterebbe nel copiare i livelli appresi dalla rete più corta e configurare i restanti strati aggiuntivi affinché eseguano una semplice **mappatura identità** (\emph{identity mapping}), lasciando inalterato il segnale in transito.

% [IMMAGINE DEL FILE DI SLIDE CON NOME "53-78"] Razionale matematico sul comportamento atteso della rete profonda e introduzione teorica alla mappatura identità — Slide 70 / Pagina 18 del PDF

Nelle reti convenzionali, tuttavia, per l'ottimizzatore è estremamente complesso addestrare molteplici strati convoluzionali non lineari affinché riproducano l'identità. ResNet aggira questo ostacolo modificando l'obiettivo di apprendimento dello strato: invece di approssimare direttamente la mappatura ottimale cercata $H(x)$, si addestra la rete a stimare la **funzione residua**:
\[ F(x) = H(x) - x \]
La mappatura originaria viene quindi riscritta introducendo una somma:
\[ \boxed{H(x) = F(x) + x} \]

% [IMMAGINE DEL FILE DI SLIDE CON NOME "53-78"] Schema a confronto tra un blocco "plain" e un blocco residuale dotato di connessione di shortcut — Slide 71 / Pagina 19 del PDF

Questa formulazione matematica si traduce nel **blocco residuale** (\emph{Residual block}):
\begin{itemize}
    \item L'input $x$ segue una connessione diretta (chiamata \textbf{shortcut} o ramo identità) che scavalca i livelli convoluzionali;
    \item In parallelo, $x$ attraversa i livelli convoluzionali del blocco, i quali calcolano la sola componente residuale $F(x)$;
    \item Prima dell'attivazione ReLU finale, si esegue una somma elemento per elemento tra l'input originario $x$ e il residuo stimato $F(x)$:
    \[ \text{Output} = \text{ReLU}(F(x) + x) \]
\end{itemize}
Se la mappatura ottimale è proprio la funzione identità, l'ottimizzatore dovrà semplicemente azzerare i pesi delle convoluzioni per ottenere $F(x) = 0$, un compito decisamente più agevole rispetto all'addestramento di uno strato plain.

% [IMMAGINE DEL FILE DI SLIDE CON NOME "53-78"] Dettaglio matematico della riformulazione del target di apprendimento tramite residuo H(x) = F(x) + x — Slide 72 / Pagina 20 del PDF

\subsection{Architettura ResNet Completa}

L'intera architettura ResNet si sviluppa collegando in cascata i blocchi residuali descritti, seguendo precise regole di progettazione:

\begin{itemize}
    \item \textbf{Struttura Del Blocco Base}: Ogni blocco residuale standard contiene al suo interno \textbf{due livelli convoluzionali $3\times3$}.

    % [IMMAGINE DEL FILE DI SLIDE CON NOME "53-78"] Rappresentazione della struttura a blocchi residuali con due livelli di convoluzione 3x3 — Slide 73 / Pagina 21 del PDF

    \item \textbf{Gestione Dei Filtri E Downsampling}: Periodicamente, il numero di filtri delle convoluzioni viene raddoppiato. In corrispondenza di questo incremento, si applica un downsampling spaziale dimezzando altezza e larghezza delle mappe delle caratteristiche (\hl{$/2$} in ogni dimensione spaziale). Questo viene ottenuto impostando uno \textbf{stride pari a 2} nella prima convoluzione del blocco interessato.

    % [IMMAGINE DEL FILE DI SLIDE CON NOME "53-78"] Schema dei blocchi residuali che evidenzia il raddoppio dei filtri e il downsampling spaziale con stride 2 — Slide 74 / Pagina 22 del PDF

    \item \textbf{Livello Convoluzionale Iniziale}: In testa alla rete è presente uno strato convoluzionale preliminare per l'elaborazione dell'input ad alta risoluzione, costituito da una convoluzione con kernel $7\times7$ e $64$ filtri, con stride 2.

    % [IMMAGINE DEL FILE DI SLIDE CON NOME "53-78"] Schema dell'architettura di ResNet con evidenziato lo strato convoluzionale iniziale 7x7 — Slide 75 / Pagina 23 del PDF

    \item \textbf{Struttura Terminale}: Analogamente a GoogLeNet, al termine della rete non vi sono strati Fully Connected densi. Subito dopo l'ultimo blocco convoluzionale viene inserito un livello di \emph{Global Average Pooling}, seguito esclusivamente da un unico livello lineare \textbf{FC 1000} per la proiezione sulle classi di output.

    % [IMMAGINE DEL FILE DI SLIDE CON NOME "53-78"] Schema dell'architettura di ResNet con evidenziato il Global Average Pooling e il livello FC 1000 terminale — Slide 76 / Pagina 24 del PDF

    \item \textbf{Profondità Disponibili}: Utilizzando questo schema costruttivo, vengono comunemente realizzate reti con profondità complessive di \textbf{34, 50, 101 o 152 livelli}.

    % [IMMAGINE DEL FILE DI SLIDE CON NOME "53-78"] Diagramma strutturale complessivo delle diverse varianti di profondità di ResNet — Slide 77 / Pagina 25 del PDF
\end{itemize}

\subsection{Variante Bottleneck Di ResNet Per Reti Profonde (ResNet-50+)}

Quando si progettano reti estremamente profonde, quali ResNet-50, ResNet-101 o ResNet-152, il blocco residuale standard formato da due strati $3\times3$ risulta computazionalmente gravoso. Si adotta allora un blocco residuale di tipo \textbf{"bottleneck"} (collo di bottiglia) per contenere il costo in termini di operazioni e memoria.

Assumendo un input di dimensione $28\times28\times256$, il blocco bottleneck si articola in tre passaggi convoluzionali:
\begin{enumerate}
    \item Una prima convoluzione \textbf{$1\times1$ con $64$ filtri}: riduce la profondità dell'input da $256$ a $64$ canali, comprimendo l'informazione;
    \item Una convoluzione intermedia \textbf{$3\times3$ con $64$ filtri}: esegue l'elaborazione spaziale dei feature map su un volume a profondità ridotta;
    \item Una terza convoluzione \textbf{$1\times1$ con $256$ filtri}: ripristina la profondità originale del canale portandola nuovamente a $256$.
\end{enumerate}
L'output di quest'ultima viene sommato alla connessione identità iniziale $x$ di dimensione $28\times28\times256$. Questa configurazione garantisce un drastico abbattimento del numero di parametri e del carico computazionale complessivo del modello.

% [IMMAGINE DEL FILE DI SLIDE CON NOME "53-78"] Schema del blocco residuale bottleneck utilizzato per le architetture ResNet molto profonde — Slide 26 / Pagina 26 del PDF

\section{Il Caso Di Studio Di ResNet E Il Blocco Bottleneck}

Nel contesto dello sviluppo di reti neurali convoluzionali profonde, l'architettura \textbf{ResNet} (proposta da He et al. nel 2015) rappresenta un punto di svolta fondamentale per superare le difficoltà legate alla svanizione del gradiente durante l'addestramento. Per reti estremamente profonde, tipicamente a partire da $50$ strati o più (ResNet-50+), gli autori introducono una variante strutturale del classico blocco residuale denominata blocco \textbf{\emph{bottleneck}} (collo di bottiglia), con l'obiettivo di ottimizzare e migliorare l'efficienza computazionale. Tale approccio condivide una filosofia simile a quella adottata nei moduli Inception dell'architettura GoogLeNet.

Nel blocco \emph{bottleneck}, la riduzione e il successivo ripristino della dimensionalità dei canali (profondità del tensore) consentono di eseguire la convoluzione spaziale più onerosa su uno spazio a dimensionalità ridotta. Consideriamo, a titolo di esempio, un tensore di input con dimensioni spaziali e numero di canali pari a $28 \times 28 \times 256$. L'elaborazione all'interno del modulo avviene attraverso tre passaggi sequenziali:

\begin{enumerate}
    \item \textbf{Compressione della dimensionalità (Livello $1 \times 1$ Conv)}: Viene applicato un livello convoluzionale con filtri di dimensione $1 \times 1$ e $64$ filtri complessivi. Questo livello proietta il tensore di input dallo spazio a $256$ canali a uno spazio ridotto a soli $64$ canali, generando un output intermedio di dimensioni:
    \[
    28 \times 28 \times 64
    \]
    \item \textbf{Convoluzione spaziale (Livello $3 \times 3$ Conv)}: La convoluzione spaziale principale, caratterizzata da filtri $3 \times 3$, opera esclusivamente su questa mappa di caratteristiche ridotta a $64$ canali. Questo consente di abbattere drasticamente il costo computazionale rispetto a una convoluzione $3 \times 3$ eseguita direttamente sui $256$ canali originari.
    \item \textbf{Ripristino della dimensionalità (Livello $1 \times 1$ Conv)}: Un ulteriore livello convoluzionale $1 \times 1$ con $256$ filtri proietta nuovamente le mappe di caratteristiche nello spazio originale a $256$ canali, producendo un tensore di output pari a:
    \[
    28 \times 28 \times 256
    \]
\end{enumerate}

Al termine di questa sequenza di tre operazioni convoluzionali, il tensore risultante viene sommato all'input originario di $28 \times 28 \times 256$ mediante la connessione residuale (\emph{skip connection}), prima di applicare l'attivazione non lineare finale.

% [IMMAGINE DEL FILE DI SLIDE CON NOME "79-101"] Schema architetturale di un blocco di tipo bottleneck in ResNet con convoluzioni 1x1 e 3x3 — Slide 79 / Pagina 1 del PDF

\subsection{Dettagli Di Addestramento Di ResNet Nella Pratica}

L'addestramento di un'architettura complessa e profonda come ResNet richiede l'adozione di specifiche e rigorose metodologie di regolarizzazione, ottimizzazione e inizializzazione dei parametri. Le scelte metodologiche adottate dagli autori per garantire la convergenza includono:

\begin{itemize}
    \item \textbf{Batch Normalization}: Applicata sistematicamente dopo ogni livello convoluzionale (\emph{CONV layer}) per stabilizzare la distribuzione degli input ai livelli successivi e accelerare la convergenza.
    \item \textbf{Inizializzazione Xavier/2}: Nota anche come inizializzazione di He, specificamente progettata per reti che utilizzano funzioni di attivazione di tipo ReLU.
    \item \textbf{Algoritmo di Ottimizzazione}: Ottimizzazione basata su discesa stocastica del gradiente con momentum (\textbf{SGD + Momentum}) impostato a $0.9$.
    \item \textbf{Politica del Tasso di Apprendimento (\emph{Learning Rate})}: Valore iniziale fissato a $0.1$, con una strategia di decadimento che prevede la divisione del fattore per $10$ ogni volta che l'errore calcolato sul set di validazione raggiunge un plateau.
    \item \textbf{Dimensione del Mini-batch}: Impostata a un valore pari a $256$.
    \item \textbf{Decadimento dei Pesi (\emph{Weight Decay})}: Regolarizzazione $L_2$ applicata ai pesi con un coefficiente pari a $10^{-5}$ ($1\text{e-}5$).
    \item \textbf{Esclusione del Dropout}: Non viene utilizzato alcun meccanismo di \emph{dropout} all'interno della rete, in quanto la normalizzazione dei batch (\emph{Batch Normalization}) svolge già una sufficiente azione regolarizzatrice.
\end{itemize}

\subsection{Risultati Sperimentali E Prestazioni}

Grazie all'introduzione delle connessioni residuali, ResNet è in grado di addestrare modelli estremamente profondi senza incorrere nel fenomeno della degradazione delle prestazioni dell'addestramento, che storicamente affliggeva i modelli tradizionali all'aumentare del numero di strati. I risultati sperimentali evidenziano le seguenti pietre miliari:

\begin{itemize}
    \item Capacità di addestrare con successo reti fino a $152$ strati sul dataset ImageNet e fino a $1202$ strati sul dataset CIFAR.
    \item Le reti più profonde riescono a raggiungere un errore di addestramento (\emph{training error}) sistematicamente inferiore, coerentemente con le aspettative teoriche sul potenziale di rappresentazione dei modelli ad alta capacità.
    \item Dominio assoluto nelle competizioni di computer vision del 2015, dove il gruppo di ricerca MSRA ha conquistato il primo posto in tutti e cinque i principali tracciati previsti dalle sfide ILSVRC e COCO 2015.
\end{itemize}

% [IMMAGINE DEL FILE DI SLIDE CON NOME "79-101"] Tabella dei risultati del team MSRA alle competizioni ILSVRC e COCO 2015 — Slide 81 / Pagina 3 del PDF

In particolare, il modello vincitore della classificazione ILSVRC 2015 ha registrato un tasso di errore \textbf{top-5 pari al $3.6\%$}, un valore inferiore al livello di accuratezza stimato per le prestazioni umane, quantificato al $5.1\%$ nello studio di Russakovsky del 2014.

% [IMMAGINE DEL FILE DI SLIDE CON NOME "79-101"] Annuncio del superamento della performance umana alla competizione ILSVRC 2015 — Slide 82 / Pagina 4 del PDF

\newpage

\section{Analisi Della Complessità E Confronto Tra Architetture}

L'evoluzione storica delle reti neurali convoluzionali nell'ambito della competizione ImageNet (\emph{ImageNet Large Scale Visual Recognition Challenge} - ILSVRC) mostra una chiara e progressiva riduzione del tasso di errore top-5, strettamente legata all'aumento della profondità dei modelli proposti di anno in anno:

\begin{itemize}
    \item \textbf{ILSVRC'10}: Modelli superficiali (\emph{shallow}), tasso di errore pari al $28.2\%$.
    \item \textbf{ILSVRC'11}: Modelli superficiali, tasso di errore pari al $25.8\%$.
    \item \textbf{ILSVRC'12 (AlexNet)}: Primo modello profondo con $8$ strati, riduzione drastica dell'errore al $16.4\%$.
    \item \textbf{ILSVRC'13}: Architettura a $8$ strati perfezionata, tasso di errore dell'$11.7\%$.
    \item \textbf{ILSVRC'14 (VGG)}: Incremento della profondità a $19$ strati, tasso di errore del $7.3\%$.
    \item \textbf{ILSVRC'14 (GoogLeNet)}: Introduzione dei blocchi Inception su $22$ strati, tasso di errore del $6.7\%$.
    \item \textbf{ILSVRC'15 (ResNet)}: Sfruttamento delle connessioni skip per raggiungere $152$ strati, stabilendo il record con un errore del $3.57\%$.
\end{itemize}

% [IMMAGINE DEL FILE DI SLIDE CON NOME "79-101"] Grafico dell'evoluzione del tasso d'errore top-5 dei vincitori di ILSVRC dal 2010 al 2015 in funzione del numero di strati — Slide 83 / Pagina 5 del PDF

Un'analisi sistematica condotta da Canziani et al. (2017) nel lavoro \emph{"An Analysis of Deep Neural Network Models for Practical Applications"} mette in relazione l'accuratezza top-1 registrata su ImageNet con la complessità computazionale (espressa in miliardi di operazioni, G-Ops) e l'occupazione di memoria (rappresentata graficamente dalla dimensione delle sfere corrispondenti ai parametri del modello).

% [IMMAGINE DEL FILE DI SLIDE CON NOME "79-101"] Grafici di confronto tra accuratezza Top-1, costo computazionale in G-Ops e dimensione della memoria per diverse CNN — Slide 84 / Pagina 6 del PDF

Dall'analisi integrata di questi grafici emergono considerazioni chiave sulle proprietà delle diverse famiglie di reti:

\begin{itemize}
    \item \textbf{Inception-v4}: Questa architettura unisce i principi dei moduli Inception con le connessioni residuali di ResNet (\textbf{ResNet + Inception!}), posizionandosi nella fascia più alta in termini di accuratezza top-1, pur mantenendo un costo computazionale e un impatto di memoria bilanciati.

    % [IMMAGINE DEL FILE DI SLIDE CON NOME "79-101"] Evidenziazione di Inception-v4 come combinazione di ResNet e Inception con alta accuratezza — Slide 85 / Pagina 7 del PDF

    \item \textbf{VGG-16 / VGG-19}: Rappresentano i modelli con il \textbf{maggior consumo di memoria} e il \textbf{maggior numero di operazioni} complessive (con un numero di parametri che supera i $130$ milioni). Nonostante l'elevata impronta computazionale, l'accuratezza top-1 risulta inferiore rispetto a soluzioni più moderne ed efficienti.

    % [IMMAGINE DEL FILE DI SLIDE CON NOME "79-101"] Evidenziazione del consumo di memoria e operazioni delle reti VGG — Slide 86 / Pagina 8 del PDF

    \item \textbf{GoogLeNet}: Si conferma come l'architettura \textbf{più efficiente} dal punto di vista del rapporto tra prestazioni e risorse computazionali, mantenendo un numero di parametri estremamente ridotto.

    % [IMMAGINE DEL FILE DI SLIDE CON NOME "79-101"] Evidenziazione dell'efficienza di GoogLeNet nel grafico accuratezza vs operazioni — Slide 87 / Pagina 9 del PDF

    \item \textbf{AlexNet}: Richiede un volume di calcoli ridotto (\emph{smaller compute}) ma rimane inefficiente in termini di memoria rispetto alle prestazioni fornite, mostrando il livello di accuratezza più basso tra i modelli analizzati.

    % [IMMAGINE DEL FILE DI SLIDE CON NOME "79-101"] Evidenziazione delle limitazioni di AlexNet in accuratezza e memoria — Slide 88 / Pagina 10 del PDF

    \item \textbf{ResNet}: Offre un livello di efficienza moderato che varia sensibilmente in base alla profondità della configurazione scelta (da ResNet-18 a ResNet-152), garantendo al contempo le massime prestazioni in termini di accuratezza.

    % [IMMAGINE DEL FILE DI SLIDE CON NOME "79-101"] Posizionamento della famiglia ResNet come standard per accuratezza ed efficienza scalabile — Slide 89 / Pagina 11 del PDF
\end{itemize}

\subsection{Tempi Di Esecuzione E Consumo Energetico}

Un ulteriore fattore di confronto cruciale per le applicazioni in tempo reale, specialmente in contesti di robotica e sistemi embedded, è rappresentato dal tempo richiesto per l'esecuzione del passaggio in avanti (\emph{forward pass}) e dal corrispondente consumo di potenza elettrica dell'hardware.

% [IMMAGINE DEL FILE DI SLIDE CON NOME "79-101"] Grafici del tempo di esecuzione del forward pass (ms per immagine) e del consumo energetico (W) in funzione della dimensione del batch — Slide 90 / Pagina 12 del PDF

Dall'analisi dei dati si rileva che il tempo di elaborazione per singola immagine (espresso in millisecondi) cresce significativamente per modelli complessi come VGG-16 e VGG-19, mentre si mantiene su livelli ottimali per modelli veloci come ENet, BN-NIN e AlexNet. Le varianti di ResNet si collocano in una fascia intermedia, dove la scalabilità del modello permette di bilanciare le prestazioni temporali con l'accuratezza richiesta dall'applicazione specifica. Il consumo energetico netto (espresso in Watt) fluttua anch'esso in relazione all'ottimizzazione del modello e alla dimensione dei batch di dati elaborati simultaneamente (\emph{batch size}).

\newpage

\section{Altre Architetture Da Conoscere E Sviluppi Successivi}

Oltre alle pietre miliari della classificazione d'immagini, vi sono approcci teorici e architetture alternative che hanno influenzato lo sviluppo dei moderni modelli convoluzionali.

% [IMMAGINE DEL FILE DI SLIDE CON NOME "79-101"] Slide di transizione verso la rassegna di altre architetture CNN storiche e moderne — Slide 91 / Pagina 13 del PDF

\subsection{Network In Network (NiN)}

Proposta da Lin et al. nel 2014, l'architettura \textbf{Network in Network} (NiN) introduce concetti radicalmente innovativi che hanno anticipato la progettazione delle moderne reti deep learning:

\begin{itemize}
    \item \textbf{Livello mlpconv}: Sostituisce il classico livello convoluzionale lineare con una struttura interna detta "micro-rete" (\emph{micronetwork}). Questa ha lo scopo di calcolare caratteristiche e rappresentazioni locali molto più astratte rispetto a un filtro lineare standard.
    \item \textbf{Struttura della Micro-rete}: La micro-rete è realizzata mediante un perceptron multistrato (MLP), che equivale matematicamente all'applicazione di molteplici livelli convoluzionali con filtri di dimensione $1 \times 1$ inseriti all'interno dello stesso blocco.
    \item \textbf{Precursore dei Blocchi Moderni}: Rappresenta il precursore teorico e l'ispirazione filosofica principale sia per lo sviluppo dei moduli Inception in GoogLeNet, sia per l'implementazione dei livelli \emph{bottleneck} usati in ResNet per la manipolazione della profondità dei canali.
\end{itemize}

% [IMMAGINE DEL FILE DI SLIDE CON NOME "79-101"] Confronto strutturale tra un livello convoluzionale lineare classico e il livello mlpconv di NiN basato su multilayer perceptron — Slide 92 / Pagina 14 del PDF

\newpage

\section{Evoluzioni E Ottimizzazioni Di ResNet}

Il successo di ResNet ha stimolato la nascita di numerose ricerche volte a ottimizzarne la struttura interna, migliorandone le prestazioni o l'efficienza computazionale.

\subsection{Mappature Di Identità In Deep Residual Networks}

He et al. (2016) hanno proposto una revisione formale del design del blocco residuale originale. Questa variante, nota come \textbf{pre-attivazione}, si focalizza sulla creazione di un percorso privo di ostacoli per il flusso dell'informazione:

\begin{itemize}
    \item Sposta le operazioni di Batch Normalization (BN) e l'attivazione non lineare ReLU \textbf{prima} del livello convoluzionale (da cui il termine pre-attivazione).
    \item Questo accorgimento assicura che il segnale che attraversa la connessione skip (identità) possa propagarsi liberamente senza subire alcuna alterazione non lineare diretta, migliorando la stabilità matematica del gradiente durante la backpropagation.
    \item Sperimentalmente, questa configurazione produce prestazioni superiori su reti molto profonde.
\end{itemize}

% [IMMAGINE DEL FILE DI SLIDE CON NOME "79-101"] Diagramma di blocco a confronto tra la formulazione originale di ResNet e la variante con pre-attivazione (Identity Mapping) — Slide 93 / Pagina 15 del PDF

\subsection{Wide Residual Networks}

Zagoruyko et al. (2016) mettono in discussione l'ipotesi che la sola profondità sia l'ingrediente fondamentale per il successo delle reti residuali:

\begin{itemize}
    \item Gli autori propongono l'uso di \textbf{blocchi residuali larghi} (\emph{Wide Residual Networks}), incrementando il numero di filtri in ciascun livello di un fattore $k$ (ottenendo $F \times k$ filtri totali, rispetto ai soli $F$ filtri del modello standard).
    \item Un modello Wide ResNet profondo solo $50$ strati riesce a superare in accuratezza una ResNet classica caratterizzata da $152$ strati.
    \item L'incremento della larghezza del blocco, rispetto all'estensione della profondità della rete, si rivela computazionalmente più conveniente poiché le operazioni di calcolo risultano altamente parallelizzabili sui moderni hardware di calcolo (GPU).
\end{itemize}

% [IMMAGINE DEL FILE DI SLIDE CON NOME "79-101"] Schema a confronto tra un blocco residuale di base e un blocco residuale largo con fattore di moltiplicazione dei filtri — Slide 94 / Pagina 16 del PDF

\subsection{ResNeXt (Aggregated Residual Transformations)}

Proposta da Xie et al. nel 2016, l'architettura \textbf{ResNeXt} introduce una nuova dimensione nello spazio di progettazione delle reti neurali, definita come \textbf{cardinalità}:

\begin{itemize}
    \item Invece di limitarsi ad aumentare la profondità o la larghezza, il blocco residuale viene suddiviso internamente in molteplici percorsi di elaborazione paralleli ed omogenei.
    \item Il numero totale di questi rami paralleli rappresenta la cardinalità del blocco. Questo principio riprende la filosofia multi-percorso dei blocchi Inception, ma ne semplifica l'implementazione rendendo i canali di calcolo strutturalmente identici tra loro.
    \item Ad esempio, un blocco di tipo bottleneck standard con $256$ canali di input viene convertito in un modulo ad alta cardinalità con ben $32$ percorsi paralleli indipendenti, in cui ciascun percorso elabora localmente un numero ridotto di canali (ad esempio, con filtri interni impostati a dimensione ridotta, come $1 \times 1, 4 \to 3 \times 3, 4 \to 1 \times 1, 256$), accumulando poi i contributi in un'unica somma finale prima del skip connection.
\end{itemize}

% [IMMAGINE DEL FILE DI SLIDE CON NOME "79-101"] Rappresentazione grafica di un blocco ResNeXt caratterizzato da 32 percorsi di elaborazione paralleli rispetto a un blocco ResNet convenzionale — Slide 95 / Pagina 17 del PDF

\subsection{Deep Networks With Stochastic Depth}

Sviluppata da Huang et al. nel 2016, la tecnica della \textbf{profondità stocastica} (\emph{Stochastic Depth}) propone una soluzione elegante sia per mitigare la svanizione del gradiente sia per velocizzare le tempistiche di addestramento:

\begin{itemize}
    \item Durante la fase di addestramento, un sottoinsieme di blocchi residuali viene disattivato in modo stocastico a ogni iterazione (passaggio in avanti e all'indietro).
    \item I blocchi esclusi vengono completamente bypassati sostituendo il loro percorso di calcolo con una semplice funzione identità, riducendo di fatto la profondità della rete "vista" dall'algoritmo in quel passo d'addestramento.
    \item In fase di inferenza (\emph{test time}), la rete viene utilizzata nella sua interezza, sfruttando tutti i blocchi precedentemente addestrati per garantire la massima capacità rappresentativa del modello.
\end{itemize}

% [IMMAGINE DEL FILE DI SLIDE CON NOME "79-101"] Schema del processo di esclusione dinamica e stocastica dei blocchi residuali durante l'addestramento — Slide 96 / Pagina 18 del PDF

\newpage

\section{Oltre Le Architetture Residuali: FractalNet E DenseNet}

La ricerca scientifica ha esplorato strutture topologiche alternative a ResNet, dimostrando che l'addestramento di reti estremamente profonde può essere affrontato anche senza l'ausilio di esplicite connessioni residuali additive.

\subsection{FractalNet}

Proposta da Larsson et al. nel 2017, l'architettura \textbf{FractalNet} introduce un approccio basato sulla geometria frattale per l'elaborazione dell'informazione:

\begin{itemize}
    \item L'assunto fondamentale del lavoro è che la capacità di addestrare reti ultra-profonde derivi principalmente dalla transizione progressiva ed efficace tra percorsi superficiali e percorsi profondi, rendendo l'uso esplicito delle connessioni residuali non strettamente necessario.
    \item La rete adotta un'architettura ricorsiva (frattale) che genera in parallelo sia percorsi brevi (\emph{shallow}) che percorsi molto lunghi (\emph{deep}) verso i livelli di output.
    \item Per regolarizzare il modello durante l'addestramento, viene introdotta una tecnica che prevede l'esclusione casuale di interi sotto-percorsi di calcolo (\emph{dropping out sub-paths}).
    \item Al momento dell'inferenza, la rete viene valutata integrando l'intera struttura frattale sviluppata.
\end{itemize}

% [IMMAGINE DEL FILE DI SLIDE CON NOME "79-101"] Schema della regola di espansione frattale e architettura a blocchi paralleli di FractalNet — Slide 97 / Pagina 19 del PDF

\subsection{DenseNet (Densely Connected Convolutional Networks)}

Huang et al. (2017) hanno introdotto una diversa tipologia di connessione multilivello con l'architettura \textbf{DenseNet}:

\begin{itemize}
    \item All'interno di ciascun **blocco denso** (\emph{dense block}), ogni singolo livello riceve come input le mappe di caratteristiche prodotte da tutti i livelli precedenti e, a sua volta, trasmette le proprie mappe di caratteristiche in input a tutti i livelli successivi in modo strettamente feedforward.
    \item Diversamente da ResNet, in cui le mappe vengono sommate, DenseNet effettua una \textbf{concatenazione} delle caratteristiche lungo la dimensione dei canali.
    \item Questo fitto schema di connessioni presenta tre vantaggi principali: attenua fortemente il problema della svanizione del gradiente, rafforza la propagazione delle caratteristiche lungo la rete e incentiva significativamente il riutilizzo delle feature (\emph{feature reuse}), riducendo la ridondanza complessiva del modello.
\end{itemize}

% [IMMAGINE DEL FILE DI SLIDE CON NOME "79-101"] Schema di un blocco denso basato su concatenazione progressiva e architettura generale di DenseNet con blocchi di transizione — Slide 20 / Pagina 20 del PDF

\newpage

\section{Architetture Efficienti: SqueezeNet}

In molti contesti reali, come sui robot mobili o sui dispositivi della Internet of Things (IoT), le limitazioni di memoria e di capacità di calcolo impongono l'uso di modelli estremamente compatti.

L'architettura \textbf{SqueezeNet} (Iandola et al., 2017) è stata progettata specificamente per raggiungere un livello di accuratezza paragonabile ad AlexNet su ImageNet, riducendo di oltre $50$ volte il numero di parametri complessivi e mantenendo la dimensione totale del modello al di sotto di $0.5$ MB.

Questo risultato viene conseguito mediante l'adozione strutturata di moduli denominati \textbf{Fire modules}. Ciascun modulo Fire è composto da due fasi distinte:

\begin{enumerate}
    \item \textbf{Fase di Squeeze (Compressione)}: Un livello convoluzionale che impiega esclusivamente filtri di dimensione $1 \times 1$. Questo livello ha lo scopo di comprimere il numero di canali di input, riducendo drasticamente il volume dei dati inviati al livello successivo.
    \item \textbf{Fase di Expand (Espansione)}: Il tensore compresso viene inviato a un livello di espansione che elabora i dati in parallelo attraverso una combinazione di filtri $1 \times 1$ e filtri spaziali $3 \times 3$, le cui mappe di caratteristiche vengono infine concatenate all'uscita del modulo.
\end{enumerate}

Grazie a questa scomposizione e a successive tecniche di compressione dei parametri, SqueezeNet permette di ottenere un file di modello fino a $510$ volte più piccolo rispetto ad AlexNet, salvaguardando l'efficacia predittiva.

% [IMMAGINE DEL FILE DI SLIDE CON NOME "79-101"] Struttura interna di un modulo Fire in SqueezeNet composto da squeeze layer e expand layer — Slide 21 / Pagina 21 del PDF

\newpage

\section{Sintesi Delle Architetture CNN}

\subsection{Casi Di Studio E Modelli Chiave}

Il panorama delle reti neurali convoluzionali analizzato si articola in una serie di tappe fondamentali che ogni progettista di sistemi di percezione deve conoscere:

\begin{itemize}
    \item \textbf{Modelli Fondamentali (Casi di Studio)}: AlexNet, VGG, GoogLeNet e ResNet.
    \item \textbf{Architetture Specialistiche ed Evoluzioni}: Network in Network (NiN), Wide ResNet, ResNeXt, Stochastic Depth, DenseNet, FractalNet e SqueezeNet.
\end{itemize}

% [IMMAGINE DEL FILE DI SLIDE CON NOME "79-101"] Tabella riassuntiva dei principali modelli CNN di riferimento e delle varianti avanzate — Slide 22 / Pagina 22 del PDF

\subsection{Tendenze E Linee Guida Per Lo Studio}

Dall'analisi complessiva dello stato dell'arte emergono alcune importanti conclusioni metodologiche:

\begin{itemize}
    \item I modelli **VGG, GoogLeNet e ResNet** sono ampiamente diffusi, ampiamente testati in contesti industriali e accademici, e sono immediatamente accessibili all'interno delle principali librerie di modelli pre-addestrati (\emph{model zoos}).
    \item **ResNet** si conferma attualmente come la scelta di riferimento ottimale (\emph{best default}) per l'avvio di qualsiasi progetto o applicazione pratica nell'ambito della computer vision.
    \item Si registra una costante tendenza di ricerca verso lo sviluppo di reti estremamente profonde.
    \item La ricerca contemporanea si concentra intensamente sulla progettazione strutturale dei blocchi convoluzionali, sull'implementazione ottimale delle connessioni a salto (\emph{skip connections}) e sul miglioramento del flusso del gradiente.
    \item I filoni di ricerca più recenti analizzano con rigore la reale necessità della profondità rispetto alla larghezza dei canali di calcolo, unitamente all'impatto e alla formulazione matematica delle connessioni residuali.
\end{itemize}

% AA Slides 11.2 Proprietà

\section{Introduzione Alle Reti Neurali Convoluzionali}

Le \textbf{Reti Neurali Convoluzionali} (\emph{Convolutional Neural Networks}, o \textbf{CNN}) costituiscono una classe specializzata di architetture neurali progettate specificamente per elaborare dati strutturati secondo una topologia a griglia, come nel caso tipico delle immagini bidimensionali.

A livello macroscopico, un'architettura convoluzionale riceve un volume di input e lo elabora attraverso una sequenza di operazioni locali per estrarre mappe di caratteristiche, ridurne progressivamente la dimensionalità spaziale e infine mappare tali informazioni su uno o più strati completamente connessi per produrre l'output finale.

% [IMMAGINE DEL FILE DI SLIDE CON NOME "1-16"] Rappresentazione ad alto livello di una CNN (convoluzioni, subsampling e strati fully connected) — Slide 1 / Pagina 1 del PDF

Nel diagramma classico ideato da LeCun et al. (1998), l'elaborazione segue un flusso ben definito:
\begin{enumerate}
    \item \textbf{Input}: Il segnale di ingresso (ad esempio, l'immagine di un carattere come la lettera ``K'').
    \item \textbf{Convoluzioni} (\emph{Convolutions}): L'applicazione di filtri locali per generare le mappe delle immagini (\textbf{Image Maps} o \emph{Feature Maps}).
    \item \textbf{Sottocampionamento} (\textbf{Subsampling} o \emph{Pooling}): Una fase di riduzione spaziale per rendere la rappresentazione più robusta alle variazioni geometriche locali.
    \item \textbf{Strati Completamente Connessi} (\textbf{Fully Connected}): I livelli terminali che aggregano le caratteristiche estratte per compiere la classificazione o regressione finale, producendo l'output del modello.
\end{enumerate}

\subsection{Struttura E Terminologia Dei Livelli Convoluzionali}

Nella letteratura scientifica e nei framework di sviluppo si riscontrano due diverse convenzioni terminologiche per descrivere la struttura interna di un livello in una CNN. È fondamentale comprendere entrambe per evitare ambiguità concettuali.

% [IMMAGINE DEL FILE DI SLIDE CON NOME "1-16"] Confronto tra la terminologia complessa e semplice di un livello convoluzionale — Slide 2 / Pagina 2 del PDF

Il pattern architetturale più comune in una CNN è composto dalla sequenza di:
\begin{itemize}
    \item Un livello convoluzionale (\textbf{Convolutional layer});
    \item Un livello di non linearità (\textbf{Non linearity layer});
    \item Un livello di \emph{pooling} opzionale (\textbf{Optional pooling layer});
    \item Un livello di normalizzazione opzionale (\textbf{Optional normalization layer}).
\end{itemize}

Questi stadi possono essere raggruppati secondo due diverse prospettive:

\subsubsection{Terminologia Complessa del Livello}
In questa formulazione (spesso utilizzata in testi teorici come il manuale di Goodfellow et al.), l'intero blocco viene considerato come un unico grande ``Livello Convoluzionale'' suddiviso internamente in tre stadi distinti:
\begin{enumerate}
    \item \textbf{Stadio di Convoluzione} (\emph{Convolution stage}): Esegue una trasformazione affine lineare sull'input attraverso l'operazione di convoluzione.
    \item \textbf{Stadio del Rilevatore} (\emph{Detector stage}): Applica una funzione di attivazione non lineare (ad esempio, la funzione rettificata lineare, comunemente nota come \textbf{ReLU}).
    \item \textbf{Stadio di Pooling} (\emph{Pooling stage}): Esegue una riduzione spaziale del volume di attivazione prima di passare l'output al livello successivo.
\end{enumerate}

\subsubsection{Terminologia Semplice del Livello}
Nei moderni framework di programmazione (come PyTorch o TensorFlow), si preferisce una scomposizione più granulare in cui ogni singola operazione matematica è identificata come un livello a sé stante:
\begin{enumerate}
    \item \textbf{Livello di Convoluzione} (\emph{Convolution layer}): Responsabile unicamente della trasformazione affine.
    \item \textbf{Livello Rilevatore} (\emph{Detector layer}): Applica esclusivamente la non linearità (es. ReLU).
    \item \textbf{Livello di Pooling} (\emph{Pooling layer}): Svolge l'operazione di riduzione dimensionale spaziale.
\end{enumerate}

\newpage

\section{Il Livello Convoluzionale E La Formulazione Matematica}

Le reti convoluzionali si distinguono dalle reti completamente connesse poiché utilizzano l'operazione matematica di convoluzione in almeno uno dei loro livelli, in sostituzione della classica moltiplicazione di matrici.

\subsection{Definizione Matematica Di Convoluzione}

Nel dominio continuo del tempo e dello spazio, la convoluzione tra una funzione di input $x(t)$ e una funzione peso o nucleo $w(t)$ (definita \textbf{kernel} o filtro) è espressa dall'integrale:
\[
s(t) = \int x(a)w(t - a)\,da
\]
L'operazione viene denotata in forma compatta con l'operatore asterisco:
\[
s(t) = (x * w)(t)
\]
Nel contesto dell'elaborazione delle immagini digitali, operiamo in un dominio discreto in cui i segnali sono campionati su una griglia bidimensionale. Di conseguenza, l'integrale continuo si trasforma in una doppia sommatoria discreta. Data un'immagine bidimensionale $I$ e un kernel bidimensionale $K$, la convoluzione discreta è definita come:
\[
S(i, j) = (I * K)(i, j) = \sum_{m} \sum_{n} I(m, n) K(i - m, j - n)
\]
In questa formulazione, gli indici del kernel subiscono un'inversione di segno (riflessione), che garantisce importanti proprietà matematiche.

\subsection{Convoluzione E Cross-Correlazione Nei Framework Di Deep Learning}

L'operazione di convoluzione discreta gode della proprietà commutativa. Ciò significa che l'ordine degli operandi può essere scambiato senza variare il risultato finale:
\[
S(i, j) = (K * I)(i, j) = \sum_{m} \sum_{n} I(i - m, j - n) K(m, n)
\]
Tuttavia, molte librerie di machine learning implementano una variante leggermente diversa, nota come \textbf{cross-correlazione} (\emph{cross-correlation}), pur continuando a chiamarla formalmente ``convoluzione''. La formula della cross-correlazione non prevede il ribaltamento del kernel rispetto all'origine:
\[
S(i, j) = (I * K)(i, j) = \sum_{m} \sum_{n} I(i + m, j + n) K(m, n)
\]
Dal punto di vista pratico dell'apprendimento, questa differenza non costituisce un problema. Poiché i pesi del kernel $K$ sono parametri liberi appresi dall'algoritmo di ottimizzazione tramite retropropagazione (\emph{backpropagation}), l'algoritmo si adatterà automaticamente apprendendo una versione speculare del kernel ideale, compensando così l'assenza del ribaltamento spaziale previsto dalla convoluzione pura.

\subsection{Rappresentazione Visiva E Calcolo Passo-Passo}

Per comprendere concretamente l'applicazione della cross-correlazione bidimensionale (il tipo di convoluzione standard utilizzato nelle CNN), consideriamo un esempio con un input di dimensioni $3 \times 4$ e un kernel di dimensioni $2 \times 2$.

% [IMMAGINE DEL FILE DI SLIDE CON NOME "1-16"] Esempio di calcolo di una convoluzione 2D (cross-correlazione) — Slide 5 / Pagina 5 del PDF

Siano definiti la matrice di input $\bm{I}$ e la matrice del kernel $\bm{K}$:
\[
\bm{I} = \begin{pmatrix}
a & b & c & d \\
e & f & g & h \\
i & j & k & l
\end{pmatrix}, \quad
\bm{K} = \begin{pmatrix}
w & x \\
y & z
\end{pmatrix}
\]
Applicando il kernel scorso sull'immagine con passo unitario (\emph{stride} pari a 1) e senza alcuna imbottitura ai bordi (\emph{padding} pari a 0), la matrice di output risultante $\bm{S}$ avrà dimensioni $2 \times 3$. Ciascun elemento di $\bm{S}$ corrisponde al prodotto scalare elemento per elemento (\emph{Hadamard product}) tra il kernel e la sottomatrice locale dell'input coperta dal kernel in quella specifica posizione:

\begin{align*}
\bm{S}_{1,1} &= aw + bx + ey + fz \\
\bm{S}_{1,2} &= bw + cx + fy + gz \\
\bm{S}_{1,3} &= cw + dx + gy + hz \\
\bm{S}_{2,1} &= ew + fx + iy + jz \\
\bm{S}_{2,2} &= fw + gx + jy + kz \\
\bm{S}_{2,3} &= gw + hx + ky + lz
\end{align*}

Questo calcolo esplicita come la convoluzione combini le informazioni spazialmente contigue dell'input per generare ciascun valore della mappa di attivazione di output.

\newpage

\section{Proprietà Fondamentali E Motivazioni Architetturali}

Le reti convoluzionali possiedono tre proprietà fondamentali che le rendono estremamente efficaci nell'elaborazione delle immagini rispetto ai modelli densi tradizionali:
\begin{enumerate}
    \item L'\textbf{interazione sparsa} (\emph{sparse interaction} o \emph{sparse connectivity});
    \item La \textbf{condivisione dei parametri} (\emph{parameter sharing});
    \item La \textbf{rappresentazione equivariante} (\emph{equivariant representation}).
\end{enumerate}

\subsection{Interazione Sparsa E Connettività Locale}

Nei livelli lineari di una rete neurale tradizionale (\emph{fully connected}), l'interazione tra l'input e l'output è descritta da una moltiplicazione matrice-vettore. Ogni singola unità di output interagisce direttamente con ogni singola unità di input, richiedendo un parametro distinto per descrivere tale connessione.

Al contrario, i livelli di una rete convoluzionale presentano \textbf{interazioni sparse} (connettività sparsa o pesi sparsi). Questo comportamento si ottiene imponendo che la dimensione spaziale del kernel sia nettamente inferiore rispetto a quella dell'input.

% [IMMAGINE DEL FILE DI SLIDE CON NOME "1-16"] Connettività sparsa vista dal basso (influenza di un input sugli output) — Slide 7 / Pagina 7 del PDF

\begin{itemize}
    \item \textbf{Prospettiva dal basso} (\emph{view from below}): Nel livello convoluzionale (con kernel di dimensione spaziale pari a 3), l'unità di input $x_3$ influisce solo su un sottoinsieme locale di output, nello specifico $s_2$, $s_3$ e $s_4$. Nel livello denso classico, lo stesso input $x_3$ influisce direttamente su tutte le unità di output da $s_1$ a $s_5$.
\end{itemize}

% [IMMAGINE DEL FILE DI SLIDE CON NOME "1-16"] Connettività sparsa vista dall'alto (influenza degli input su un singolo output) — Slide 8 / Pagina 8 del PDF

\begin{itemize}
    \item \textbf{Prospettiva dall'alto} (\emph{view from above}): Se analizziamo l'unità di output $s_3$, notiamo che in una CNN essa è influenzata unicamente da una porzione localizzata di input, ovvero il set $\{x_2, x_3, x_4\}$. In una rete densa, l'unità $s_3$ riceve connessioni dirette da ogni singolo elemento di input $\{x_1, x_2, x_3, x_4, x_5\}$.
\end{itemize}

\subsection{Il Concetto Di Campo Ricettivo}

L'insieme di unità di ingresso che influenzano lo stato di una determinata unità $s_3$ posta a una generica profondità $k$ prende il nome di \textbf{campo ricettivo} (\emph{receptive field}) di quel neurone.

% [IMMAGINE DEL FILE DI SLIDE CON NOME "1-16"] Campo ricettivo diretto di un neurone — Slide 9 / Pagina 9 del PDF

Mentre nello strato convoluzionale adiacente il campo ricettivo di $s_3$ si limita strettamente alle unità $\{x_2, x_3, x_4\}$, la situazione evolve in modo significativo quando si considerano architetture composte da più strati in cascata.

% [IMMAGINE DEL FILE DI SLIDE CON NOME "1-16"] Espansione del campo ricettivo attraverso strati convoluzionali successivi — Slide 10 / Pagina 10 del PDF

Si consideri un'architettura a tre livelli: lo strato di input $x$, lo strato nascosto $h$ e lo strato di output $g$. Ipotizzando filtri di ampiezza spaziale pari a 3:
\begin{itemize}
    \item L'unità $g_3$ possiede un campo ricettivo immediato nello strato precedente pari a $\{h_2, h_3, h_4\}$.
    \item Ciascuna di queste unità nascoste, a sua volta, possiede un campo ricettivo nell'input originario:
    \begin{align*}
        h_2 &\leftarrow \{x_1, x_2, x_3\} \\
        h_3 &\leftarrow \{x_2, x_3, x_4\} \\
        h_4 &\leftarrow \{x_3, x_4, x_5\}
    \end{align*}
    \item Aggregando queste relazioni, si osserva che l'unità $g_3$ nello strato più profondo risente dell'influenza dell'intero volume di input $\{x_1, x_2, x_3, x_4, x_5\}$.
\end{itemize}

Questo fenomeno mostra come, nonostante l'uso di connessioni locali e sparse in ciascun livello singolo, lo strato finale di una rete convoluzionale profonda sia in grado di integrare informazioni globali provenienti da tutto l'input spaziale originario.

\subsection{Benefici Computazionali Dell'Interazione Sparsa}

L'introduzione della connettività sparsa garantisce vantaggi quantificabili sia in termini di memoria sia di tempo di calcolo:
\begin{itemize}
    \item \textbf{Efficienza statistica e di memoria}: Un'immagine digitale può contenere milioni di pixel. Tuttavia, le caratteristiche locali più significative (come i bordi o i contorni) possono essere rilevate in modo ottimale con kernel di piccole dimensioni (poche decine o centinaia di pixel). Ciò riduce drasticamente il numero di parametri da memorizzare.
    \item \textbf{Complessità computazionale}: Supponendo di avere $m$ unità di input e $n$ unità di output:
    \begin{itemize}
        \item In una rete densa classica basata sulla moltiplicazione di matrici, la complessità dell'operazione è descritta da:
        \[
        \mathcal{O}(n \times m)
        \]
        \item Utilizzando un livello convoluzionale con un kernel di dimensione spaziale limitata pari a $k$, la complessità si riduce a:
        \[
        \mathcal{O}(k \times n)
        \]
        \item Poiché nelle tipiche applicazioni reali la dimensione del kernel è molto piccola (comunemente $k = 3$) rispetto al numero totale di input ($k \ll m$), il guadagno computazionale risultante è estremamente elevato.
    \end{itemize}
\end{itemize}

\subsection{Condivisione Dei Parametri}

La \textbf{condivisione dei parametri} (\emph{parameter sharing}) consiste nell'impiegare lo stesso parametro per più funzioni di connessione all'interno dello stesso modello.

Nelle reti densamente connesse tradizionali, ogni elemento della matrice dei pesi viene moltiplicato per un elemento dell'input e successivamente non viene più riutilizzato per calcolare altre unità dello stesso livello. Nelle reti convoluzionali, invece, ogni singolo coefficiente del kernel viene applicato sistematicamente in ogni posizione spaziale dell'input.

% [IMMAGINE DEL FILE DI SLIDE CON NOME "1-16"] Confronto della condivisione dei parametri (pesi) tra CNN e rete completamente connessa — Slide 12 / Pagina 12 del PDF

Questo principio implica che, anziché dover apprendere un insieme di parametri indipendente per ciascuna coordinata spaziale dell'immagine, la rete apprende un solo set di pesi valido per l'intera estensione spaziale dell'input. Tale proprietà si rivela straordinariamente efficiente dal punto di vista della memoria e della prevenzione del sovraddestramento (\emph{overfitting}).

\subsection{Rappresentazione Equivariante Alla Traslazione}

La particolare struttura basata sulla condivisione dei parametri dota il livello convoluzionale della proprietà di \textbf{equivarianza alla traslazione}.

Una funzione $f$ si dice equivariante rispetto a una trasformazione $g$ se l'applicazione della trasformazione all'input produce sul risultato finale lo stesso effetto che si otterrebbe applicando la trasformazione direttamente all'output:
\[
f(g(x)) = g(f(x))
\]
Nel contesto delle immagini, se si applica una traslazione spaziale all'input e successivamente si esegue la convoluzione, l'output finale risulterà traslato esattamente della medesima quantità rispetto a quello che si sarebbe ottenuto eseguendo prima la convoluzione e poi la traslazione spaziale. Questa caratteristica assicura che se una specifica caratteristica saliente (ad esempio un bordo) si sposta all'interno dell'immagine, la sua rappresentazione nella mappa delle caratteristiche risultante si sposterà in modo corrispondente, mantenendo intatta l'informazione di coerenza spaziale.

Si segnala tuttavia un limite teorico rilevante: l'operazione di convoluzione \textbf{non è naturalmente equivariante} rispetto ad altre tipologie di trasformazione geometrica, quali le variazioni di scala o le rotazioni dell'immagine. Per gestire correttamente questo tipo di modifiche geometriche sull'input, è necessario ricorrere ad altri meccanismi dedicati.

\newpage

\section{Analisi Dimensionale E Calcolo Dei Parametri Di Output}

Durante la progettazione di un'architettura convoluzionale, è essenziale calcolare accuratamente la variazione delle dimensioni dei volumi di attivazione e il numero di parametri liberi introdotti da ogni livello.

% [IMMAGINE DEL FILE DI SLIDE CON NOME "1-16"] Connessione locale lungo le dimensioni spaziali e totale lungo la profondità del canale — Slide 14 / Pagina 14 del PDF

Come mostrato in figura, un aspetto chiave risiede nel fatto che ciascun neurone si connette localmente lungo le sole coordinate spaziali (altezza e larghezza), ma si estende per l'**intera profondità** del canale del volume di input.

Consideriamo due esempi numerici pratici basati su questa regola:
\begin{itemize}
    \item \textbf{Esempio 1}: Dato un volume di input avente dimensioni pari a $[32 \times 32 \times 3]$ (corrispondente a un'immagine a colori RGB del dataset CIFAR-10) e un kernel con campo ricettivo (dimensione spaziale del filtro) pari a $5 \times 5$. Ciascun neurone del livello convoluzionale possederà pesi distribuiti su una regione locale estesa per l'intera profondità dell'input, ovvero una porzione $[5 \times 5 \times 3]$. Il numero di pesi (\emph{weights}) per ciascun neurone sarà dunque pari a:
    \[
    5 \times 5 \times 3 = 75 \text{ pesi} \quad (+ 1 \text{ parametro di bias})
    \]
    L'estensione della connettività lungo l'asse della profondità deve essere obbligatoriamente pari a 3, in virtù del fatto che il volume di ingresso possiede 3 canali.

    \item \textbf{Esempio 2}: Si consideri un volume di input intermedio con dimensioni spaziali e di canale pari a $[16 \times 16 \times 20]$. Applicando un'operazione di convoluzione con un campo ricettivo avente dimensione spaziale pari a $3 \times 3$, ogni singolo neurone del livello convoluzionale stabilirà un numero di connessioni con il volume di input pari a:
    \[
    3 \times 3 \times 20 = 180 \text{ connessioni}
    \]
    Anche in questo secondo scenario, la connettività si mantiene confinata localmente nello spazio geometrico ($3 \times 3$), ma attraversa interamente lo spessore del canale d'ingresso ($20$).
\end{itemize}

\subsection{Iperparametri Di Controllo}

La dimensione finale spaziale del volume di output è regolata da tre iperparametri fondamentali:
\begin{enumerate}
    \item \textbf{Profondità} (\emph{depth}): Corrisponde al numero totale di filtri distinti impiegati nel livello convoluzionale. Ciascun filtro si specializza nell'apprendimento di una specifica caratteristica visiva dell'input. Di conseguenza, la profondità del volume di output coinciderà esattamente con il numero di filtri applicati.
    \item \textbf{Passo} (\emph{stride}): Rappresenta l'intervallo di scorrimento spaziale del filtro lungo l'input. Con uno stride pari a 1, i filtri avanzano di un solo pixel alla volta. Con uno stride pari a 2, i filtri saltano di due pixel alla volta, determinando un effetto di sottocampionamento (\emph{downsampling}) spaziale del volume di attivazione.
    \item \textbf{Zero-padding}: Consiste nell'aggiungere una cornice di zeri attorno al perimetro esterno del volume di ingresso, con l'obiettivo di preservare la dimensione spaziale e controllare gli effetti di bordo indotti dal filtro.
\end{enumerate}

\subsection{Formula Delle Dimensioni Di Output}

Definiamo le variabili matematiche associate a un livello convoluzionale:
\begin{itemize}
    \item $H_i \times W_i \times D_i$: Dimensioni spaziali e di profondità del volume di input;
    \item $K = (K_x, K_y)$: Dimensioni spaziali del filtro (kernel);
    \item $P = (P_x, P_y)$: Quantità di zero-padding applicata lungo le direzioni verticale e orizzontale;
    \item $S = (S_x, S_y)$: Passo di scorrimento (\emph{stride}) lungo gli assi verticale e orizzontale;
    \item $C$: Numero di filtri utilizzati (profondità dell'output).
\end{itemize}

La dimensione spaziale del volume di output è data dal prodotto cartesiano delle dimensioni calcolate indipendentemente per ciascun asse spaziale:

\[
\boxed{\left( \frac{H_i - K_x + 2P_x}{S_x} + 1 \right) \times \left( \frac{W_i - K_y + 2P_y}{S_y} + 1 \right) \times C}
\]

Affinché la rete sia valida e ben definita dal punto di vista dell'implementazione computazionale, i valori derivati dalle frazioni devono risultare rigorosamente interi. Non è infatti ammissibile ottenere dimensioni frazionarie per i pixel (ad esempio, un valore di output pari a 4.5).

\subsection{Calcolo Dei Parametri Totali Di Un Livello Convoluzionale}

Dall'analisi delle proprietà geometriche si deducono due regole fondamentali per il conteggio dei parametri del modello:
\begin{itemize}
    \item La profondità del volume di output non è in alcun modo influenzata dalla profondità del volume di input ($D_i$), bensì dipende unicamente dal numero di filtri configurati ($C$).
    \item Al contrario, il numero di parametri associati a ciascun singolo filtro dipende direttamente dalla profondità del volume di input ($D_i$).
\end{itemize}

Dato un livello convoluzionale caratterizzato da $C$ filtri aventi dimensioni spaziali $K = (K_x, K_y)$ applicati su un volume con profondità di ingresso $D_i$, la quantità complessiva di parametri addestrabili è regolata dalla formula:

\[
\boxed{\text{Parametri Totali} = K_x \times K_y \times D_i \times C + C \text{ bias}}
\]

In questa espressione, il termine $K_x \times K_y \times D_i \times C$ quantifica i pesi della matrice di convoluzione, mentre il termine $+C$ indica il singolo bias associato a ciascuno dei $C$ filtri di output.

\subsection{Gestione Dei Bordi Tramite Zero-Padding}

L'adozione dello zero-padding consente di gestire l'interazione del filtro in prossimità dei bordi geometrici del volume di input:
\begin{itemize}
    \item \textbf{Convoluzione Valida} (\emph{valid convolution}): Si verifica quando non viene applicato alcun tipo di padding ($P = 0$). In questa modalità, il kernel può posizionarsi solo dove l'intera matrice del filtro è contenuta completamente all'interno dei confini fisici del volume di input. Di conseguenza, a ogni passaggio del livello la dimensione spaziale della mappa di output subirà una contrazione geometrica legata alla dimensione del filtro utilizzato.
\end{itemize}

\section{Gestione Dei Bordi E Convoluzioni Speciali}

Nel capitolo precedente abbiamo visto come l'applicazione di un'operazione di convoluzione riduca progressivamente le dimensioni spaziali della mappa di attivazione di output. Per controllare accuratamente questa contrazione spaziale e preservare l'informazione geometrica, si ricorre a strategie specifiche di imbottitura (\emph{padding}) dei bordi e a configurazioni particolari del nucleo convoluzionale.

\subsection{Imbottitura Dei Bordi: Same Convolutions E Full Convolutions}

La formula generale per determinare le dimensioni spaziali dell'output lungo un asse, dati la dimensione dell'input $H_i$, la dimensione del filtro $K_x$, lo zero-padding $P_x$ e lo stride $S_x$, è:
\[
H_o = \frac{H_i - K_x + 2P_x}{S_x} + 1
\]
A seconda di come viene calibrato il parametro del padding $P_x$, si distinguono tre scenari operativi principali: la convoluzione \emph{valid} (priva di imbottitura, $P = 0$), la convoluzione \emph{same} e la convoluzione \emph{full}.

\subsubsection{Same Convolution}
Si definisce \textbf{same convolution} l'operazione in cui viene aggiunto un numero di zeri sufficiente a garantire che la dimensione spaziale dell'output sia esattamente identica a quella dell'input ($H_o = H_i$).

Ipotizzando uno stride unitario ($S_x = 1$), possiamo ricavare analiticamente il valore di $P_x$ necessario a soddisfare questa condizione uguagliando l'altezza di output a quella di input:
\begin{align*}
H_i - K_x + 2P_x + 1 &= H_i \\
2P_x &= K_x - 1 \\
\end{align*}
Risolvendo rispetto a $P_x$, otteniamo la relazione di progetto per l'imbottitura:
\[
\boxed{P_x = \frac{K_x - 1}{2}}
\]
Questa equazione evidenzia perché nella pratica si prediligano quasi esclusivamente filtri con dimensioni spaziali dispari (es. $3 \times 3$, $5 \times 5$): l'uso di un valore di $K_x$ dispari garantisce infatti che la quantità di padding $P_x$ da applicare simmetricamente su entrambi i lati del bordo sia un numero intero.

Tuttavia, un limite intrinseco della convoluzione \emph{same} risiede nel fatto che i pixel posizionati in prossimità del bordo dell'immagine di input partecipano a un numero inferiore di sotto-finestre di convoluzione rispetto ai pixel centrali, esercitando di conseguenza una minore influenza sulle caratteristiche estratte rispetto a questi ultimi.

\subsubsection{Full Convolution}
Per ovviare a questa asimmetria d'influenza, la \textbf{full convolution} prevede l'inserimento di una quantità di zero-padding tale da consentire a ciascun pixel dell'input di essere visitato esattamente lo stesso numero di volte $t$ in ogni direzione spaziale dal filtro di convoluzione. L'applicazione di questa configurazione determina un'espansione del volume di output, producendo un'immagine finale avente altezza (o larghezza) pari a:
\[
\boxed{H_o = H_i + t - 1}
\]

\subsection{Convoluzione 1x1}

Nelle applicazioni standard di elaborazione dei segnali bidimensionali, l'adozione di un filtro con dimensione spaziale pari a $1 \times 1$ non avrebbe alcun senso logico, riducendosi a una semplice operazione di riscalatura lineare del singolo pixel per una costante.

Nelle reti convoluzionali, tuttavia, l'operazione di \textbf{convoluzione 1x1} assume un'importanza cruciale. Ciò è dovuto al fatto che le CNN non operano su matrici bidimensionali isolate, ma su volumi tridimensionali in cui la terza dimensione rappresenta la profondità del canale (\emph{depth}).

Poiché i filtri di convoluzione si estendono obbligatoriamente per l'intera profondità del canale di input, un filtro $1 \times 1$ applicato a un volume di dimensioni $[H \times W \times D_i]$ esegue di fatto un prodotto scalare tridimensionale (\emph{3-dimensional dot product}) lungo l'asse dei canali per ciascuna coordinata spaziale dell'immagine.
Ad esempio, se l'input ha dimensioni $[32 \times 32 \times 3]$ (un'immagine RGB), l'applicazione di un filtro $1 \times 1$ corrisponde all'esecuzione di una combinazione lineare pesata dei 3 canali di colore per ciascun pixel spaziale, fungendo da efficace operatore di proiezione e di riduzione (o espansione) della dimensionalità dei canali.

\newpage

\section{Lo Stadio Di Attivazione (Rilevazione)}

Subito dopo l'applicazione del livello convoluzionale lineare (lo stadio di trasformazione affine), l'output grezzo viene immesso all'interno di una funzione di attivazione non lineare. Questo passaggio fondamentale prende il nome di \textbf{stadio di rilevamento} (\emph{detection stage}), poiché il suo scopo è quello di mappare le risposte lineari dei filtri su uno spazio di attivazione non lineare, discriminando la presenza o l'assenza di determinate caratteristiche visive (\emph{features}) locali.

Le principali funzioni di attivazione non lineari impiegate in questo stadio includono:
\begin{itemize}
    \item \textbf{ReLU} (\emph{Rectified Linear Unit}), definita come $f(x) = \max(0, x)$;
    \item \textbf{Leaky ReLU}, che preserva un piccolo gradiente anche per valori di input negativi;
    \item \textbf{Sigmoide} ($\sigma$), che schiaccia l'output nell'intervallo $(0, 1)$;
    \item \textbf{Tanh} (tangente iperbolica), che mappa i valori nell'intervallo $(-1, 1)$;
    \item \textbf{Maxout}, una generalizzazione che apprende una funzione lineare a tratti selezionando il massimo tra più combinazioni lineari dell'input.
\end{itemize}

\newpage

\section{Lo Stadio Di Pooling}

Il terzo stadio fondamentale di un blocco convoluzionale tipico è rappresentato dall'operazione di \textbf{pooling}.

\subsection{Definizione E Tipi Di Pooling}

La funzione di pooling sostituisce l'output della rete in corrispondenza di una determinata coordinata spaziale con una statistica sintetica calcolata sulle uscite adiacenti, racchiuse all'interno di un intorno rettangolare localizzato. Le funzioni di pooling più diffuse in letteratura sono:
\begin{itemize}
    \item \textbf{Max pooling}: Estrae il valore massimo registrato all'interno del vicinato rettangolare considerato. Rappresenta la variante più comune in quanto seleziona la risposta di attivazione più forte prodotta dal rilevatore.
    \item \textbf{Average pooling} (\emph{avg pooling}): Calcola il valore medio delle attivazioni all'interno della finestra locale.
    \item \textbf{Pooling basato su norma L2}: Calcola la norma euclidea dei valori contenuti nel vicinato rettangolare.
    \item \textbf{Media ponderata}: Assegna pesi differenti ai pixel dell'intorno in base alla loro distanza geometrica dal pixel centrale.
\end{itemize}

\subsection{Invarianza Alla Traslazione Locale}

In tutti gli scenari applicativi, il pooling svolge il compito essenziale di rendere la rappresentazione interna del modello approssimativamente \textbf{invariante a piccole traslazioni} dell'input spaziale. Dire che una rappresentazione è invariante alla traslazione locale significa che, qualora l'input subisse una traslazione di entità modesta, i valori risultanti dalla maggior parte delle unità di pooling non subirebbero variazioni.

Questa proprietà si rivela di estrema utilità quando l'interesse primario della rete risiede nel verificare la \emph{presenza} di una determinata caratteristica visiva, piuttosto che nel determinarne la collocazione spaziale con precisione micrometrica.
Ad esempio, nel compito di rilevare la presenza di un volto all'interno di un'immagine, non è necessario determinare la posizione esatta dei singoli fotorecettori oculari con accuratezza al singolo pixel; è sufficiente accertare la presenza di un occhio nella porzione sinistra del volto e di un secondo occhio nella porzione destra.

% [IMMAGINE DEL FILE DI SLIDE CON NOME "17-31"] Confronto tra i valori dello stadio del rilevatore (detector stage) e dello stadio di pooling (pooling stage) per mostrare l'invarianza locale — Slide 5 / Pagina 5 del PDF

Il meccanismo di invarianza locale può essere compreso analiticamente osservando il comportamento numerico di uno stadio di rilevamento e del rispettivo stadio di pooling (con finestra di ampiezza 3 e stride unitario) a fronte di uno slittamento dell'input:
\begin{enumerate}
    \item Nel primo caso, lo stadio di rilevazione produce i valori $[0.1, 1.0, 0.2, 0.1]$. Il max-pooling calcolato su una finestra mobile di 3 elementi restituisce i valori aggregati $[1.0, 1.0, 1.0, 0.2]$.
    \item Supponiamo ora che lo stesso segnale subisca una traslazione verso destra, portando lo stadio di rilevazione a registrare i valori $[0.3, 0.1, 1.0, 0.2]$. L'applicazione dello stesso operatore di max-pooling produce l'array $[0.3, 1.0, 1.0, 1.0]$.
\end{enumerate}
Come si può notare, nonostante lo spostamento spaziale del picco di attivazione (pari a $1.0$) all'interno dello stadio del rilevatore, i valori centrali restituiti dallo stadio di pooling rimangono costanti e pari a $1.0$, dimostrando concretamente la stabilità della rappresentazione rispetto a piccole variazioni di posizione.

\subsection{Apprendimento Di Invarianze Trasformazionali}

Mentre il pooling calcolato su regioni spaziali contigue produce invarianza alla sola traslazione, se l'operazione di pooling viene eseguita sulle uscite di convoluzioni parametrizzate separatamente, la rete acquisisce la capacità di apprendere a quali trasformazioni geometriche diventare invariante.

% [IMMAGINE DEL FILE DI SLIDE CON NOME "17-31"] Pooling su convoluzioni parametrizzate separatamente per apprendere l'invarianza alle rotazioni — Slide 6 / Pagina 6 del PDF

Questo principio è illustrato nell'estrazione di caratteristiche da simboli grafici (come il numero ``5'' scritto a mano):
\begin{itemize}
    \item Si predispongono tre filtri convoluzionali distinti, ciascuno addestrato a reagire a una specifica rotazione del carattere (ad esempio: inclinato a sinistra, orientamento standard, inclinato a destra).
    \item Ciascun filtro alimenta una propria unità di rilevazione lineare indipendente.
    \item Convogliando le risposte di questi tre rilevatori all'interno di una singola unità di pooling, quest'ultima produrrà un'attivazione elevata indipendentemente dal fatto che l'immagine di input contenga il carattere ``5'' in posizione verticale o ruotato. In questo modo, l'architettura apprende l'invarianza rispetto alla rotazione del simbolo.
\end{itemize}

\subsection{Efficienza Computazionale E Formule Dimensionali}

Poiché l'operazione di pooling sintetizza le risposte all'interno di un intero vicinato, è possibile ridurre il numero di unità di pooling necessarie rispetto alle unità di rilevamento originarie dello stadio precedente.

% [IMMAGINE DEL FILE DI SLIDE CON NOME "17-31"] Riduzione spaziale tramite pooling con stride maggiore di 1 — Slide 7 / Pagina 7 del PDF

Questo obiettivo si consegue distanziando le finestre di pooling di un passo (\emph{stride}) $k > 1$ anziché $1$ pixel. Così facendo, il livello successivo riceverà all'incirca $k$ volte meno input da elaborare, incrementando significativamente l'efficienza computazionale globale del sistema.

La configurazione standard e più diffusa prevede l'impiego di filtri di pooling con dimensione spaziale pari a $2 \times 2$ applicati con un passo (\emph{stride}) pari a 2. Questa combinazione dimezza la risoluzione dell'immagine sia in larghezza che in altezza, operando un campionamento in cui ogni operazione di massimo (nel caso del Max pooling) viene calcolata su un blocco di 4 pixel contigui per ciascun canale. Si noti che la profondità del volume non viene alterata.

In termini matematici generali, dato un volume di ingresso di dimensioni $H_i \times W_i \times C_i$, un filtro di pooling di dimensione spaziale $F_p$ e uno stride $S_p$, le dimensioni del volume di output sono regolate dalle seguenti relazioni:
\[
\boxed{W_o = \frac{W_i - F_p}{S_p} + 1}
\]
\[
\boxed{H_o = \frac{H_i - F_p}{S_p} + 1}
\]
\[
\boxed{C_o = C_i}
\]
È fondamentale evidenziare che il livello di pooling è una funzione matematica deterministica fissa e, pertanto, \textbf{non possiede alcun parametro libero} da apprendere durante la fase di addestramento.

\newpage

\section{Livelli Di Normalizzazione E Strati Completamente Connessi}

\subsection{Normalizzazione}

Inizialmente, i livelli di normalizzazione erano stati introdotti all'interno delle architetture convoluzionali per emulare alcuni schemi di inibizione laterale tipici del cervello biologico umano. Tuttavia, con l'evoluzione dei modelli di apprendimento profondo, tali soluzioni sono cadute in disuso ed è estremamente raro riscontrare l'impiego di questi strati di normalizzazione biologica all'interno delle moderne architetture CNN.

\subsection{Strati Completamente Connessi E Equivalenza Con I Livelli Convoluzionali}

I livelli completamente connessi (\emph{fully connected} o strati densi) rappresentano la struttura classica delle reti multistrato, in cui ogni singola unità neurale stabilisce una connessione con ogni elemento del volume di input.

La sola differenza concettuale tra un livello completamente connesso (\textbf{FC}) e un livello convoluzionale (\textbf{CONV}) risiede nel fatto che i neuroni del livello convoluzionale si connettono esclusivamente a una porzione localizzata dell'input (connettività locale) e condividono i propri parametri (condivisione dei pesi). Nonostante questa differenza, è sempre matematicamente possibile convertire un livello convoluzionale in uno completamente connesso, e viceversa.

\subsubsection{Conversione da CONV a FC}
Un livello convoluzionale può essere espresso in forma di livello completamente connesso. Per fare ciò, è sufficiente costruire una matrice dei pesi di grandi dimensioni che risulti prevalentemente sparsa (ossia composta da valori pari a zero ovunque, tranne che in corrispondenza dei blocchi definiti dalla connettività locale del kernel) e in cui i valori non nulli all'interno dei diversi blocchi siano vincolati a essere identici tra loro, riflettendo la condivisione dei parametri tipica della convoluzione.

\subsubsection{Conversione da FC a CONV}
Un livello completamente connesso può essere convertito in un livello convoluzionale equivalente. Per fare ciò, è sufficiente configurare un filtro di convoluzione la cui dimensione spaziale sia esattamente uguale a quella del volume di input spaziale ($K_x = H_i, K_y = W_i$) e impostare il numero totale di filtri pari al numero di unità nascoste (\emph{hidden units}) dello strato FC che si desidera emulare.

Per chiarire questo passaggio, analizziamo un esempio pratico:
\begin{itemize}
    \item Supponiamo di avere un volume di input di dimensioni $[12 \times 12 \times 64]$.
    \item Desideriamo sostituire un livello completamente connesso avente $1024$ unità nascoste.
    \item A tale scopo, è sufficiente introdurre un livello convoluzionale composto da $1024$ filtri di dimensione spaziale $12 \times 12$, applicati con uno stride pari a 1 e un padding pari a 0.
\end{itemize}
Verifichiamo analiticamente le dimensioni spaziali dell'output applicando la formula nota:
\[
H_o = \frac{12 - 12 + 2(0)}{1} + 1 = 1
\]
\[
W_o = \frac{12 - 12 + 2(0)}{1} + 1 = 1
\]
Poiché il numero di filtri (canali di output) è pari a $1024$, il volume finale di attivazione risulterà essere $[1 \times 1 \times 1024]$. Questo volume è matematicamente equivalente a un vettore unidimensionale di $1024$ elementi, dimostrando la perfetta equivalenza tra le due formulazioni.

\newpage

\section{Pattern Di Progettazione Dell'Architettura}

\subsection{Regole Di Composizione Degli Strati}

La progettazione di una rete convoluzionale efficace si basa sul concatenamento logico di blocchi ripetuti secondo un'alternanza ben definita. Il pattern architetturale più comune prevede la sovrapposizione di diversi blocchi composti da livelli convoluzionali seguiti da funzioni ReLU, seguiti a intervalli regolari da livelli di pooling per ridurre la risoluzione spaziale, ripetendo questo schema fino a quando l'immagine non è stata ridotta a dimensioni minime. A questo punto, si effettua la transizione verso i livelli completamente connessi per mappare le caratteristiche finali sulle classi di output.

Il pattern generale può essere formalizzato tramite la seguente notazione:
\[
\boxed{\text{INPUT} \rightarrow [[\text{CONV} \rightarrow \text{RELU}]^{*N} \rightarrow \text{POOL?}]^{*M} \rightarrow [\text{FC} \rightarrow \text{RELU}]^{*K} \rightarrow \text{FC}}
\]
All'interno di questa formulazione:
\begin{itemize}
    \item Il simbolo $*$ denota la ripetizione del blocco racchiuso tra parentesi;
    \item $\text{POOL?}$ indica la presenza opzionale di un livello di pooling;
    \item I vincoli tipici sui parametri di ripetizione sono: $N \ge 0$ (con un valore solitamente contenuto entro $N \le 3$), $M \ge 0$, $K \ge 0$ (con un valore solitamente inferiore a $K < 3$).
\end{itemize}

Alcune istanziazioni comuni di questo schema includono:
\begin{itemize}
    \item $\text{INPUT} \rightarrow \text{FC}$: Rappresenta la configurazione di un classificatore lineare semplice (dove $N = M = K = 0$).
    \item $\text{INPUT} \rightarrow \text{CONV} \rightarrow \text{RELU} \rightarrow \text{FC}$
    \item $\text{INPUT} \rightarrow [\text{CONV} \rightarrow \text{RELU} \rightarrow \text{POOL}]^{*2} \rightarrow \text{FC} \rightarrow \text{RELU} \rightarrow \text{FC}$: Configurazione classica che prevede un singolo livello convoluzionale interposto prima di ogni operazione di pooling.
    \item $\text{INPUT} \rightarrow [\text{CONV} \rightarrow \text{RELU} \rightarrow \text{CONV} \rightarrow \text{RELU} \rightarrow \text{POOL}]^{*3} \rightarrow [\text{FC} \rightarrow \text{RELU}]^{*2} \rightarrow \text{FC}$: Architettura in cui si sovrappongono più livelli convoluzionali prima di applicare il pooling. Questa scelta progettuale è indicata per reti profonde, poiché l'accumulo di più livelli convoluzionali consecutivi consente di estrarre e sviluppare caratteristiche visive estremamente complesse prima di applicare l'operazione di pooling, che per sua natura è distruttiva sul piano delle informazioni spaziali di dettaglio.
\end{itemize}

\subsection{Confronto Tra Stack Di Piccoli Filtri E Filtri Singoli Grandi}

Una regola fondamentale nella moderna progettazione di reti convoluzionali suggerisce di prediligere una pila (\emph{stack}) di livelli convoluzionali con filtri di dimensioni ridotte rispetto a un singolo livello convoluzionale basato su un filtro con un campo ricettivo ampio.

Per comprendere questa scelta, analizziamo il confronto tra l'applicazione di tre strati convoluzionali consecutivi con filtri spaziali $3 \times 3$ (intervallati da funzioni di attivazione non lineari) e l'applicazione di un singolo strato convoluzionale basato su un filtro $7 \times 7$.

\subsubsection{Analisi del Campo Ricettivo}
\begin{itemize}
    \item Ciascun neurone appartenente al primo livello convoluzionale $3 \times 3$ possiede una vista di area $3 \times 3$ sull'input originario.
    \item Un neurone del secondo livello convoluzionale $3 \times 3$ possiede una vista $3 \times 3$ sulle attivazioni del livello precedente. Di conseguenza, per transizione geometrica, esso possiede un campo ricettivo equivalente a un'area $5 \times 5$ sull'input iniziale.
    \item Un neurone del terzo livello convoluzionale $3 \times 3$ possiede a sua volta una vista $3 \times 3$ sul secondo livello, che si traduce in un campo ricettivo finale equivalente a un'area di $7 \times 7$ sull'input originario.
\end{itemize}
Dal punto di vista dell'estensione del campo ricettivo spaziale, i due approcci sono geometricamente equivalenti. Tuttavia, la soluzione basata sulla pila di tre filtri piccoli presenta due enormi vantaggi.

\subsubsection{Maggiore Espressività}
I tre strati convoluzionali contengono al loro interno tre funzioni di attivazione non lineari distinte (una dopo ciascun livello). Al contrario, la soluzione a filtro singolo grande dispone di una sola non linearità. La presenza di più funzioni non lineari consente alla rete di apprendere relazioni matematiche e caratteristiche visive molto più articolate ed espressive.

\subsubsection{Riduzione del Numero dei Parametri}
Ipotizziamo che tutti i volumi di attivazione intermedi e di input possiedano un numero di canali pari a $C$. Calcoliamo il numero di parametri associati a entrambe le configurazioni:
\begin{itemize}
    \item \textbf{Filtro singolo $7 \times 7$}: Un singolo livello con $C$ filtri di dimensioni $7 \times 7$ agenti su un canale di ingresso di profondità $C$ introduce un numero di parametri pari a:
    \[
    C \times (7 \times 7 \times C) = 49C^2
    \]
    \item \textbf{Stack di tre filtri $3 \times 3$}: Ognuno dei tre livelli convoluzionali è composto da $C$ filtri di dimensioni $3 \times 3$ applicati su canali di profondità $C$. Il numero totale di parametri della pila è dato da:
    \[
    3 \times (C \times (3 \times 3 \times C)) = 27C^2
    \]
\end{itemize}
Il confronto evidenzia come la pila di tre piccoli filtri comporti una drastica riduzione dei parametri da addestrare ($27C^2$ contro $49C^2$), migliorando notevolmente l'efficienza statistica del modello e contrastando il rischio di overfitting.

Tuttavia, questa scelta presenta uno svantaggio pratico: durante la fase di addestramento tramite retropropagazione (\emph{backpropagation}), la rete sarà costretta a memorizzare in GPU tutte le mappe di attivazione intermedie prodotte dai livelli convoluzionali della pila, comportando un maggiore consumo di memoria di lavoro rispetto alla soluzione a filtro singolo.

\subsection{Linee Guida Pratiche E Regole Empiriche Per Il Dimensionamento}

Per agevolare il dimensionamento ottimale delle architetture convoluzionali, si fa affidamento ad alcune regole pratiche consolidate:

\begin{itemize}
    \item \textbf{Livello di Input}: Le dimensioni spaziali dell'immagine in ingresso dovrebbero essere divisibili per 2 un numero elevato di volte. I valori più comuni includono 32 (es. CIFAR-10), 64, 96 (es. STL-10), 224 (es. ImageNet), 384 e 512.
    \item \textbf{Livelli Convoluzionali}: Si consiglia l'impiego di filtri di dimensioni ridotte (pari a $3 \times 3$ o al massimo $5 \times 5$), impostando uno stride unitario ($S = 1$) e calibrando lo zero-padding in modo tale da non alterare la risoluzione spaziale del volume di input. Applicando la regola generale:
    \[
    P = \frac{F - 1}{2}
    \]
    si deduce che per un filtro con dimensione $F = 3$ occorre impostare un padding $P = 1$, mentre per un filtro con dimensione $F = 5$ occorre impostare $P = 2$. Filtri di dimensioni maggiori (ad esempio $7 \times 7$ o superiori) vengono adottati quasi esclusivamente nel primissimo livello convoluzionale, ovvero quello a diretto contatto con l'immagine di ingresso.
    \item \textbf{Livelli di Pooling}: Hanno il compito esclusivo di effettuare il sottocampionamento spaziale del volume. La configurazione standard prevede il max-pooling con filtri di dimensioni $2 \times 2$ ($F = 2$) e stride pari a 2 ($S = 2$). Questa impostazione dimezza sia la larghezza che l'altezza dell'immagine, scartando esattamente il $75\%$ delle attivazioni del volume originario. Una seconda opzione, leggermente meno comune, consiste nell'utilizzare filtri di pooling $3 \times 3$ con stride pari a 2. È altamente sconsigliato l'impiego di filtri di pooling con dimensioni superiori a 3, poiché l'operazione risulterebbe eccessivamente aggressiva e causerebbe una perdita distruttiva di informazioni, compromettendo le prestazioni complessive del modello.
\end{itemize}

\subsection{Compromessi Architetturali E Vincoli Di Memoria}

Nonostante l'efficacia delle regole di dimensionamento descritte, la quantità di memoria fisica richiesta dalle attivazioni e dai relativi gradienti può crescere molto rapidamente, soprattutto nei primi livelli della rete, scontrandosi con i limiti strutturali delle GPU.

Per illustrare questa problematica, analizziamo un caso concreto:
\begin{itemize}
    \item \textbf{Input}: Un'immagine di dimensioni $224 \times 224 \times 3$.
    \item \textbf{Architettura iniziale}: Tre livelli convoluzionali consecutivi composti ciascuno da $64$ filtri di dimensioni $3 \times 3$, applicati con un padding pari a 1 per preservare le dimensioni dell'input.
    \item \textbf{Calcolo delle attivazioni}: Ciascuno dei tre livelli genera un volume di attivazione intermedio avente dimensioni pari a $[224 \times 224 \times 64]$.
    \item \textbf{Conteggio totale}: Il numero complessivo di singole attivazioni da memorizzare per una singola immagine è pari a:
    \[
    3 \times (224 \times 224 \times 64) = 9.633.792 \approx 10 \text{ milioni di attivazioni}
    \]
    Ipotizzando l'uso del formato standard a virgola mobile a 32 bit (in cui ogni float occupa 4 byte), la memorizzazione di queste sole attivazioni richiede circa $38.5\text{ MB}$ di memoria. Poiché durante l'addestramento è necessario conservare sia le attivazioni della fase di \emph{forward} sia i rispettivi gradienti per la fase di \emph{backward}, il requisito di memoria raddoppia, arrivando a circa \textbf{72 MB} di memoria per singola immagine processata.
\end{itemize}

Poiché le schede grafiche sono spesso soggette a colli di bottiglia legati alla memoria fisica disponibile, si rende necessario scendere a compromessi progettuali. Tali compromessi vengono applicati quasi esclusivamente all'altezza del **primissimo livello convoluzionale** della rete, allo scopo di abbattere rapidamente la risoluzione spaziale dell'input prima di addentrarsi nei livelli successivi.

Due esempi storici di questo tipo di compromesso sono rappresentati dalle seguenti architetture:
\begin{itemize}
    \item \textbf{ZF Net}: Utilizza nel primo livello filtri di dimensioni spaziali pari a $7 \times 7$ applicati con uno stride pari a 2.
    \item \textbf{AlexNet}: Adotta nel primo livello filtri con dimensioni spaziali pari a $11 \times 11$ applicati con uno stride pari a 4.
\end{itemize}
L'adozione di passi di scorrimento elevati ($S \ge 2$) e di filtri spazialmente ampi al debutto della rete consente di ridurre drasticamente la mappa spaziale iniziale, abbattendo la richiesta di memoria per i livelli successivi e rendendo l'addestramento sostenibile sui dispositivi hardware.

% AA Slides 11.3
\section{Transfer Learning E Fine Tuning}
Nel contesto applicativo dell'apprendimento profondo, addestrare un'intera rete neurale convoluzionale (\textbf{ConvNet} o \textbf{CNN}) partendo da zero (\emph{from scratch}) con un'inizializzazione casuale dei pesi è una pratica estremamente rara. Questa difficoltà risiede principalmente nella disponibilità dei dati: è infatti raro disporre di un insieme di dati (\emph{dataset}) di dimensioni sufficientemente grandi da consentire una convergenza stabile senza incorrere in gravi fenomeni di \emph{overfitting}.

Al contrario, l'approccio standard prevede il pre-addestramento (\emph{pretraining}) di una ConvNet su un dataset estremamente vasto, come \textbf{ImageNet} (che contiene circa $1{,}2$ milioni di immagini suddivise in $1000$ categorie distinte). Successivamente, la rete convoluzionale così addestrata viene impiegata per il compito specifico di interesse (\emph{task of interest}), agendo o come inizializzazione per i pesi o come un estrattore di caratteristiche a pesi congelati (\textbf{estrattore di caratteristiche fisse}).

\subsection{I Tre Scenari Principali Di Transfer Learning}
L'apprendimento per trasferimento (\textbf{Transfer Learning}) si articola principalmente in tre scenari operativi.

\subsubsection{La Rete Convoluzionale Come Estrattore Di Caratteristiche Fisse}
Il primo scenario prevede l'utilizzo della ConvNet come estrattore di caratteristiche fisse (\emph{fixed feature extractor}). In questo caso, si prende una rete convoluzionale pre-addestrata su ImageNet e si rimuove l'ultimo livello completamente connesso (\textbf{fully-connected layer} o \textbf{FC}), i cui output originali rappresentano i punteggi delle $1000$ classi del dataset di pre-addestramento.

La porzione restante della rete viene trattata come un estrattore statico di caratteristiche per il nuovo dataset, mantenendo i pesi congelati (\emph{frozen}) durante tutta la fase di addestramento. Ad esempio, utilizzando l'architettura \textbf{AlexNet}, questo processo consente di calcolare un vettore a $4096$ dimensioni (\textbf{vettore $4096$-D}) per ciascuna immagine in ingresso. Questo vettore racchiude le attivazioni dello strato nascosto immediatamente precedente al classificatore finale. Tali rappresentazioni prendono il nome di \textbf{codici CNN} (\emph{CNN codes}).

Per garantire prestazioni ottimali, è fondamentale che questi codici subiscano una sogliatura tramite la funzione di attivazione \textbf{ReLU} (ossia vengano filtrati ponendo a zero i valori negativi, \emph{ReLUd}), qualora tale sogliatura fosse presente durante la fase di addestramento iniziale su ImageNet (circostanza che si verifica nella quasi totalità dei casi). Una volta estratti i codici a $4096$ dimensioni per tutte le immagini del nuovo dataset, si procede all'addestramento di un classificatore lineare semplice, come una \textbf{SVM lineare} (Support Vector Machine) o un classificatore \textbf{Softmax}, dedicato al nuovo compito.

% [IMMAGINE DEL FILE DI SLIDE CON NOME "Slides Transfer Learning"] Schema di VGG16 utilizzata come estrattore di caratteristiche fisse: i blocchi inferiori mantengono i pesi congelati mentre l'ultimo livello fully-connected viene rimpiazzato e addestrato da zero con un nuovo classificatore (Softmax o L2 Loss) — Slide 2 / Pagina 2 del PDF

\subsubsection{Sintonizzazione Fine Della Rete Convoluzionale (\emph{Fine-Tuning})}
Il secondo scenario consiste nella sintonizzazione fine (\textbf{\emph{fine-tuning}}) della ConvNet. Questa strategia non si limita a sostituire e riaddestrare il classificatore posto in cima alla rete, ma prevede l'aggiornamento dei pesi della rete pre-addestrata continuando il processo di retropropagazione dell'errore (\textbf{\emph{backpropagation}}) sul nuovo dataset.

È possibile procedere sintonizzando tutti i livelli della ConvNet oppure, per mitigare i problemi di \emph{overfitting}, congelare i primi strati della rete e aggiornare solo una porzione superiore (strati più profondi o ad alto livello). Questa scelta è motivata dall'evidenza empirica che i primi strati di una rete convoluzionale apprendono caratteristiche generiche (\textbf{caratteristiche generiche}), come rilevatori di bordi (\emph{edge detectors}) o di macchie di colore (\emph{color blob detectors}), utili per una vasta gamma di compiti visivi. Al contrario, gli strati successivi diventano progressivamente più specifici rispetto ai dettagli delle classi contenute nel dataset originale. Ad esempio, nel caso di ImageNet, che include numerose razze canine, una porzione significativa della capacità rappresentativa dei livelli più profondi della ConvNet potrebbe essere specializzata esclusivamente nella differenziazione tra specifiche razze di cani.

% [IMMAGINE DEL FILE DI SLIDE CON NOME "Slides Transfer Learning"] Schema di VGG16 configurata per il fine-tuning: i livelli convoluzionali profondi vengono aggiornati tramite retropropagazione insieme al nuovo classificatore lineare — Slide 3 / Pagina 3 del PDF

\subsubsection{Modelli Pre-Addestrati (\emph{Pretrained Models})}
Dato che l'addestramento da zero di moderne reti convoluzionali su ImageNet richiede dalle due alle tre settimane di calcolo distribuito su più GPU, è prassi comune rilasciare pubblicamente i \emph{checkpoint} finali dei pesi delle reti. Questo consente ad altri utenti di usufruire dei modelli già pronti per compiti di sintonizzazione fine. Ad esempio, la libreria \textbf{Caffe} mette a disposizione un archivio condiviso noto come \textbf{Model Zoo}, all'interno del quale la comunità scientifica e gli sviluppatori possono condividere i pesi delle proprie reti.

% [IMMAGINE DEL FILE DI SLIDE CON NOME "Slides Transfer Learning"] Schema di VGG16 che illustra l'importazione di pesi pre-addestrati scaricati dal web per essere impiegati come base di partenza — Slide 4 / Pagina 4 del PDF

\newpage

\subsection{Quando E Come Eseguire Il Fine-Tuning}
La decisione su quale tipo di apprendimento per trasferimento adottare di fronte a un nuovo dataset dipende essenzialmente da due fattori cardine:
\begin{enumerate}
    \item \textbf{La dimensione del nuovo dataset}: classificabile come piccolo (\emph{small}) o grande (\emph{big}).
    \item \textbf{La similarità con il dataset originale}: ad esempio, se il nuovo dataset è simile a ImageNet in termini di contenuto delle immagini e di classi, oppure se è molto differente (come nel caso di immagini acquisite al microscopio).
\end{enumerate}

Tenendo a mente che le caratteristiche estratte dalla ConvNet sono generiche nei primi strati e altamente specifiche del dataset originale negli strati più profondi, si possono delineare quattro scenari di riferimento per guidare la scelta:

\begin{enumerate}[leftmargin=*]
    \item \textbf{Nuovo dataset piccolo e simile al dataset originale:}
    Essendo il dataset di piccole dimensioni, l'esecuzione del fine-tuning sull'intera rete è sconsigliata a causa dell'elevato rischio di \emph{overfitting}. Tuttavia, poiché il dataset è simile a quello di pre-addestramento, possiamo assumere che le caratteristiche di alto livello estratte dalla rete siano rilevanti anche per il nuovo compito. Di conseguenza, la strategia ottimale consiste nell'addestrare un classificatore lineare direttamente sui codici CNN estratti, mantenendo i pesi della ConvNet pre-addestrata completamente congelati durante l'addestramento.

    \item \textbf{Nuovo dataset grande e simile al dataset originale:}
    Disponendo di una grande quantità di dati, si riduce notevolmente il rischio di \emph{overfitting}. È quindi possibile procedere con fiducia alla sintonizzazione fine dei pesi dell'intera rete (\emph{fine-tuning} completo), partendo dall'inizializzazione fornita dal modello pre-addestrato.

    \item \textbf{Nuovo dataset piccolo ma molto differente dal dataset originale:}
    La ridotta dimensione dei dati impone nuovamente di evitare il fine-tuning completo per prevenire l'\emph{overfitting}. Tuttavia, poiché il dataset è molto differente, non è ottimale addestrare il classificatore lineare a partire dalle attivazioni degli strati finali della rete, in quanto queste contengono caratteristiche troppo specifiche del dataset originale. La soluzione più efficace consiste nell'addestrare un classificatore lineare (come una SVM) a partire dalle attivazioni di strati intermedi e più precoci della rete, scartando i livelli finali di estrazione delle caratteristiche al fine di sfruttare unicamente le rappresentazioni più generiche.

    \item \textbf{Nuovo dataset grande e molto differente dal dataset originale:}
    Data la grande dimensione del dataset, si potrebbe teoricamente considerare l'addestramento di una nuova ConvNet partendo da zero. Ciononostante, nella pratica si rivela quasi sempre estremamente vantaggioso inizializzare i pesi della rete con quelli di un modello pre-addestrato. In questa situazione, la disponibilità di dati consente di effettuare una sintonizzazione fine sull'intera architettura, avendo cura di adottare tassi di apprendimento adeguati.
\end{enumerate}

\newpage

\subsection{Consigli Pratici Per L'Apprendimento Per Trasferimento}
Durante l'implementazione pratica del Transfer Learning, vi sono alcuni aspetti aggiuntivi da considerare:

\begin{itemize}
    \item \textbf{Vincoli derivanti dai modelli pre-addestrati:}
    L'impiego di una rete pre-addestrata introduce alcuni vincoli sull'architettura che è possibile adottare per il nuovo compito. Ad esempio, non è consentito rimuovere arbitrariamente strati convoluzionali dalla rete originale senza alterare la coerenza dei flussi dei dati.

    Tuttavia, alcune modifiche strutturali sono semplici da implementare. Grazie alla proprietà di \textbf{condivisione dei parametri} (\emph{parameter sharing}), è possibile applicare una rete pre-addestrata a immagini di dimensioni spaziali differenti rispetto a quelle usate nel pretraining. Questa flessibilità è intrinseca nei \textbf{livelli convoluzionali e di pooling}, la cui funzione di propagazione in avanti (\emph{forward function}) è indipendente dalla dimensione spaziale del volume di input, purché i passi (\emph{strides}) siano compatibili con le dimensioni geometriche.

    Per quanto concerne i \textbf{livelli completamente connessi} (FC), questa flessibilità può essere preservata poiché un livello FC può essere formalmente convertito in un livello convoluzionale. Ad esempio, in AlexNet, il volume di pooling finale che precede il primo livello FC ha dimensioni pari a $6 \times 6 \times 512$. Di conseguenza, il livello FC che elabora tale volume equivale a un livello convoluzionale avente una dimensione del campo ricettivo (\emph{receptive field}) pari a $6 \times 6$, applicato con un'imbottitura (\emph{padding}) pari a $0$.

    \item \textbf{Gestione dei tassi di apprendimento (\emph{Learning Rates}):}
    È una pratica comune impostare un tasso di apprendimento significativamente più piccolo per i pesi della ConvNet che vengono sintonizzati, rispetto al tasso di apprendimento applicato al nuovo classificatore lineare (il quale viene invece inizializzato in modo del tutto casuale per calcolare i punteggi delle classi sul nuovo dataset). Questo accorgimento si giustifica con il fatto che i pesi pre-addestrati della ConvNet sono già posizionati in una regione ottimale dello spazio dei parametri; pertanto, si vuole evitare di distorcerli troppo rapidamente o eccessivamente, in particolar modo durante le prime fasi in cui il sovrastante classificatore lineare si sta ancora stabilizzando a partire dalla sua inizializzazione casuale.
\end{itemize}

% AA Slides 11.4
\section{Introduzione E Informazioni Amministrative}

Benvenuti alla Lezione 11 del corso, intitolata \textbf{Detection and Segmentation} (Rilevamento e Segmentazione). Di seguito vengono presentate le informazioni organizzative del corso e l'introduzione alla piattaforma didattica HyperQuest.

\subsection{Informazioni Organizzative}

\begin{itemize}
    \item I compiti scritti di metà corso (\emph{midterms}) sono attualmente in fase di correzione. Si invitano gli studenti a non discutere dei quesiti d'esame fino alla settimana successiva, poiché alcuni studenti devono ancora sostenerlo.
    \item Il secondo assegnamento (\emph{A2}) è attualmente in fase di correzione.
    \item Le scadenze per le pietre miliari del progetto di corso (\emph{project milestones}) sono fissate per martedì 16 maggio.
\end{itemize}

\subsection{La Piattaforma Di Apprendimento HyperQuest}

Viene introdotto l'esperimento didattico \textbf{HyperQuest}, una piattaforma di gioco interattiva volta ad approfondire l'intuizione e la comprensione pratica degli studenti riguardo all'addestramento e alla regolazione degli iperparametri nelle reti neurali.

% [IMMAGINE DEL FILE DI SLIDE CON NOME "11-3"] Schermata di login della piattaforma didattica HyperQuest — Slide 3 / Pagina 3 del PDF

% [IMMAGINE DEL FILE DI SLIDE CON NOME "11-4"] Sondaggio iniziale sui dati dello studente e sull'esperienza pregressa — Slide 4 / Pagina 4 del PDF

L'interfaccia richiede una registrazione iniziale e la compilazione di un breve sondaggio relativo allo status accademico dello studente (studente triennale, magistrale, dottorato o non iscritto a un corso di laurea) e al livello di esperienza nell'addestramento di reti neurali (espresso in anni di esperienza e tipologie di architetture utilizzate, come reti completamente connesse, CNN, reti ricorrenti, reti per compiti di visione o di NLP).

% [IMMAGINE DEL FILE DI SLIDE CON NOME "11-5"] Regolamento e istruzioni di base per l'utilizzo di HyperQuest — Slide 5 / Pagina 5 del PDF

Il regolamento e le istruzioni operative di HyperQuest prevedono le seguenti fasi:
\begin{itemize}
    \item A ogni studente viene fornito un dataset casuale. L'obiettivo è addestrare una rete neurale per compiti di classificazione, cercando di ottenere la massima accuratezza di validazione (\emph{validation accuracy}) possibile.
    \item Nella prima fase, lo studente seleziona la configurazione iniziale della rete.
    \item Nella seconda fase, si monitora il processo di addestramento e si ha la possibilità di regolare manualmente gli iperparametri a ogni epoca.
\end{itemize}

% [IMMAGINE DEL FILE DI SLIDE CON NOME "11-7"] Configurazione iniziale dei parametri e statistiche del dataset in HyperQuest — Slide 7 / Pagina 7 del PDF

All'avvio, la piattaforma mostra le statistiche del dataset generato:
\begin{itemize}
    \item \textbf{Classi}: 10.
    \item \textbf{Dimensione dei dati di input}: $[3, 32, 32]$.
    \item \textbf{Esempi per split}: 8500 campioni per il \emph{training set} e 1500 per il \emph{validation set}.
\end{itemize}

Lo studente deve quindi definire la configurazione iniziale della rete selezionando tra i seguenti valori discreti:
\begin{itemize}
    \item \textbf{Larghezza della rete CNN} (\emph{CNN network width}): 32, 64 o 128 filtri.
    \item \textbf{Tasso di apprendimento} (\emph{Learning rate}): 0.1, 0.01 o 0.001.
    \item \textbf{Profondità della rete CNN} (\emph{CNN network depth}): 2, 4 o 8 strati.
    \item \textbf{Tasso di Dropout} (\emph{Dropout rate}): 0 o 0.5.
\end{itemize}

% [IMMAGINE DEL FILE DI SLIDE CON NOME "11-8"] Interfaccia di monitoraggio dell'addestramento e grafici di loss/accuracy alla prima epoca — Slide 8 / Pagina 8 del PDF

% [IMMAGINE DEL FILE DI SLIDE CON NOME "11-9"] Interfaccia di monitoraggio con cronologia delle configurazioni e regolazione attiva degli iperparametri — Slide 9 / Pagina 9 del PDF

Durante l'addestramento, che si sviluppa fino a un massimo di 60 azioni (con un limite di 30 epoche totali), lo studente monitora le metriche in tempo reale. Le azioni disponibili comprendono:
\begin{itemize}
    \item \textbf{Continue}: prosegue l'addestramento per le epoche successive mantenendo o modificando i parametri.
    \item \textbf{Done training}: conclude la sessione di addestramento salvando il risultato migliore.
    \item \textbf{Revert}: consente di annullare gli ultimi passaggi e tornare indietro nello stato di addestramento (fino a un massimo di 5 volte consecutive).
\end{itemize}

La regolazione attiva degli iperparametri permette di modificare dinamicamente i seguenti valori a ogni passo di addestramento:
\begin{itemize}
    \item Regolarizzazione tramite decadimento dei pesi (\emph{weight decay}).
    \item Profondità della rete CNN.
    \item Tasso di dropout.
    \item Tasso di apprendimento.
    \item Larghezza della rete CNN.
\end{itemize}

I grafici cartesiani mostrano l'andamento della funzione di perdita (\emph{loss}) e dell'accuratezza (\emph{accuracy}) sia per il set di addestramento che per quello di validazione, mantenendo una cronologia dettagliata di tutte le configurazioni testate. Ulteriori dettagli operativi su questa attività verranno pubblicati sulla piattaforma Piazza nel pomeriggio.

\newpage

\section{Ripasso Della Lezione Precedente: Reti Neurali Ricorrenti}

\subsection{Architetture E Applicazioni Delle Reti Ricorrenti}

Nella lezione precedente sono state presentate le \textbf{Reti Neurali Ricorrenti} (RNN), modelli in grado di elaborare sequenze di dati di lunghezza variabile grazie all'introduzione di uno stato interno che si aggiorna in modo ricorsivo. A seconda della struttura dei dati di input e di output, le reti ricorrenti possono essere configurate secondo diversi schemi architetturali:

% [IMMAGINE DEL FILE DI SLIDE CON NOME "11-11"] Schemi delle architetture ricorrenti (one-to-one, one-to-many, many-to-one, many-to-many) e diagramma di un modello linguistico a livello di carattere — Slide 11 / Pagina 11 del PDF

\begin{itemize}
    \item \textbf{One-to-one}: corrisponde alla rete feedforward classica, in cui un singolo input produce un singolo output (ad esempio, la classificazione standard di immagini).
    \item \textbf{One-to-many}: a fronte di un singolo input, la rete genera una sequenza di elementi in uscita (tipico nell'\emph{image captioning}, dove da un'immagine si ricava una descrizione testuale).
    \item \textbf{Many-to-one}: una sequenza di input viene elaborata per restituire un singolo output finale (ad esempio, l'analisi del sentimento di un testo).
    \item \textbf{Many-to-many (asincrono)}: una sequenza di input viene interamente codificata prima che la rete inizi a generare una sequenza di output (utilizzato nella traduzione automatica \emph{sequence-to-sequence}).
    \item \textbf{Many-to-many (sincrono)}: ogni elemento della sequenza di input produce contemporaneamente un elemento della sequenza di output (ad esempio, la classificazione dei singoli frame di un video).
\end{itemize}

Un esempio concreto è il \textbf{modello linguistico a livello di carattere} (\emph{character-level language model}). Data una sequenza di caratteri di input, il modello stima la probabilità del carattere successivo. Fornendo ad esempio la parola "hello" una lettera alla volta:
\begin{itemize}
    \item Ciascun carattere di input viene codificato tramite un vettore \emph{one-hot}.
    \item Lo strato nascosto calcola l'aggiornamento dello stato tramite la relazione ricorsiva:
    \[
    \bm{h}_t = \tanh(\bm{W}_{hh} \bm{h}_{t-1} + \bm{W}_{xh} \bm{x}_t)
    \]
    \item Lo strato di output calcola i punteggi (\emph{scores}) non normalizzati per il carattere successivo mediante una proiezione lineare:
    \[
    \bm{y}_t = \bm{W}_{hy} \bm{h}_t
    \]
\end{itemize}

% [IMMAGINE DEL FILE DI SLIDE CON NOME "11-12"] Esempi di testi generati da una rete ricorrente addestrata su LaTeX, Shakespeare e codice Linux — Slide 12 / Pagina 12 del PDF

% [IMMAGINE DEL FILE DI SLIDE CON NOME "11-13"] Esempi di image captioning prodotti tramite reti ricorrenti — Slide 13 / Pagina 13 del PDF

Questo meccanismo di generazione consente alle reti ricorrenti di produrre testi coerenti e strutturati, come codice sorgente C del kernel Linux, dialoghi in stile shakespeariano o articoli scientifici in formato LaTeX. Inoltre, nel compito dell'image captioning (come descritto nel lavoro di Karpathy et al., 2015), l'integrazione di una CNN per l'estrazione delle caratteristiche visive e di una RNN permette di generare descrizioni testuali accurate a partire da immagini reali.

\subsection{Dal Modello Vanilla Alla Long Short-Term Memory}

Sebbene il modello ricorrente classico (noto come \textbf{Vanilla RNN}) presenti una formulazione elegante, esso mostra limiti significativi nell'apprendimento di dipendenze a lungo termine a causa dei fenomeni di dissolvimento ed esplosione del gradiente. La struttura della Vanilla RNN prevede infatti la concatenazione dello stato precedente $\bm{h}_{t-1}$ con il vettore di input corrente $\bm{x}_t$, seguita da una trasformazione lineare e dall'applicazione di una non-linearità tangente iperbolica ($\tanh$):
\[
\bm{h}_t = \tanh\left(\bm{W} \begin{bmatrix} \bm{h}_{t-1} \\ \bm{x}_t \end{bmatrix}\right)
\]

% [IMMAGINE DEL FILE DI SLIDE CON NOME "11-14"] Confronto architetturale tra Vanilla RNN e LSTM — Slide 14 / Pagina 14 del PDF

Per ovviare a tali problemi, l'architettura \textbf{Long Short-Term Memory} (LSTM) introduce uno stato di cella interno $\bm{c}_t$ che funge da memoria a lungo termine, preservando il flusso del gradiente in modo additivo. L'interazione con lo stato di cella è regolata da quattro porte (\emph{gates}) logiche, calcolate congiuntamente a partire dal vecchio stato nascosto $\bm{h}_{t-1}$ e dall'input $\bm{x}_t$:
\[
\begin{pmatrix} f \\ i \\ g \\ o \end{pmatrix} = \begin{pmatrix} \sigma \\ \sigma \\ \tanh \\ \sigma \end{pmatrix} \left( \bm{W} \begin{bmatrix} \bm{h}_{t-1} \\ \bm{x}_t \end{bmatrix} \right)
\]

Ciascun vettore svolge una funzione specifica all'interno della cella:
\begin{itemize}
    \item \textbf{Porta di dimenticanza} ($f$, \emph{forget gate}): determina quanta parte dello stato di cella precedente $\bm{c}_{t-1}$ debba essere conservata.
    \item \textbf{Porta di ingresso} ($i$, \emph{input gate}): decide quali elementi della nuova informazione debbano essere scritti nello stato di cella.
    \item \textbf{Candidato di aggiornamento} ($g$, \emph{gate gate}): rappresenta i nuovi valori candidati, calcolati tramite la funzione $\tanh$, pronti per essere inseriti nello stato di cella.
    \item \textbf{Porta di uscita} ($o$, \emph{output gate}): stabilisce quale porzione dello stato di cella debba essere riflessa nello stato nascosto $\bm{h}_t$.
\end{itemize}

L'aggiornamento dello stato di cella avviene combinando linearmente la vecchia memoria e il nuovo candidato:
\[
\bm{c}_t = f \odot \bm{c}_{t-1} + i \odot g
\]

Infine, lo stato nascosto corrente $\bm{h}_t$ viene ricavato applicando una tangente iperbolica allo stato di cella e filtrandone il risultato tramite la porta di uscita:
\[
\boxed{\bm{h}_t = o \odot \tanh(\bm{c}_t)}
\]
dove l'operatore $\odot$ indica il prodotto elemento per elemento (\emph{Hadamard product}).

\newpage

\section{Riconoscimento Visivo: Oltre La Classificazione Singola}

\subsection{I Diversi Task Della Visione Computazionale}

Finora, lo studio delle reti neurali convoluzionali si è concentrato sulla \textbf{classificazione delle immagini} (\emph{Image Classification}). In questa configurazione, l'immagine d'ingresso viene elaborata da una CNN per produrre un singolo vettore di punteggi associati alle classi di interesse (ad esempio, identificando la presenza di un gatto).

% [IMMAGINE DEL FILE DI SLIDE CON NOME "11-16"] Pipeline classica di classificazione delle immagini tramite CNN — Slide 16 / Pagina 16 del PDF

Tuttavia, una comprensione approfondita di una scena visiva richiede di andare oltre la classificazione globale. Vengono quindi definiti quattro compiti principali di computer vision, caratterizzati da diversi livelli di complessità spaziale e semantica:

% [IMMAGINE DEL FILE DI SLIDE CON NOME "11-17"] Panoramica dei principali task di computer vision: segmentazione semantica, classificazione + localizzazione, object detection, instance segmentation — Slide 17 / Pagina 17 del PDF

\begin{enumerate}
    \item \textbf{Segmentazione Semantica} (\emph{Semantic Segmentation}): consiste nell'assegnare un'etichetta di categoria a ciascun pixel dell'immagine. In questo compito non si distinguono le diverse istanze di una stessa classe (ad esempio, più cani presenti nell'immagine verranno associati a un'unica maschera di pixel della classe "Dog"). L'attenzione è rivolta esclusivamente ai pixel, senza la nozione di singolo oggetto.
    \item \textbf{Classificazione e Localizzazione} (\emph{Classification + Localization}): si concentra su un singolo oggetto preminente nell'immagine. Oltre a classificarlo, l'obiettivo è individuare la sua posizione spaziale racchiudendolo all'interno di una scatola di delimitazione (\textbf{bounding box}).
    \item \textbf{Rilevamento di Oggetti} (\emph{Object Detection}): estende il concetto di localizzazione a scenari con oggetti multipli. La rete deve individuare tutte le istanze degli oggetti appartenenti alle classi di interesse presenti nella scena, posizionando una bounding box attorno a ciascuna di esse e assegnando la relativa etichetta di classe.
    \item \textbf{Segmentazione delle Istanze} (\emph{Instance Segmentation}): rappresenta l'unione tra la segmentazione semantica e l'object detection. L'algoritmo deve non solo classificare i pixel, ma anche distinguere le singole istanze individuali appartenenti alla stessa categoria (ad esempio, colorando con maschere differenti due gatti adiacenti).
\end{enumerate}

\newpage

\section{Segmentazione Semantica}

\subsection{Definizione E Formulazione Del Problema}

La segmentazione semantica mira a etichettare ogni singolo pixel dell'immagine con la rispettiva classe di appartenenza (come cielo, erba, strada, gatto, ecc.).

% [IMMAGINE DEL FILE DI SLIDE CON NOME "11-19"] Esempi di segmentazione semantica per gatti, mucche, erba, alberi e cielo — Slide 19 / Pagina 19 del PDF

Un vincolo fondamentale di questo approccio risiede nel fatto che non vi è alcuna distinzione tra istanze diverse della stessa categoria. Se due mucche pascolano vicine sul prato, la segmentazione semantica produrrà un'unica regione contigua etichettata come "Cow", senza delineare il confine fisico che separa i due animali.

\subsection{Approccio Tramite Finestra Scorrevole}

Una prima intuizione ingenua per risolvere il problema consiste nell'utilizzare un metodo a \textbf{finestra scorrevole} (\emph{sliding window}).

% [IMMAGINE DEL FILE DI SLIDE CON NOME "11-20"] Illustrazione dell'approccio basato su finestra scorrevole per la segmentazione semantica — Slide 20 / Pagina 20 del PDF

Il funzionamento di questa tecnica può essere riassunto nei seguenti passaggi:
\begin{itemize}
    \item Data un'immagine di input, si considera ciascun pixel individualmente.
    \item Si estrae una porzione quadrata locale dell'immagine (\emph{patch}) centrata sul pixel di interesse.
    \item Si somministra tale patch a una rete convoluzionale standard per classificarne il pixel centrale (ad esempio, determinando se appartiene alla classe "mucca" o "erba").
\end{itemize}

Tuttavia, questo approccio presenta un limite insormontabile:
\textbf{Inefficienza computazionale estrema.} Poiché i pixel adiacenti condividono gran parte dell'informazione visiva, le patch estratte per pixel vicini si sovrappongono quasi interamente. Il metodo a finestra scorrevole non riutilizza in alcun modo le caratteristiche condivise tra queste patch sovrapposte, costringendo la rete a ripetere calcoli ridondanti per milioni di volte su una singola immagine.

\subsection{Reti Completamente Convoluzionali A Risoluzione Piena}

Per ovviare all'inefficienza del metodo a finestra scorrevole, si propone l'idea di progettare una rete composta esclusivamente da strati convoluzionali, in grado di generare predizioni per tutti i pixel simultaneamente. Questo tipo di architettura prende il nome di \textbf{Rete Completamente Convoluzionale} (\emph{Fully Convolutional Network}, FCN).

% [IMMAGINE DEL FILE DI SLIDE CON NOME "11-22"] Architettura di una rete completamente convoluzionale (FCN) operante a risoluzione piena — Slide 22 / Pagina 22 del PDF

La struttura di una FCN a risoluzione piena si articola come segue:
\begin{itemize}
    \item \textbf{Input}: un'immagine di dimensioni $3 \times H \times W$, dove $3$ indica i canali colore RGB, $H$ l'altezza e $W$ la larghezza.
    \item \textbf{Strati convoluzionali}: una sequenza di convoluzioni che preservano la risoluzione spaziale, producendo mappe di caratteristiche intermedie di dimensione $D \times H \times W$ (con $D$ numero di canali o filtri).
    \item \textbf{Output}: un tensore di punteggi di dimensione $C \times H \times W$, dove $C$ rappresenta il numero complessivo di classi del dataset.
\end{itemize}

Per ottenere la mappa di predizione finale di dimensioni $H \times W$, si applica l'operatore $\text{argmax}$ lungo la dimensione dei canali di output per ciascuna coordinata spaziale $(i, j)$:
\[
\boxed{\hat{y}_{i, j} = \argmax_{c \in \{1, \dots, C\}} S_{c, i, j}}
\]
dove $S_{c, i, j}$ indica il punteggio relativo alla classe $c$ nel pixel posizionato alla riga $i$ e colonna $j$.

Sebbene questo schema risolva il problema della ridondanza dei calcoli, introduce a sua volta una nuova problematica critica: eseguire operazioni di convoluzione mantenendo la risoluzione spaziale originaria $H \times W$ lungo tutta la profondità della rete risulta estremamente esigente in termini di memoria e tempo di calcolo, rendendo l'approccio impraticabile per immagini ad alta definizione.

\subsection{Reti Completamente Convoluzionali Con Downsampling E Upsampling}

Per coniugare l'efficienza computazionale con la necessità di predizioni ad alta risoluzione spaziale, l'approccio consolidato consiste nel progettare la rete alternando fasi di \textbf{downsampling} (riduzione della risoluzione) e \textbf{upsampling} (aumento della risoluzione) all'interno dell'architettura stessa.

% [IMMAGINE DEL FILE DI SLIDE CON NOME "11-24"] Architettura FCN con percorsi di downsampling e upsampling all'interno della rete — Slide 24 / Pagina 24 del PDF

Questo design prevede un collo di bottiglia (\emph{bottleneck}) centrale e si sviluppa secondo le seguenti fasi:
\begin{enumerate}
    \item \textbf{Fase di Downsampling} (Contrazione): l'immagine di input di dimensione $3 \times H \times W$ viene elaborata tramite convoluzioni e operazioni di pooling per ridurne progressivamente le dimensioni spaziali e aumentarne la profondità del canale. Si passa ad esempio da una risoluzione iniziale a una intermedia ad alta risoluzione (\emph{High-res}: $D_1 \times H/2 \times W/2$), poi a una media risoluzione (\emph{Med-res}: $D_2 \times H/4 \times W/4$), fino a raggiungere una bassa risoluzione al livello del collo di bottiglia (\emph{Low-res}: $D_3 \times H/8 \times W/8$). Questa riduzione consente di applicare i filtri convoluzionali su rappresentazioni spaziali compatte, abbattendo drasticamente il costo computazionale.
    \item \textbf{Fase di Upsampling} (Espansione): l'informazione compressa viene proiettata nuovamente verso l'alto per ripristinare la dimensione spaziale originaria. Attraverso strati di upsampling dedicati, si risale alla media risoluzione ($D_2 \times H/4 \times W/4$), poi alla alta risoluzione ($D_1 \times H/2 \times W/2$), fino a produrre la mappa di predizione finale avente risoluzione $H \times W$.
\end{enumerate}

Questo paradigma architetturale rappresenta la base dei moderni sistemi di segmentazione ed è stato formalizzato in lavori pionieristici della letteratura scientifica, tra cui:
\begin{itemize}
    \item Long, Shelhamer e Darrell, \emph{"Fully Convolutional Networks for Semantic Segmentation"}, CVPR 2015 [24].
    \item Noh et al., \emph{"Learning Deconvolution Network for Semantic Segmentation"}, ICCV 2015 [24].
\end{itemize}

\section{L'Idea Della Segmentazione Semantica: Reti Fully Convolutional}

L'obiettivo principale della segmentazione semantica consiste nel classificare ciascun pixel di un'immagine di input assegnandogli una specifica categoria di appartenenza (ad esempio distinguendo aree corrispondenti a erba, cielo, mucche o alberi). Per affrontare questo compito, l'approccio classico prevede la progettazione di una rete interamente convoluzionale (FCN, \emph{Fully Convolutional Network}) composta da una sequenza di strati convoluzionali in grado di elaborare un input di dimensioni $3 \times H \times W$ e produrre una mappa di predizioni di dimensioni $H \times W$.

% [IMMAGINE DEL FILE DI SLIDE CON NOME "25-48"] Diagramma dell'architettura Fully Convolutional per segmentazione semantica con fasi di downsampling e upsampling — Slide 25 / Pagina 1 del PDF

Per gestire in modo efficiente la complessità computazionale e consentire ai filtri di catturare un contesto spaziale più ampio (aumentando il campo ricettivo o \emph{receptive field}), la struttura di questa rete viene suddivisa in due fasi principali:

\begin{itemize}[noitemsep]
    \item \textbf{Downsampling}: Questa fase riduce progressivamente la risoluzione spaziale delle mappe delle feature mentre ne aumenta la profondità dei canali. Partendo da mappe ad alta risoluzione di dimensioni $D_1 \times H/2 \times W/2$, si passa a risoluzioni intermedie $D_2 \times H/4 \times W/4$ fino ad arrivare a mappe a bassa risoluzione $D_3 \times H/4 \times W/4$ (o inferiori). Le operazioni tipiche per ottenere il downsampling includono strati di \textbf{pooling} o convoluzioni con passo maggiore di uno (\textbf{strided convolution}).
    \item \textbf{Upsampling}: È l'operazione inversa rispetto al downsampling, necessaria per ripristinare la risoluzione spaziale originaria $H \times W$ a partire dalle feature compresse a bassa risoluzione, consentendo così la classificazione finale pixel per pixel.
\end{itemize}

I contributi teorici fondamentali per questa architettura risalgono ai lavori di Long, Shelhamer e Darrell (``Fully Convolutional Networks for Semantic Segmentation'', CVPR 2015) e Noh et al. (``Learning Deconvolution Network for Semantic Segmentation'', ICCV 2015).

\newpage

\subsection{Tecniche Di Upsampling In-Network: Unpooling}

Per effettuare l'upsampling all'interno della rete senza introdurre parametri apprendibili, si ricorre a tecniche geometriche fisse note come \textbf{unpooling}. Consideriamo come esempio un blocco di feature di input di dimensioni $2 \times 2$:

\[
\begin{pmatrix}
1 & 2 \\
3 & 4
\end{pmatrix}
\]

L'obiettivo è espandere spazialmente questa matrice per ottenere un output di dimensioni $4 \times 4$. Esistono due modalità principali per eseguire tale operazione in modo deterministico:

% [IMMAGINE DEL FILE DI SLIDE CON NOME "25-48"] Rappresentazione visiva delle tecniche di unpooling Nearest Neighbor e Bed of Nails — Slide 26 / Pagina 2 del PDF

\begin{enumerate}[noitemsep]
    \item \textbf{Nearest Neighbor (Vicino Più Vicino)}: Ciascun valore della matrice di input viene replicato all'interno di una sotto-regione quadrata di dimensioni pari al fattore di upsampling (in questo caso $2 \times 2$). La matrice di output risultante è:
    \[
    \begin{pmatrix}
    1 & 1 & 2 & 2 \\
    1 & 1 & 2 & 2 \\
    3 & 3 & 4 & 4 \\
    3 & 3 & 4 & 4
    \end{pmatrix}
    \]

    \item \textbf{Bed of Nails (Letto Di Chiodi)}: Il valore di ciascun pixel di input viene copiato in una posizione fissa prestabilita all'interno del corrispondente blocco di output (ad esempio, l'angolo in alto a sinistra), mentre tutti gli altri pixel adiacenti all'interno dello stesso blocco vengono azzerati. La matrice di output risultante assume la forma:
    \[
    \begin{pmatrix}
    1 & 0 & 2 & 0 \\
    0 & 0 & 0 & 0 \\
    3 & 0 & 4 & 0 \\
    0 & 0 & 0 & 0
    \end{pmatrix}
    \]
\end{enumerate}

\newpage

\subsection{Max Unpooling}

Il \textbf{Max Unpooling} rappresenta un'evoluzione asimmetrica e intelligente del concetto di unpooling, progettata per essere utilizzata in coppia con gli strati di \textbf{Max Pooling} posizionati nella parte speculare di downsampling della rete.

% [IMMAGINE DEL FILE DI SLIDE CON NOME "25-48"] Flusso di lavoro del Max Unpooling con trasferimento degli indici di posizione dallo strato di pooling — Slide 27 / Pagina 3 del PDF

Il principio di funzionamento si basa sulla memorizzazione delle coordinate spaziali dei valori massimi durante il cammino di discesa:

\begin{itemize}[noitemsep]
    \item \textbf{Fase di Max Pooling}: Durante il downsampling, quando viene applicato il filtro di Max Pooling su una mappa di feature di input (ad esempio di dimensioni $4 \times 4$), la rete non si limita a estrarre il valore massimo di ogni sotto-regione $2 \times 2$ per produrre l'output $2 \times 2$, ma memorizza attivamente gli indici corrispondenti alla posizione esatta da cui ciascun massimo è stato selezionato.

    Considerando l'esempio numerico fornito nelle slide, si parta dalla matrice di input:
    \[
    \begin{pmatrix}
    1 & 2 & \textcolor{red}{6} & 3 \\
    3 & \textcolor{red}{5} & 2 & 1 \\
    1 & 2 & 2 & 1 \\
    \textcolor{red}{7} & 3 & 4 & \textcolor{red}{8}
    \end{pmatrix}
    \]
    I valori massimi selezionati per ciascun quadrante $2 \times 2$ sono rispettivamente $5$ (in basso a destra nel primo blocco), $6$ (in alto a sinistra nel secondo), $7$ (in basso a sinistra nel terzo) e $8$ (in basso a destra nel quarto). L'output del pooling sarà dunque:
    \[
    \begin{pmatrix}
    5 & 6 \\
    7 & 8
    \end{pmatrix}
    \]
    La rete conserva memoria delle posizioni di questi quattro elementi.

    \item \textbf{Fase di Max Unpooling}: Nel cammino di risalita (upsampling), lo strato di Max Unpooling riceve una mappa di feature a bassa risoluzione (ad esempio $2 \times 2$), generata dalle elaborazioni degli strati intermedi della rete:
    \[
    \begin{pmatrix}
    1 & 2 \\
    3 & 4
    \end{pmatrix}
    \]
    Per ricostruire la mappa $4 \times 4$, i valori dell'input corrente vengono inseriti esattamente nelle medesime posizioni spaziali che erano state memorizzate durante la fase di pooling associata. Tutti gli altri elementi della matrice di output vengono impostati a zero:
    \[
    \begin{pmatrix}
    0 & 0 & \textcolor{red}{2} & 0 \\
    0 & \textcolor{red}{1} & 0 & 0 \\
    0 & 0 & 0 & 0 \\
    \textcolor{red}{3} & 0 & 0 & \textcolor{red}{4}
    \end{pmatrix}
    \]
\end{itemize}

Questo meccanismo consente di preservare e recuperare l'informazione spaziale di dettaglio (come i contorni degli oggetti), che andrebbe altrimenti persa durante le riduzioni dimensionali del pooling, senza la necessità di introdurre parametri addestrabili aggiuntivi.

\newpage

\subsection{Upsampling Apprendibile: Convoluzione Trasposta}

La \textbf{convoluzione trasposta} (\emph{transpose convolution}) definisce un meccanismo di upsampling \textbf{apprendibile}, in cui i pesi del filtro convoluzionale vengono ottimizzati durante la fase di addestramento tramite la discesa del gradiente.

Per comprendere la transizione logica verso la convoluzione trasposta, è utile richiamare il comportamento di una convoluzione standard bidimensionale.

% [IMMAGINE DEL FILE DI SLIDE CON NOME "25-48"] Rappresentazione visiva di una convoluzione standard 2D con stride 1 e stride 2 — Slides 28-33 / Pagine 4-9 del PDF

In una normale convoluzione $3 \times 3$ con stride 1 e padding 1 applicata a un input $4 \times 4$:
\begin{itemize}[noitemsep]
    \item Il filtro si sovrappone a una sotto-regione dell'input.
    \item Viene calcolato il prodotto scalare (dot product) tra i pesi del filtro e i valori dell'input coperti.
    \item Il risultato fornisce un singolo valore scalare nella corrispondente posizione della matrice di output $4 \times 4$.
    \item Lo spostamento della finestra di convoluzione avviene passo dopo passo (stride 1).
\end{itemize}

Se si imposta lo stride pari a 2 (convoluzione standard $3 \times 3$, stride 2, pad 1):
\begin{itemize}[noitemsep]
    \item La finestra si sposta di due pixel sull'input per ogni pixel prodotto in output.
    \item Il rapporto tra il movimento nell'input e nell'output è regolato dallo stride.
    \item Questo processo riduce la dimensione spaziale, mappando l'input $4 \times 4$ in un output $2 \times 2$.
\end{itemize}

\newpage

\subsubsection{Meccanismo Di Calcolo Della Convoluzione Trasposta 2D}

La convoluzione trasposta inverte questo rapporto di scala spaziale. Consideriamo l'applicazione di una convoluzione trasposta $3 \times 3$ con stride 2 e padding 1, che accetta in input una mappa di feature $2 \times 2$ e produce in output una mappa $4 \times 4$.

% [IMMAGINE DEL FILE DI SLIDE CON NOME "25-48"] Schema del calcolo della convoluzione trasposta con stride 2 tramite somma di filtri pesati sovrapposti — Slides 34-38 / Pagine 10-14 del PDF

Il calcolo si sviluppa secondo i seguenti passaggi:
\begin{enumerate}[noitemsep]
    \item Si considera un singolo valore di pixel dall'input (ad esempio, il valore nell'angolo in alto a sinistra).
    \item Questo valore scalare viene utilizzato come coefficiente di moltiplicazione per pesare tutti gli elementi del filtro convoluzionale $3 \times 3$.
    \item La matrice risultante da questa moltiplicazione scalare-matrice viene copiata nella corrispondente area di destinazione dell'output.
    \item Ci si sposta al pixel successivo dell'input. Lo stride pari a 2 indica che la finestra di scrittura sull'output si deve spostare di \textbf{due pixel} per ogni singolo pixel di spostamento nell'input.
    \item Nelle zone della matrice di output in cui i filtri pesati si sovrappongono spazialmente, i contributi sovrapposti vengono sommati tra loro (\emph{sum where output overlaps}).
\end{enumerate}

In letteratura scientifica, questa operazione è identificata anche con altre nomenclature:
\begin{itemize}[noitemsep]
    \item \textbf{Deconvolution} (termine sconsigliato in quanto matematicamente non corrisponde all'inversa della convoluzione);
    \item \textbf{Upconvolution};
    \item \textbf{Fractionally strided convolution} (poiché lo spostamento di un pixel nell'output corrisponde idealmente a uno spostamento frazionario nell'input);
    \item \textbf{Backward strided convolution}.
\end{itemize}

\newpage

\subsubsection{Esempio Unidimensionale Di Convoluzione Trasposta}

Per comprendere matematicamente l'interazione tra i segnali e la generazione delle sovrapposizioni, consideriamo un esempio unidimensionale (1D) con stride 2.

Siano definiti:
\begin{itemize}[noitemsep]
    \item Un vettore di input di dimensione 2: $\bm{a} = \begin{bmatrix} a \\ b \end{bmatrix}$
    \item Un filtro convoluzionale di dimensione 3: $\bm{f} = \begin{bmatrix} x \\ y \\ z \end{bmatrix}$
\end{itemize}

% [IMMAGINE DEL FILE DI SLIDE CON NOME "25-48"] Esempio di calcolo unidimensionale di convoluzione trasposta 1D con stride 2 — Slide 39 / Pagina 15 del PDF

Il primo elemento $a$ viene moltiplicato per il filtro $\bm{f}$, fornendo il contributo:
\[
a \cdot \begin{bmatrix} x \\ y \\ z \end{bmatrix} = \begin{bmatrix} ax \\ ay \\ az \end{bmatrix}
\]

Il secondo elemento $b$ viene parimenti moltiplicato per il filtro $\bm{f}$ e, a causa dello stride pari a 2, il suo contributo viene posizionato nell'output con uno slittamento di due posizioni a destra:
\[
b \cdot \begin{bmatrix} x \\ y \\ z \end{bmatrix} = \begin{bmatrix} bx \\ by \\ bz \end{bmatrix} \quad \text{(slittato di 2)}
\]

Sommando i contributi parziali nelle zone di sovrapposizione (in questo caso, la terza cella riceve sia il termine $az$ che il termine $bx$), si ottiene il vettore finale di output di dimensione 5:
\[
\bm{o} = \begin{bmatrix} ax \\ ay \\ \textcolor{red}{az + bx} \\ by \\ bz \end{bmatrix}
\]

Nel caso in cui si desideri un output che sia esattamente il doppio della dimensione dell'input (ovvero di dimensione 4), si effettua un'operazione di ritaglio (\emph{crop}) rimuovendo l'ultimo pixel in eccesso dal vettore di output.

\newpage

\subsection{Rappresentazione Algebrica: La Convoluzione Come Moltiplicazione Tra Matrici}

Qualsiasi operazione di convoluzione lineare unidimensionale o bidimensionale può essere formulata algebricamente sotto forma di moltiplicazione matrice-vettore:

\[
\bm{x} * \bm{a} = \bm{X}\bm{a}
\]

dove $\bm{a}$ è il vettore contenente i valori di input, opportunamente esteso mediante operazioni di zero-padding, e $\bm{X}$ è una matrice strutturata sparsa (matrice convoluzionale o di Toeplitz) i cui elementi sono definiti dai pesi del filtro di convoluzione.

\subsubsection{Caso Con Stride Pari A 1}

Consideriamo un esempio di convoluzione 1D con filtro di dimensione 3, stride 1 e padding 1. Dato un vettore di input di dimensione 4, lo estendiamo con padding di zeri alle estremità ottenendo $\tilde{\bm{a}} = [0, a, b, c, d, 0]^T$. L'operazione di convoluzione tramite il filtro $[x, y, z]$ si esprime come:

% [IMMAGINE DEL FILE DI SLIDE CON NOME "25-48"] Moltiplicazione matriciale associata a una convoluzione 1D con stride 1 e alla sua trasposta — Slides 40-41 / Pagine 16-17 del PDF

\[
\bm{X}\tilde{\bm{a}} = \begin{bmatrix}
x & y & z & 0 & 0 & 0 \\
0 & x & y & z & 0 & 0 \\
0 & 0 & x & y & z & 0 \\
0 & 0 & 0 & x & y & z
\end{bmatrix}
\begin{bmatrix}
0 \\ a \\ b \\ c \\ d \\ 0
\end{bmatrix} =
\begin{bmatrix}
ay + bz \\
ax + by + cz \\
bx + cy + dz \\
cx + dy
\end{bmatrix}
\]

La convoluzione trasposta corrisponde all'operazione algebrica di moltiplicazione per la matrice trasposta $\bm{X}^T$:

\[
\bm{x} *^T \bm{a} = \bm{X}^T\bm{a}
\]

Moltiplicando l'input di dimensione 4 per la trasposta $\bm{X}^T$ di dimensioni $6 \times 4$, otteniamo un output di dimensione 6:

\[
\bm{X}^T\bm{a} = \begin{bmatrix}
x & 0 & 0 & 0 \\
y & x & 0 & 0 \\
z & y & x & 0 \\
0 & z & y & x \\
0 & 0 & z & y \\
0 & 0 & 0 & z
\end{bmatrix}
\begin{bmatrix}
a \\ b \\ c \\ d
\end{bmatrix} =
\begin{bmatrix}
ax \\
ay + bx \\
az + by + cx \\
bz + cy + dx \\
cz + dy \\
dz
\end{bmatrix}
\]

Quando lo stride è pari a 1, la convoluzione trasposta si comporta esattamente come una convoluzione standard, differendo esclusivamente per le regole geometriche di padding applicate alle estremità.

\newpage

\subsubsection{Caso Con Stride Maggiore Di 1}

Nel caso in cui si utilizzi uno stride superiore a 1 (ad esempio stride pari a 2), la matrice convoluzionale $\bm{X}$ si ottiene selezionando una riga ogni due rispetto alla corrispondente matrice con stride 1. Per un filtro di dimensione 3, stride 2 e padding 1, la moltiplicazione assume la forma:

% [IMMAGINE DEL FILE DI SLIDE CON NOME "25-48"] Formulazione matriciale della convoluzione standard e della convoluzione trasposta con stride 2 — Slides 42-43 / Pagine 18-19 del PDF

\[
\bm{X}\tilde{\bm{a}} = \begin{bmatrix}
x & y & z & 0 & 0 & 0 \\
0 & 0 & x & y & z & 0
\end{bmatrix}
\begin{bmatrix}
0 \\ a \\ b \\ c \\ d \\ 0
\end{bmatrix} =
\begin{bmatrix}
ay + bz \\
bx + cy + dz
\end{bmatrix}
\]

La corrispondente operazione di convoluzione trasposta si esegue moltiplicando un vettore di input $[a, b]^T$ per la matrice trasposta $\bm{X}^T$ di dimensioni $6 \times 2$:

\[
\bm{X}^T\bm{a} = \begin{bmatrix}
x & 0 \\
y & 0 \\
z & x \\
0 & y \\
0 & z \\
0 & 0
\end{bmatrix}
\begin{bmatrix}
a \\ b
\end{bmatrix} =
\begin{bmatrix}
ax \\
ay \\
az + bx \\
by \\
bz \\
0
\end{bmatrix}
\]

Risulta evidente che, quando lo $\text{stride} > 1$, la convoluzione trasposta non può più essere ricondotta a una convoluzione standard. Essa definisce un'operazione intrinsecamente differente, caratterizzata dalla proiezione distanziata delle componenti e dalla somma dei relativi termini nelle zone di sovrapposizione.

\newpage

\subsection{L'Architettura Fully Convolutional Completa}

Integrando i meccanismi di riduzione e di espansione dimensionale esaminati, l'architettura complessiva di una rete Fully Convolutional per la segmentazione semantica può essere così sintetizzata:

\begin{itemize}[noitemsep]
    \item Un cammino di discesa (\textbf{downsampling}) che sfrutta strati di pooling o convoluzioni con stride per estrarre caratteristiche semantiche ad alta astrazione e ridurre la dimensionalità spaziale.
    \item Un cammino simmetrico di risalita (\textbf{upsampling}) basato su moduli di \textbf{unpooling} o \textbf{convoluzioni trasposte con stride} per ripristinare progressivamente la risoluzione geometrica originaria delle mappe di feature, fino a ottenere la mappa di predizione finale di dimensioni $H \times W$.
\end{itemize}

% [IMMAGINE DEL FILE DI SLIDE CON NOME "25-48"] Schema completo di una FCN che integra downsampling e upsampling per la segmentazione semantica — Slide 44 / Pagina 20 del PDF

\newpage

\section{Classificazione E Localizzazione}

Nel campo del riconoscimento visivo tramite deep learning, si distinguono diverse tipologie di compiti a seconda del livello di granularità spaziale e del numero di istanze da individuare:

% [IMMAGINE DEL FILE DI SLIDE CON NOME "25-48"] Schema comparativo tra classificazione, localizzazione, object detection e segmentazione semantica/delle istanze — Slide 45 / Pagina 21 del PDF

\begin{itemize}[noitemsep]
    \item \textbf{Classificazione (Classification)}: Identificazione della categoria semantica principale presente all'interno dell'immagine (ad esempio, determinare se nell'immagine è presente un gatto).
    \item \textbf{Classificazione e Localizzazione (Classification + Localization)}: Tratta lo scenario in cui è presente un singolo oggetto di interesse (\emph{single object}). L'obiettivo è sia determinarne la classe di appartenenza sia individuarne la posizione spaziale, definendo i confini di una scatola di delimitazione o \textbf{bounding box}.
    \item \textbf{Rilevamento di Oggetti (Object Detection)}: Estende la localizzazione al caso in cui vi siano molteplici oggetti appartenenti a classi differenti all'interno della stessa scena (\emph{multiple objects}), richiedendo l'identificazione e la delimitazione di ogni singola istanza.
    \item \textbf{Segmentazione delle Istanze (Instance Segmentation)}: Unisce i concetti di rilevamento degli oggetti e segmentazione semantica, individuando le singole istanze presenti e definendo per ciascuna di esse una maschera di pixel precisa, anziché limitarsi a una bounding box rettangolare.
\end{itemize}

\newpage

\subsection{Localizzazione Come Problema Di Regressione E Loss Multitask}

Per implementare congiuntamente la classificazione e la localizzazione di un singolo oggetto all'interno di una rete neurale, si formula la localizzazione come un **problema di regressione**.

L'architettura della rete si organizza secondo la seguente pipeline:

\begin{enumerate}[noitemsep]
    \item L'immagine di input viene elaborata da una rete convoluzionale (\textbf{ConvNet}) deputata all'estrazione delle feature.
    \item Gli strati convoluzionali terminano in un vettore di caratteristiche completamente connesso (fully connected), tipicamente di dimensione 4096.
    \item A partire da questo vettore comune, la rete si sdoppia in due rami indipendenti (\emph{teste}):
    \begin{itemize}[noitemsep]
        \item \textbf{Testa di Classificazione}: Composta da uno strato completamente connesso che mappa il vettore da 4096 elementi alle $C$ classi di output (ad esempio, uno strato FC da 4096 a 1000). Questo ramo produce i punteggi delle classi (\emph{class scores}), che definiscono la probabilità di appartenenza dell'oggetto a ciascuna categoria (ad esempio: Gatto: 0.9, Cane: 0.05, Auto: 0.01).
        \item \textbf{Testa di Localizzazione (Regressione)}: Composta da uno strato completamente connesso che mappa il vettore da 4096 elementi a un vettore di 4 dimensioni reali (strato FC da 4096 a 4). Questo vettore rappresenta le coordinate della bounding box predetta:
        \[
        \bm{b} = \begin{bmatrix} x \\ y \\ w \\ h \end{bmatrix}
        \]
        dove $(x, y)$ sono le coordinate spaziali del centro della scatola di delimitazione, mentre $w$ e $h$ ne indicano rispettivamente la larghezza e l'altezza.
    \end{itemize}
\end{enumerate}

% [IMMAGINE DEL FILE DI SLIDE CON NOME "25-48"] Architettura di rete a due rami per compiti di classificazione (Softmax Loss) e localizzazione (L2 Loss) — Slides 46-48 / Pagine 22-24 del PDF

Durante la fase di addestramento, la rete viene ottimizzata definendo una **Loss Multitask (funzione di costo multitask)**, formata dalla combinazione lineare delle perdite associate a ciascun ramo:

\begin{itemize}[noitemsep]
    \item Per la classificazione, si applica la **Softmax Loss** (cross-entropia) tra i punteggi di classe predetti e l'etichetta corretta dell'immagine (ground-truth).
    \item Per la localizzazione, si applica la **L2 Loss** (errore quadratico medio) tra la bounding box predetta $\bm{b} = [x, y, w, h]^T$ e la bounding box reale fornita dal dataset $\bm{b}' = [x', y', w', h']^T$:
    \[
    \mathcal{L}_{reg} = \|\bm{b} - \bm{b}'\|_2^2 = (x - x')^2 + (y - y')^2 + (w - w')^2 + (h - h')^2
    \]
\end{itemize}

La loss totale che guida l'aggiornamento dei parametri della rete tramite backpropagation è data da:

\[
\mathcal{L}_{totale} = \mathcal{L}_{class} + \lambda \mathcal{L}_{reg}
\]

dove $\lambda$ rappresenta un iperparametro di regolarizzazione adoperato per bilanciare l'importanza relativa dei due compiti durante la fase di ottimizzazione.

\section{Pratica Della Classificazione E Localizzazione: Transfer Learning}

Per implementare concretamente un sistema di classificazione e localizzazione, è prassi comune non addestrare l'intera rete convoluzionale partendo da zero. L'addestramento da zero di modelli complessi richiede enormi quantità di dati e tempo di calcolo. Di conseguenza, si ricorre quasi sempre al \textbf{transfer learning} (apprendimento per trasferimento): la rete convoluzionale che funge da estrattore di caratteristiche (\emph{backbone}) viene inizializzata utilizzando pesi pre-addestrati su un dataset su larga scala, tipicamente ImageNet.

% [IMMAGINE DEL FILE DI SLIDE CON NOME "49-72"] Architettura di rete per classificazione e localizzazione con evidenza sul pre-addestramento della ConvNet tramite Transfer Learning su ImageNet — Slide 49 / Pagina 1 del PDF

Questo approccio consente di sfruttare i filtri convoluzionali già specializzati nel rilevamento di pattern visivi generici, come angoli, texture e forme complesse, velocizzando drasticamente la convergenza e migliorando l'accuratezza finale anche in presenza di un numero limitato di campioni per il compito di destinazione.

\subsection{Una Digressione: Stima Della Posa Umana}

Un'applicazione strettamente legata ai concetti di localizzazione geometrica e regressione è la \textbf{stima della posa umana} (\emph{human pose estimation}). L'obiettivo di questo compito consiste nell'individuare la configurazione spaziale del corpo di un individuo all'interno di un'immagine.

% [IMMAGINE DEL FILE DI SLIDE CON NOME "49-72"] Esempio di stima della posa umana con i 14 giunti articolari proiettati su un tennista — Slide 50 / Pagina 2 del PDF

Per formalizzare questo problema in modo geometrico, la posa del corpo umano viene comunemente descritta attraverso un insieme rigido di $14$ punti di snodo articolare (\emph{joint positions}), come proposto nel lavoro fondamentale di Johnson e Everingham, ``Clustered Pose and Nonlinear Appearance Models for Human Pose Estimation'' (BMVC 2010):
\begin{itemize}[noitemsep]
    \item Piede sinistro e destro (\emph{left / right foot})
    \item Ginocchio sinistro e destro (\emph{left / right knee})
    \item Anca sinistra e destra (\emph{left / right hip})
    \item Spalla sinistra e destra (\emph{left / right shoulder})
    \item Gomito sinistro e destro (\emph{left / right elbow})
    \item Mano sinistra e destra (\emph{left / right hand})
    \item Collo (\emph{neck})
    \item Sommità della testa (\emph{head top})
\end{itemize}

\newpage

\subsubsection{Formulazione Della Loss Per La Stima Della Posa}

Un'architettura celebre per risolvere questo problema è \textbf{DeepPose}, introdotta da Toshev e Szegedy in ``DeepPose: Human Pose Estimation via Deep Neural Networks'' (CVPR 2014). Il problema viene formulato interamente come un compito di regressione spaziale multi-variata:

% [IMMAGINE DEL FILE DI SLIDE CON NOME "49-72"] Pipeline di DeepPose che mostra la regressione delle coordinate dei 14 giunti corporei a partire da un vettore di feature di dimensione 4096 — Slides 51-52 / Pagine 3-4 del PDF

\begin{enumerate}[noitemsep]
    \item L'immagine di input contenente la persona viene elaborata da una rete convoluzionale profondamente addestrata.
    \item L'ultimo strato completamente connesso estrae un vettore rappresentativo di dimensione $4096$.
    \item Questo vettore viene mappato in un output continuo contenente le coordinate spaziali $(x, y)$ per ciascuno dei $14$ punti di snodo articolare della posa (piede sinistro, piede destro, fino alla sommità della testa).
    \item Durante l'addestramento, ciascun punto predetto viene confrontato con la corrispondente coordinata reale (\emph{ground truth}) fornita dal dataset. Si applica una funzione di costo quadratica (\textbf{L2 loss}) su ciascun giunto. La funzione di costo complessiva $\mathcal{L}_{pose}$ è data dalla somma delle perdite L2 calcolate per ciascuno dei $14$ punti articolari:
    \[
    \boxed{\mathcal{L}_{pose} = \sum_{i=1}^{14} \left( (x_i - x'_i)^2 + (y_i - y'_i)^2 \right)}
    \]
    dove $(x_i, y_i)$ rappresenta la coordinata predetta per l'articolazione $i$-esima, mentre $(x'_i, y'_i)$ indica la coordinata reale corretta di riferimento.
\end{enumerate}

\newpage

\section{Il Rilevamento Di Oggetti: Object Detection}

Come introdotto nella tassonomia generale, il \textbf{rilevamento di oggetti} (\emph{object detection}) si differenzia nettamente dalla classificazione e localizzazione a causa della complessità indotta dalla presenza di un numero variabile e non noto a priori di istanze all'interno della medesima immagine (\emph{multiple object}).

% [IMMAGINE DEL FILE DI SLIDE CON NOME "49-72"] Confronto tra compiti di segmentazione semantica, classificazione con localizzazione, object detection e segmentazione delle istanze — Slide 53 / Pagina 5 del PDF

Mentre la localizzazione assume l'esistenza di un singolo oggetto target di cui predire la bounding box, l'object detection deve localizzare e classificare simultaneamente ogni singola istanza presente nella scena.

\subsection{L'Impatto Del Deep Learning Nell'Object Detection}

Le prestazioni degli algoritmi di object detection sono state storicamente valutate su benchmark pubblici competitivi, in particolare sulla sfida \textbf{PASCAL VOC}. L'introduzione delle reti convoluzionali profonde ha segnato un punto di svolta radicale per questa metrica, misurata in termini di \textbf{mean Average Precision} (mAP):

% [IMMAGINE DEL FILE DI SLIDE CON NOME "49-72"] Andamento storico della metrica mAP sul dataset PASCAL VOC dal 2007 al 2015, evidenziando il salto prestazionale dovuto alle reti convoluzionali — Slide 54 / Pagina 6 del PDF

\begin{itemize}[noitemsep]
    \item \textbf{Prima dei deep convnets (2007-2012)}: I metodi tradizionali, basati su estrattori di caratteristiche progettati manualmente (\emph{hand-crafted features}), mostravano un progresso estremamente lento e incrementale, con valori di mAP che faticavano a salire da circa il $17\%$ del 2007 fino a stabilizzarsi intorno al $40\%$ nel 2012.
    \item \textbf{Utilizzando i deep convnets (2013-2015)}: L'adozione di architetture convoluzionali profonde ha innescato una crescita rapidissima delle prestazioni, portando il mAP a balzare immediatamente oltre il $53\%$ nel 2013, per poi superare il $62\%$ nel 2014 e raggiungere valori superiori al $75\%$ nel 2015.
\end{itemize}

Questo andamento, documentato nei dati di Ross Girshick (2015), dimostra l'efficacia superiore delle rappresentazioni gerarchiche apprese automaticamente rispetto ai vecchi approcci geometrici o statistici tradizionali.

\newpage

\subsection{Perché L'Object Detection Non Può Essere Formulato Come Semplice Regressione?}

Di fronte alla necessità di localizzare più oggetti, si potrebbe ipotizzare di applicare lo stesso schema di regressione diretta utilizzato per la singola localizzazione o per la stima della posa. Tuttavia, questo approccio si scontra con un vincolo architetturale invalicabile delle reti neurali tradizionali:

% [IMMAGINE DEL FILE DI SLIDE CON NOME "49-72"] Esempio di immagini con numero variabile di oggetti (un gatto, tre animali, molte anatre) e le relative dimensioni degli output — Slides 55-56 / Pagine 7-8 del PDF

\begin{itemize}[noitemsep]
    \item Un'immagine contenente un solo gatto richiede la predizione di un'unica bounding box, corrispondente a un vettore di $4$ numeri: $\text{CAT}: (x, y, w, h)$.
    \item Un'immagine con tre animali richiede la definizione di tre bounding box separate, per un totale di $12$ coordinate (o $16$ numeri se si considerano altre informazioni).
    \item Un'immagine contenente uno stormo composto da decine di anatre richiede la predizione simultanea di moltissime coordinate spaziali (\emph{many numbers}).
\end{itemize}

Di conseguenza, \textbf{ogni immagine richiede un numero differente di output}. Poiché le reti neurali standard sono vincolate ad avere strati di output completamente connessi con dimensioni rigidamente fisse e predeterminate in fase di progettazione, non è possibile addestrare una rete a effettuare una regressione diretta di un numero indefinito e variabile di bounding box.

\newpage

\subsection{Rilevamento Come Classificazione: Sliding Window}

Per aggirare il problema della dimensione variabile dell'output, si può tentare di convertire il rilevamento in un compito di classificazione standard applicando l'approccio della \textbf{finestra scorrevole} (\emph{sliding window}).

L'idea alla base consiste nell'estrarre sistematicamente una moltitudine di ritagli (\emph{crops}) dall'immagine di input a diverse scale, posizioni e rapporti d'aspetto, per poi sottoporre ciascun ritaglio a una rete convoluzionale (\emph{CNN}) addestrata a discriminare se l'area in esame contenga un oggetto specifico o semplicemente lo sfondo:

% [IMMAGINE DEL FILE DI SLIDE CON NOME "49-72"] Esempio di finestra scorrevole che analizza diverse regioni dell'immagine classificandole come sfondo o come oggetti specifici — Slides 57-60 / Pagine 9-12 del PDF

\begin{itemize}[noitemsep]
    \item Se la finestra scorrevole analizza un'area marginale priva di soggetti significativi (ad esempio la parte sinistra dell'immagine), la CNN produce come output: Dog? NO, Cat? NO, Background? YES.
    \item Spostando la finestra sul Golden Retriever, la rete risponderà: Dog? YES, Cat? NO, Background? NO.
    \item Spostando ulteriormente la finestra sul cane più piccolo al centro, si otterrà nuovamente: Dog? YES, Cat? NO, Background? NO.
    \item Posizionando infine la finestra sul gatto a destra, l'output sarà: Dog? NO, Cat? YES, Background? NO.
\end{itemize}

\textbf{Il collo di bottiglia computazionale}: Nonostante la semplicità concettuale, questo approccio è affetto da una drammatica inefficienza. Per garantire la localizzazione accurata di oggetti di ogni dimensione e in qualsiasi coordinata spaziale, è necessario valutare un \textbf{numero enorme di posizioni e scale geometriche}. Eseguire l'intera rete convoluzionale per ciascuna di queste finestre risulta computazionalmente insostenibile ed estremamente costoso.

% [IMMAGINE DEL FILE DI SLIDE CON NOME "49-72"] Evidenziazione del limite computazionale della sliding window dovuto all'elevato numero di scale e posizioni da valutare — Slide 61 / Pagina 13 del PDF

\newpage

\subsection{Proposte Di Regioni: Region Proposals}

Per risolvere l'inefficienza insita nella sliding window, sono stati sviluppati algoritmi di \textbf{proposta di regioni} (\emph{region proposals}). Questi metodi hanno lo scopo di esaminare rapidamente l'immagine per individuare aree ``di tipo oggetto'' (\emph{blobby regions}) che presentano un'elevata probabilità di contenere un elemento di interesse, escludendo a priori le zone chiaramente appartenenti allo sfondo.

% [IMMAGINE DEL FILE DI SLIDE CON NOME "49-72"] Esempio di proposte di regioni "blobby" generate su un'immagine contenente un gattino e un cucciolo di cane — Slide 62 / Pagina 14 del PDF

Questi algoritmi sono progettati per essere computazionalmente leggeri e veloci da eseguire su CPU. Ad esempio, l'algoritmo di **Selective Search** (Uijlings et al., ``Selective Search for Object Recognition'', IJCV 2013) è in grado di produrre circa $1000$ o $2000$ proposte di regioni candidate in pochi secondi.

Altri contributi rilevanti in quest'ambito includono:
\begin{itemize}[noitemsep]
    \item Alexe et al., ``Measuring the objectness of image windows'' (TPAMI 2012);
    \item Cheng et al., ``BING: Binarized normed gradients for objectness estimation at 300fps'' (CVPR 2014);
    \item Zitnick e Dollar, ``Edge boxes: Locating object proposals from edges'' (ECCV 2014).
\end{itemize}

\newpage

\section{L'Architettura R-CNN}

Sfruttando il concetto di proposta di regione, Girshick et al. hanno introdotto l'architettura **R-CNN** (\emph{Region-based Convolutional Neural Network}) nel lavoro ``Rich feature hierarchies for accurate object detection and semantic segmentation'' (CVPR 2014). Questo modello combina per la prima volta l'efficacia dei metodi di proposta di regioni con la capacità rappresentativa delle reti convoluzionali profonde.

% [IMMAGINE DEL FILE DI SLIDE CON NOME "49-72"] Pipeline completa del modello R-CNN dal caricamento dell'immagine all'output di classificazione e regressione — Slides 63-68 / Pagine 15-20 del PDF

La pipeline di elaborazione di R-CNN si sviluppa in modo sequenziale attraverso i seguenti passaggi fondamentali:
\begin{enumerate}[noitemsep]
    \item \textbf{Input}: Si acquisisce l'immagine originale.
    \item \textbf{Estrazione delle Region of Interest (RoI)}: Si applica un algoritmo di proposta di regioni esterno (come Selective Search) per individuare circa $2000$ regioni candidate che potrebbero ospitare un oggetto.
    \item \textbf{Adattamento Dimensionale (Warping)}: Poiché le reti convoluzionali con strati completamente connessi richiedono ingressi di dimensioni spaziali fisse e uniformi, ciascuna delle $2000$ regioni candidate (aventi originariamente rapporti d'aspetto variabili) viene ritagliata e deformata geometricamente (\emph{warped}) per essere ricondotta a una dimensione fissa predefinita.
    \item \textbf{Inoltro alla ConvNet}: Ognuna delle regioni deformate viene inoltrata singolarmente attraverso una ConvNet profonda per estrarre un vettore di caratteristiche ad alta dimensionalità.
    \item \textbf{Classificazione con SVM}: I vettori di caratteristiche estratti vengono classificati utilizzando un insieme di macchine a vettori di supporto lineari (\textbf{linear SVMs}) pre-addestrate per ciascuna classe target.
    \item \textbf{Regressione della Bounding Box (Bbox reg)}: Contemporaneamente alla classificazione, le caratteristiche estratte vengono fornite a un regressore lineare che stima i parametri di scostamento (\emph{offsets}) per affinare geometricamente la precisione della bounding box proposta inizialmente.
\end{enumerate}

\newpage

\subsection{Limitazioni E Problemi Di R-CNN}

Nonostante l'eccellente accuratezza, l'architettura R-CNN originale presenta pesanti limiti operativi che ne precludono l'impiego in tempo reale:

% [IMMAGINE DEL FILE DI SLIDE CON NOME "49-72"] Schema riassuntivo dei colli di bottiglia computazionali e del flusso di addestramento frammentato di R-CNN — Slide 69 / Pagina 21 del PDF

\begin{itemize}[noitemsep]
    \item \textbf{Obiettivi di addestramento frammentati (\emph{Ad hoc training objectives})}: Il sistema non è addestrabile in modalità \emph{end-to-end}. Il processo richiede infatti tre fasi distinte e isolate:
    \begin{enumerate}[noitemsep]
        \item Il fine-tuning della ConvNet utilizzando un classificatore softmax (minimizzando la log loss).
        \item L'addestramento post-hoc delle SVM lineari (minimizzando la hinge loss).
        \item L'addestramento post-hoc dei regressori per le bounding box tramite il metodo dei minimi quadrati (\emph{least squares}).
    \end{enumerate}
    \item \textbf{Processo di training lento e oneroso}: L'addestramento completo richiede circa $84$ ore di calcolo e necessita di un'enorme quantità di memoria su disco fisso, poiché occorre salvare preventivamente i vettori di feature estratti dalla ConvNet per ciascuna delle $2000$ regioni di ogni singola immagine prima di poter addestrare SVM e regressori.
    \item \textbf{Infezione ed esecuzione lenta}: Il tempo di inferenza è estremamente elevato, richiedendo circa $47$ secondi per singola immagine quando si adotta un backbone pesante come la rete VGG16 (Simonyan \& Zisserman, ICLR15). Questo limite invalicabile è dovuto alla necessità di eseguire la ConvNet $2000$ volte in modo indipendente per ogni immagine, una volta per ciascuna RoI. Questo collo di bottiglia è stato parzialmente affrontato in seguito dal modello SPP-net (He et al., ECCV14).
\end{itemize}

\newpage

\section{L'Architettura Fast R-CNN}

Per superare i limiti di lentezza di R-CNN, Ross Girshick ha proposto un'evoluzione denominata **Fast R-CNN** (``Fast R-CNN'', ICCV 2015). L'idea chiave consiste nel \textbf{condividere il calcolo convoluzionale} sull'intera immagine, evitando di elaborare le singole regioni in modo isolato:

% [IMMAGINE DEL FILE DI SLIDE CON NOME "49-72"] Pipeline di Fast R-CNN che illustra l'estrazione della mappa di feature conv5 globale e la successiva proiezione delle RoI — Slides 70-72 / Pagine 22-24 del PDF

\begin{enumerate}[noitemsep]
    \item \textbf{Elaborazione Unificata dell'Immagine}: L'intera immagine di input viene inoltrata una sola volta attraverso la rete convoluzionale principale per generare una mappa di caratteristiche globale ad alta risoluzione semantica (denominata mappa delle feature ``\textbf{conv5}'' dell'immagine).
    \item \textbf{Proiezione delle Region of Interest (RoI)}: Invece di ritagliare ed effettuare il warping delle regioni direttamente sull'immagine originale, le proposte di regioni (RoI) generate dall'algoritmo esterno vengono proiettate direttamente sulla mappa delle feature conv5.
\end{enumerate}

Grazie a questo accorgimento, l'estrazione delle caratteristiche convoluzionali avviene in un unico passaggio per l'intera immagine, eliminando il costo computazionale dovuto alla ripetizione del calcolo convoluzionale per ognuna delle $2000$ regioni candidate.

\section{L'Evoluzione Di Fast R-CNN E Il Modulo RoI Pooling}

Per comprendere l'evoluzione dei moderni sistemi di rilevamento degli oggetti (\emph{object detection}), è fondamentale analizzare i limiti dei primi approcci, come R-CNN, e come l'introduzione di \textbf{Fast R-CNN} (proposto da Ross Girshick nel 2015) abbia ridefinito l'efficienza di questa classe di modelli.

Nell'architettura Fast R-CNN, invece di sottoporre ciascuna regione proposta singolarmente alla rete convoluzionale, l'intera immagine di input viene elaborata in un unico passaggio (\emph{forward pass}) attraverso una rete convoluzionale profonda (\textbf{ConvNet}). Questo consente di estrarre una mappa delle caratteristiche globale ad alto livello, tipicamente indicata come mappa conv5 (\textbf{"conv5" feature map}).

Nel frattempo, le regioni candidate o regioni di interesse (\textbf{Regions of Interest, RoI}) vengono generate a partire da un metodo di proposta esterno. Queste RoI, espresse originariamente nelle coordinate dell'immagine di input, vengono proiettate direttamente sulla mappa conv5.

Per convertire queste aree proiettate, che presentano dimensioni geometriche variabili a seconda della proposta originaria, in vettori di caratteristiche a dimensione fissa, viene introdotto uno strato specifico chiamato \textbf{"RoI Pooling" layer}.

% [IMMAGINE DEL FILE DI SLIDE CON NOME "73-95"] Schema iniziale dell'architettura Fast R-CNN con evidenziato l'inserimento del modulo "RoI Pooling" — Slide 1 / Pagina 1 del PDF

Una volta che il livello di RoI Pooling ha estratto vettori di dimensioni costanti, questi vengono inviati a una serie di livelli completamente connessi (\textbf{Fully-Connected layers, FCs}).

% [IMMAGINE DEL FILE DI SLIDE CON NOME "73-95"] Schema di Fast R-CNN con l'aggiunta dello strato Fully-Connected e del classificatore Softmax — Slide 2 / Pagina 2 del PDF

All'uscita dei livelli FCs, la rete si sdoppia in due rami predittivi paralleli, realizzando un'architettura multi-task:
\begin{enumerate}
    \item Un ramo lineare seguito da un classificatore \textbf{Softmax} (\emph{Linear + softmax}), incaricato di stimare la distribuzione di probabilità dell'oggetto sulle classi di destinazione (incluso il background).
    \item Un ramo di regressione lineare dei bounding-box (\textbf{Linear Bounding-box regressors}), che calcola gli scostamenti (\emph{offset}) geometrici per rifinire la posizione e la dimensione della scatola di delimitazione di ciascuna classe.
\end{enumerate}

% [IMMAGINE DEL FILE DI SLIDE CON NOME "73-95"] Schema completo di Fast R-CNN con l'aggiunta dei regressori lineari per i bounding-box — Slide 3 / Pagina 3 del PDF

\newpage

\subsection{Addestramento Della Rete E Propagazione All'Indietro}

L'aspetto chiave di Fast R-CNN risiede nella possibilità di essere addestrato congiuntamente (\emph{end-to-end}) grazie a una \textbf{funzione di loss multi-task} (\emph{multi-task loss}), espressa come:

\[
\mathcal{L} = \mathcal{L}_{\text{class}} + \mathcal{L}_{\text{loc}}
\]

Nello specifico, per ciascuna regione di interesse, il modello ottimizza contemporaneamente:
\begin{itemize}
    \item La perdita di classificazione, definita tramite una \textbf{Log loss} applicata sulla distribuzione Softmax.
    \item La perdita di localizzazione, implementata mediante una \textbf{Smooth L1 loss} che misura lo scostamento tra i parametri del bounding-box predetti e quelli reali (\emph{ground truth}).
\end{itemize}

% [IMMAGINE DEL FILE DI SLIDE CON NOME "73-95"] Schema di addestramento di Fast R-CNN con l'introduzione della loss multi-task — Slide 4 / Pagina 4 del PDF

Durante la fase di addestramento, l'errore calcolato tramite la loss multi-task viene propagato all'indietro (\textbf{backpropagation}) attraverso l'intera architettura. Il gradiente fluisce dai rami di classificazione e regressione, attraversa i livelli completamente connessi (FCs), passa per il livello di RoI Pooling e giunge fino ai livelli convoluzionali della ConvNet iniziale. Questo flusso consente l'aggiornamento simultaneo di tutti i parametri del modello.

% [IMMAGINE DEL FILE DI SLIDE CON NOME "73-95"] Schema di backpropagation di Fast R-CNN evidenziato dalle frecce viola — Slide 5 / Pagina 5 del PDF

\newpage

\subsection{Dettaglio Funzionale Del Livello RoI Pooling}

Per comprendere la necessità matematica del livello RoI Pooling, si consideri la gestione delle dimensioni dei tensori all'interno della rete, illustrata nell'esempio di Slide 6.

Si assuma di disporre di un'immagine di input ad alta risoluzione con dimensioni $3 \times 640 \times 480$, sulla quale è definita una regione di interesse. L'immagine viene inoltrata attraverso la ConvNet, producendo una mappa di caratteristiche convoluzionali ad alta risoluzione (\emph{Hi-res conv features}) di dimensioni:

\[
512 \times 20 \times 15
\]

La regione proposta sull'immagine originale viene proiettata geometricamente su questa mappa. Poiché le proposte originarie variano in dimensione, la corrispondente proposta di regione proiettata avrà dimensioni variabili; ad esempio, potrebbe corrispondere a un volume di:

\[
512 \times 18 \times 8
\]

I livelli completamente connessi della rete si aspettano in ingresso una dimensione fissa e predefinita, tipicamente a bassa risoluzione, pari a:

\[
512 \times 7 \times 7
\]

Il livello di \textbf{RoI Pooling} risolve questa incompatibilità dimensionale operando nel modo seguente:
\begin{enumerate}
    \item Divide la regione proiettata (nell'esempio di dimensione variabile $18 \times 8$) in una griglia spaziale fissa di dimensione $7 \times 7$.
    \item Applica un'operazione di \textbf{max-pooling} in maniera indipendente all'interno di ciascuna cella della griglia e su ciascuno dei $512$ canali.
\end{enumerate}

Il risultato di questa operazione è un tensore normalizzato di dimensioni esatte $512 \times 7 \times 7$, pronto per essere elaborato dai successivi strati completamente connessi.

% [IMMAGINE DEL FILE DI SLIDE CON NOME "73-95"] Schema esplicativo del funzionamento del livello di RoI Pooling con le dimensioni dei tensori — Slide 6 / Pagina 6 del PDF

\newpage

\section{Analisi Delle Prestazioni E Limitazioni Di Fast R-CNN}

Un confronto quantitativo tra le architetture R-CNN, SPP-Net e Fast R-CNN evidenzia i notevoli miglioramenti introdotti in termini di tempi di addestramento e di esecuzione a test (\emph{test-time}).

I dati relativi al **tempo di addestramento** (misurato in ore) mostrano la seguente progressione:
\begin{itemize}
    \item \textbf{R-CNN}: $84$ ore.
    \item \textbf{SPP-Net}: $25.5$ ore.
    \item \textbf{Fast R-CNN}: $8.75$ ore.
\end{itemize}

I dati relativi al **tempo di test per singola immagine** (espresso in secondi) evidenziano un netto divario a seconda che si includa o si escluda la generazione delle proposte di regione:
\begin{itemize}
    \item \textbf{R-CNN}: $49\,\text{s}$ (incluso il calcolo delle proposte) / $47\,\text{s}$ (escluso il calcolo delle proposte).
    \item \textbf{SPP-Net}: $4.3\,\text{s}$ (incluso il calcolo delle proposte) / $2.3\,\text{s}$ (escluso il calcolo delle proposte).
    \item \textbf{Fast R-CNN}: $2.3\,\text{s}$ (incluso il calcolo delle proposte) / $0.32\,\text{s}$ (escluso il calcolo delle proposte).
\end{itemize}

% [IMMAGINE DEL FILE DI SLIDE CON NOME "73-95"] Grafici comparativi dei tempi di addestramento e di test per R-CNN, SPP-Net e Fast R-CNN — Slide 7 / Pagina 7 del PDF

Sebbene Fast R-CNN riduca il tempo di calcolo della rete a soli $0.32$ secondi per immagine, sorge un problema evidente: il tempo di esecuzione complessivo ($2.3$ secondi) è interamente dominato dal calcolo algoritmico delle proposte di regione esterne (\emph{region proposals}). Questo costituisce il collo di bottiglia computazionale del sistema.

% [IMMAGINE DEL FILE DI SLIDE CON NOME "73-95"] Grafici comparativi dei tempi di test con evidenziata la problematica del collo di bottiglia delle proposte di regione — Slide 8 / Pagina 8 del PDF

\newpage

\section{Faster R-CNN: Integrazione Della Region Proposal Network}

Per superare la limitazione imposta dal calcolo esterno delle proposte, l'architettura \textbf{Faster R-CNN} (proposta da Ren et al. nel 2015) introduce un paradigma innovativo basato sul principio di far generare le proposte direttamente alla rete convoluzionale stessa:

\[
\text{\emph{"Make CNN do proposals!"}}
\]

Questo obiettivo viene raggiunto inserendo una rete dedicata, denominata \textbf{Region Proposal Network (RPN)}, che riceve in ingresso le caratteristiche estratte dalla mappa convoluzionale comune e predice le proposte direttamente da esse, condividendo i pesi convoluzionali con il resto del rivelatore.

Il sistema viene addestrato congiuntamente ottimizzando un totale di \textbf{4 funzioni di perdita} (\emph{losses}):
\begin{enumerate}
    \item Perdita della RPN per la classificazione binaria delle ancore in oggetto / non oggetto (\emph{RPN classify object / not object}).
    \item Perdita della RPN per la regressione delle coordinate dei bounding-box delle proposte (\emph{RPN regress box coordinates}).
    \item Perdita finale per la classificazione multiclasse degli oggetti (\emph{Final classification score}).
    \item Perdita finale per la regressione di precisione delle coordinate dei bounding-box (\emph{Final box coordinates}).
\end{enumerate}

% [IMMAGINE DEL FILE DI SLIDE CON NOME "73-95"] Schema dell'architettura Faster R-CNN con l'introduzione della Region Proposal Network (RPN) — Slide 9 / Pagina 9 del PDF

L'integrazione della RPN elimina la necessità di algoritmi esterni, portando il tempo di test per singola immagine a valori estremamente competitivi:
\begin{itemize}
    \item \textbf{R-CNN}: $49\,\text{s}$
    \item \textbf{SPP-Net}: $4.3\,\text{s}$
    \item \textbf{Fast R-CNN}: $2.3\,\text{s}$
    \item \textbf{Faster R-CNN}: \hl{$0.2\,\text{s}$}
\end{itemize}

% [IMMAGINE DEL FILE DI SLIDE CON NOME "73-95"] Grafico comparativo delle prestazioni in termini di tempo di test inclusa Faster R-CNN — Slide 10 / Pagina 10 del PDF

\newpage

\section{Rilevamento Monostadio Senza Proposte: YOLO E SSD}

Una filosofia alternativa al rilevamento basato su proposte (a due stadi) è rappresentata dai modelli monostadio, come \textbf{YOLO} (\emph{You Only Look Once}) e \textbf{SSD} (\emph{Single-Shot MultiBox Detector}). L'obiettivo di questi approcci è mappare l'immagine direttamente in un tensore di output contenente sia le coordinate spaziali che le classi degli oggetti, utilizzando un'unica grande rete convoluzionale.

Il funzionamento logico di questi modelli si articola come segue:
\begin{itemize}
    \item Si considera un'immagine di input di dimensione $3 \times H \times W$.
    \item L'immagine viene suddivisa spazialmente in una griglia, ad esempio di dimensione $7 \times 7$.
    \item Per ogni cella della griglia, si assume la presenza di un insieme di \textbf{base boxes} (scatole di riferimento o ancore) centrate in quella cella. Nel caso illustrato a Slide 11, il numero di base boxes per cella è $B = 3$.
\end{itemize}

All'interno di ciascuna cella della griglia, la rete esegue le seguenti predizioni:
\begin{itemize}
    \item Per ciascuna delle $B$ base boxes, effettua la regressione verso un bounding-box finale stimando $5$ numeri reali:
    \[
    (dx, dy, dh, dw, \text{confidence})
    \]
    dove $(dx, dy)$ rappresenta l'offset spaziale del centro, $(dh, dw)$ definisce la variazione di scala di altezza e larghezza, e $\text{confidence}$ esprime il grado di confidenza del modello sulla presenza di un oggetto all'interno di quella scatola.
    \item Predice i punteggi di probabilità per ciascuna delle $C$ classi totali considerate, includendo la classe di background come categoria indipendente.
\end{itemize}

% [IMMAGINE DEL FILE DI SLIDE CON NOME "73-95"] Schema di funzionamento del metodo YOLO/SSD basato su griglia spaziale e base boxes — Slide 11 / Pagina 11 del PDF

La dimensione del tensore di output finale della rete convoluzionale è definita in base alla discretizzazione spaziale della griglia, al numero di ancore e al numero di classi:

\[
\boxed{7 \times 7 \times (5 \times B + C)}
\]

Questo approccio consente di effettuare il rilevamento tramite un singolo passo in avanti della rete, riducendo drasticamente la complessità computazionale rispetto alle architetture basate su proposte di regioni.

% [IMMAGINE DEL FILE DI SLIDE CON NOME "73-95"] Schema di YOLO/SSD che mostra il passaggio diretto dall'immagine di input al tensore dei punteggi finali — Slide 12 / Pagina 12 del PDF

\newpage

\section{Fattori Di Progettazione E Trade-Off Nel Rilevamento Di Oggetti}

La progettazione di un sistema di object detection richiede la selezione e la configurazione di molteplici componenti strutturali, che influenzano direttamente il bilanciamento tra accuratezza e velocità del modello.

Le variabili principali da considerare includono:
\begin{itemize}
    \item \textbf{Base Network} (Estrattore di caratteristiche convoluzionali): la scelta spazia tra reti ad alta capacità computazionale o ottimizzate per dispositivi mobili, come VGG16, ResNet-101, Inception V2, Inception V3, Inception ResNet e MobileNet.
    \item \textbf{Object Detection Architecture}: la scelta dell'architettura complessiva, come Faster R-CNN, R-FCN o SSD.
    \item \textbf{Image Size}: la risoluzione spaziale delle immagini fornite in input.
    \item \textbf{\# Region Proposals}: il numero di proposte di regioni generate ed elaborate.
\end{itemize}

Dalle analisi condotte in letteratura (ad esempio, Huang et al., CVPR 2017) emergono i seguenti concetti chiave (\emph{takeaways}):
\begin{itemize}
    \item I modelli basati su \textbf{Faster R-CNN} si collocano nella fascia ad alta precisione, risultando complessivamente più lenti ma maggiormente accurati.
    \item I modelli basati su \textbf{SSD} sono significativamente più veloci nell'elaborazione, ma mostrano un'accuratezza inferiore.
\end{itemize}

% [IMMAGINE DEL FILE DI SLIDE CON NOME "73-95"] Panoramica delle variabili di progettazione e dei trade-off velocità/accuratezza per i moderni rilevatori di oggetti — Slide 13 / Pagina 13 del PDF

\newpage

\section{Oltre Il Rilevamento Di Oggetti: Dense Captioning}

Un'estensione naturale che unisce la visione artificiale e l'elaborazione del linguaggio naturale è costituita dal compito di \textbf{Dense Captioning} (descritta come l'unione tra rilevamento di oggetti e generazione di didascalie).

A differenza della classica didascalizzazione globale di un'immagine, il Dense Captioning mira a individuare molteplici regioni significative all'interno di una scena e a generare per ciascuna di esse una descrizione testuale dettagliata e localizzata.

% [IMMAGINE DEL FILE DI SLIDE CON NOME "73-95"] Esempio di output del task di Dense Captioning applicato a scene complesse — Slide 14 / Pagina 14 del PDF

L'architettura di riferimento, denominata \textbf{DenseCap} (proposta da Johnson, Karpathy e Fei-Fei nel 2016), si articola nei seguenti passaggi logici:
\begin{enumerate}
    \item L'immagine di input di dimensione $3 \times W \times H$ viene elaborata da una \textbf{CNN} per estrarre una mappa di caratteristiche convoluzionali di dimensione $C \times W' \times H'$.
    \item La mappa viene fornita a un livello di localizzazione (\textbf{Localization Layer}) in cui:
    \begin{itemize}
        \item Un livello convoluzionale (\emph{Conv}) genera simultaneamente le proposte di regione ($4k \times W' \times H'$) e i rispettivi punteggi associati ($k \times W' \times H'$).
        \item Tramite un'operazione di campionamento (\emph{Sampling}), si selezionano le $B$ migliori proposte di regione (dimensione $B \times 4$).
        \item Un generatore di griglia (\textbf{Grid Generator}) definisce una griglia di campionamento di dimensione $B \times X \times Y \times 2$.
        \item Un campionatore bilineare (\textbf{Bilinear Sampler}) estrae le caratteristiche corrispondenti alle regioni selezionate, producendo un volume normalizzato di caratteristiche di regione pari a $B \times 512 \times 7 \times 7$.
    \end{itemize}
    \item Tali caratteristiche di regione (di dimensione generica $B \times C \times X \times Y$) vengono fornite a una rete di riconoscimento (\textbf{Recognition Network}), che produce codici vettoriali compatti per ciascuna regione di dimensione $B \times D$.
    \item Infine, una rete ricorrente di tipo \textbf{LSTM} riceve questi vettori e genera sequenzialmente le didascalie testuali per ciascuna regione identificata.
\end{enumerate}

% [IMMAGINE DEL FILE DI SLIDE CON NOME "73-95"] Architettura dettagliata della rete DenseCap per il Dense Captioning — Slide 15 / Pagina 15 del PDF

% [IMMAGINE DEL FILE DI SLIDE CON NOME "73-95"] Esempio pratico di Dense Captioning su un fotogramma d'interni — Slide 16 / Pagina 16 del PDF

\newpage

\section{Segmentazione Delle Istanze E Mask R-CNN}

Per comprendere il compito di \textbf{Segmentazione delle Istanze} (\emph{Instance Segmentation}), è opportuno confrontarlo con gli altri principali task di percezione visiva:
\begin{itemize}
    \item \textbf{Segmentazione Semantica} (\emph{Semantic Segmentation}): etichetta ogni singolo pixel dell'immagine con la classe di appartenenza (es. erba, gatto, albero, cielo), senza distinguere tra diverse istanze della medesima classe. Non vi è la nozione di singolo oggetto (\emph{"No objects, just pixels"}).
    \item \textbf{Classificazione + Localizzazione}: assume la presenza di un singolo oggetto nell'immagine, predicendone la classe e delimitandolo con un singolo bounding-box.
    \item \textbf{Rilevamento di Oggetti} (\emph{Object Detection}): individua ed evidenzia tramite bounding-box individuali tutte le istanze degli oggetti presenti nell'immagine, anche se appartenenti alla stessa categoria (es. cane, cane, gatto).
    \item \textbf{Segmentazione delle Istanze} (\emph{Instance Segmentation}): unisce i vantaggi del rilevamento degli oggetti e della segmentazione semantica, identificando ciascuna istanza singola e predicendo per essa una maschera di segmentazione precisa a livello di pixel.
\end{itemize}

% [IMMAGINE DEL FILE DI SLIDE CON NOME "73-95"] Confronto visivo tra i task di Segmentazione Semantica, Classificazione + Localizzazione, Rilevamento di Oggetti e Segmentazione delle Istanze — Slide 17 / Pagina 17 del PDF

L'architettura di riferimento per questo compito è \textbf{Mask R-CNN} (He et al., 2017), la quale estende Faster R-CNN aggiungendo un ramo parallelo dedicato alla predizione delle maschere di segmentazione.

Rispetto alle architetture precedenti, Mask R-CNN introduce il modulo \textbf{RoI Align} in sostituzione del RoI Pooling. Il modulo RoI Align evita la quantizzazione delle coordinate delle RoI preservando la precisione spaziale delle mappe di caratteristiche, che vengono estratte con dimensioni fisse pari a:

\[
256 \times 14 \times 14
\]

Da questo punto, l'elaborazione si divide in due rami paralleli:
\begin{enumerate}
    \item Un ramo deputato al calcolo dei punteggi di classificazione per le $C$ classi e alla definizione delle coordinate spaziali dei bounding-box per ciascuna classe ($4 \times C$).
    \item Un ramo convoluzionale costituito da strati di convoluzione consecutivi che elaborano il tensore di dimensione $256 \times 14 \times 14$ per predire una maschera binaria di dimensione:
    \[
    C \times 14 \times 14
    \]
    ovvero una maschera di pixel $14 \times 14$ per ciascuna delle $C$ classi di oggetti considerate.
\end{enumerate}

% [IMMAGINE DEL FILE DI SLIDE CON NOME "73-95"] Architettura di rete di Mask R-CNN con il modulo RoI Align e la predizione della maschera — Slide 18 / Pagina 18 del PDF

L'applicazione di questo metodo consente di ottenere maschere di pixel estremamente precise anche in contesti con occlusioni o elevata densità di oggetti.

% [IMMAGINE DEL FILE DI SLIDE CON NOME "73-95"] Esempi qualitativi di risultati ottenuti con l'architettura Mask R-CNN — Slide 19 / Pagina 19 del PDF

\newpage

\subsection{Estensione Alla Stima Della Posa Humana}

Oltre alla segmentazione delle istanze, Mask R-CNN si rivela un'architettura flessibile in grado di eseguire anche la \textbf{stima della posa umana} (\emph{human pose estimation}).

Per implementare questa funzionalità, l'architettura viene integrata inserendo nel primo ramo la predizione delle \textbf{coordinate dei giunti} (\emph{joint coordinates}) o punti chiave del corpo umano (\emph{keypoints}). La rete tratta ciascun giunto corporeo come una categoria a sé stante e modella la predizione come una maschera di localizzazione di tipo \emph{one-hot}, consentendo di ricostruire lo scheletro dei soggetti individuati all'interno dell'immagine.

% [IMMAGINE DEL FILE DI SLIDE CON NOME "73-95"] Architettura di Mask R-CNN estesa alla predizione della posa mediante le coordinate dei giunti — Slide 20 / Pagina 20 del PDF

% [IMMAGINE DEL FILE DI SLIDE CON NOME "73-95"] Risultati visivi della stima della posa umana ottenuti con Mask R-CNN — Slide 21 / Pagina 21 del PDF

\newpage

\section{Riepilogo Dei Task Di Visione Artificiale}

In sintesi, i quattro principali compiti di percezione visiva analizzati possono essere classificati in base alla granularità spaziale e alla tipologia di output richiesto:

\begin{table}[H]
\centering
\begin{tabularx}{\textwidth}{lXX}
\toprule
\textbf{Task} & \textbf{Granularità dell'Output} & \textbf{Distinzione delle Istanze} \\
\midrule
\textbf{Segmentazione Semantica} & Singoli pixel dell'immagine & No (tutte le istanze della stessa classe condividono l'etichetta) \\
\textbf{Classificazione + Localizzazione} & Singolo bounding-box globale & Singolo oggetto predominante nell'immagine \\
\textbf{Rilevamento di Oggetti} & Molteplici bounding-box & Sì (ogni oggetto ha una propria etichetta e scatola di delimitazione) \\
\textbf{Segmentazione delle Istanze} & Maschere di pixel per ciascun oggetto & Sì (ogni oggetto ha una maschera e confini spaziali definiti) \\
\bottomrule
\end{tabularx}
\caption{Confronto riassuntivo delle caratteristiche dei principali compiti di visione artificiale.}
\label{tab:recap_tasks}
\end{table}

% [IMMAGINE DEL FILE DI SLIDE CON NOME "73-95"] Schema riassuntivo dei diversi compiti di visione artificiale — Slide 22 / Pagina 22 del PDF

% AA Slides 11.5
\section{Region Of Interest Pooling (RoI)}

Le moderne architetture dedicate alla rilevazione degli oggetti (\emph{object detection}) si strutturano tipicamente su due fasi principali:

\begin{itemize}
    \item \textbf{Proposta delle regioni (Region proposal)}: Data un'immagine in input, l'obiettivo è individuare tutte le possibili aree all'interno delle quali potrebbero essere collocati degli oggetti. L'output di questa fase consiste in un elenco di scatole di delimitazione (\emph{bounding boxes}) che indicano le posizioni più probabili degli oggetti stessi. Queste aree vengono comunemente denominate proposte di regione (\emph{region proposals}) o regioni di interesse (\textbf{Region of Interest}, o \textbf{RoI}).

    \item \textbf{Classificazione finale (Final classification)}: Per ciascuna proposta di regione generata nella fase precedente, il sistema deve stabilire se l'area appartenga a una delle classi di interesse (\emph{target classes}) oppure allo sfondo (\emph{background}).
\end{itemize}

% [IMMAGINE DEL FILE DI SLIDE CON NOME "ROI E YOLO"] Schema del flusso di elaborazione del Region of Interest Pooling (RoI Pooling) — Slide 2 / Pagina 2 del PDF

Nella fase di proposta delle regioni, è fondamentale generare un quantitativo elevato di regioni di interesse. Questo requisito è essenziale poiché, se un oggetto non viene correttamente individuato durante il primo stadio, non vi sarà alcuna possibilità di classificarlo nel secondo. Tuttavia, la gestione di un numero così cospicuo di proposte solleva diverse problematiche:

\begin{itemize}
    \item La generazione e l'elaborazione di un numero elevato di regioni di interesse comporta un significativo carico computazionale, con conseguenti problemi di prestazioni.
    \item Tale approccio risulta subottimale dal punto di vista della velocità di elaborazione (\emph{processing speed}).
    \item L'addestramento dell'intera rete non può avvenire in modalità \emph{end-to-end}.
\end{itemize}

Per ovviare a queste limitazioni e ottimizzare il flusso di lavoro, si introduce lo strato di \textbf{RoI Pooling} (\emph{RoI pooling layer}).

\newpage

\subsection{Funzionamento Dello Strato Di RoI Pooling}

Lo strato di RoI pooling accetta due input principali:

\begin{enumerate}
    \item Una mappa delle feature (\emph{feature map}) a dimensione fissa, ottenuta mediante una rete convoluzionale profonda composta da diversi strati di convoluzione e di \emph{max pooling}.
    \item Una matrice di dimensioni $N \times 5$ che rappresenta l'elenco delle regioni di interesse rilevate, in cui $N$ indica il numero di RoI. La prima colonna di questa matrice specifica l'indice dell'immagine di riferimento, mentre le restanti quattro colonne contengono le coordinate spaziali dei vertici superiore sinistro e inferiore destro di ciascuna regione.
\end{enumerate}

% [IMMAGINE DEL FILE DI SLIDE CON NOME "ROI E YOLO"] Esempio di un'immagine raffigurante un gatto con numerose bounding boxes candidate — Slide 4 / Pagina 4 del PDF

L'operazione di RoI pooling agisce estraendo, per ciascuna regione di interesse presente nella lista di input, la sezione corrispondente della mappa delle feature di input. Questa sezione viene successivamente riscalata a una dimensione d'uscita predefinita (ad esempio, $7 \times 7$). Il processo di riscalamento avviene secondo i seguenti passaggi:

\begin{itemize}
    \item Suddivisione della proposta di regione in sezioni di uguale dimensione (il numero di tali sezioni deve essere identico alla dimensione spaziale dell'output desiderato).
    \item Ricerca del valore massimo (\emph{max value}) all'interno di ciascuna sezione.
    \item Scrittura di questi valori massimi nel buffer di output.
\end{itemize}

L'aspetto chiave di questo metodo è che, a partire da un elenco di rettangoli di dimensioni tra loro eterogenee, si ottiene rapidamente un insieme di mappe di feature aventi tutte la medesima dimensione predefinita. La dimensione dell'output del RoI pooling non dipende in alcun modo dalla dimensione della mappa delle feature originale o dalle dimensioni iniziali delle proposte di regione.

Il beneficio principale risiede nell'incremento della velocità di calcolo. In presenza di svariate proposte di oggetti all'interno del medesimo fotogramma, è possibile riutilizzare la mappa delle feature di input precedentemente calcolata per tutte le RoI. Poiché il calcolo dei filtri convoluzionali nei primi stadi di una rete è estremamente costoso, questo meccanismo di condivisione consente di ridurre drasticamente i tempi di elaborazione.

\newpage

\subsection{Esempio Pratico Di Riscalamento Con RoI Pooling}

Per comprendere la logica matematica dell'operazione, consideriamo un esempio dettagliato in cui si desidera ridurre una proposta di regione a un output finale di dimensione $2 \times 2$ tramite max pooling.

% [IMMAGINE DEL FILE DI SLIDE CON NOME "ROI E YOLO"] Matrice 8x8 che rappresenta la mappa delle feature di input con valori numerici — Slide 6 / Pagina 6 del PDF

Si parta da una mappa delle feature di input di dimensione $8 \times 8$ contenente i valori numerici indicati nella Slide 6.

% [IMMAGINE DEL FILE DI SLIDE CON NOME "ROI E YOLO"] Mappa delle feature con una proposta di regione (region proposal) di dimensione 5x7 evidenziata — Slide 7 / Pagina 7 del PDF

Su questa mappa viene definita una proposta di regione (\emph{region proposal}) di dimensioni $5 \times 7$, che copre le righe da 4 a 8 e le colonne da 1 a 7, come mostrato nella Slide 7.

% [IMMAGINE DEL FILE DI SLIDE CON NOME "ROI E YOLO"] Suddivisione della proposta di regione in quattro sezioni di pooling di dimensioni asimmetriche — Slide 8 / Pagina 8 del PDF

Dato che la regione proposta ha dimensione $5 \times 7$ e la griglia di output desiderata deve essere $2 \times 2$, lo strato suddivide la RoI in quattro sezioni (\emph{pooling sections}) disposte su una griglia $2 \times 2$, come illustrato nella Slide 8. Poiché 5 e 7 non sono perfettamente divisibili per 2, la divisione avviene creando sezioni di dimensioni asimmetriche:
\begin{itemize}
    \item L'altezza di 5 righe viene suddivisa in una fascia superiore di 2 righe e una fascia inferiore di 3 righe.
    \item La larghezza di 7 colonne viene suddivisa in una fascia sinistra di 3 colonne e una fascia destra di 4 colonne.
\end{itemize}

Le quattro sotto-sezioni risultanti contengono i seguenti elementi:

\begin{enumerate}
    \item \textbf{Sezione Superiore Sinistra} (dimensione $2 \times 3$, righe 4--5, colonne 1--3):
    \[
    \begin{pmatrix}
    0.85 & 0.34 & 0.76 \\
    0.32 & 0.74 & 0.21
    \end{pmatrix}
    \]
    Il valore massimo di questa sezione è:
    \[
    \max(0.85, 0.34, 0.76, 0.32, 0.74, 0.21) = \boxed{0.85}
    \]

    \item \textbf{Sezione Superiore Destra} (dimensione $2 \times 4$, righe 4--5, colonne 4--7):
    \[
    \begin{pmatrix}
    0.84 & 0.29 & 0.75 & 0.62 \\
    0.39 & 0.34 & 0.03 & 0.33
    \end{pmatrix}
    \]
    Il valore massimo di questa sezione è:
    \[
    \max(0.84, 0.29, 0.75, 0.62, 0.39, 0.34, 0.03, 0.33) = \boxed{0.84}
    \]

    \item \textbf{Sezione Inferiore Sinistra} (dimensione $3 \times 3$, righe 6--8, colonne 1--3):
    \[
    \begin{pmatrix}
    0.20 & 0.14 & 0.16 \\
    0.19 & 0.69 & 0.09 \\
    0.83 & 0.24 & 0.97
    \end{pmatrix}
    \]
    Il valore massimo di questa sezione è:
    \[
    \max(0.20, 0.14, 0.16, 0.19, 0.69, 0.09, 0.83, 0.24, 0.97) = \boxed{0.97}
    \]

    \item \textbf{Sezione Inferiore Destra} (dimensione $3 \times 4$, righe 6--8, colonne 4--7):
    \[
    \begin{pmatrix}
    0.13 & 0.73 & 0.65 & 0.96 \\
    0.86 & 0.88 & 0.07 & 0.01 \\
    0.04 & 0.24 & 0.35 & 0.50
    \end{pmatrix}
    \]
    Il valore massimo di questa sezione è:
    \[
    \max(0.13, 0.73, 0.65, 0.96, 0.86, 0.88, 0.07, 0.01, 0.04, 0.24, 0.35, 0.50) = \boxed{0.96}
    \]
\end{enumerate}

% [IMMAGINE DEL FILE DI SLIDE CON NOME "ROI E YOLO"] Matrice di output 2x2 risultante dall'operazione di RoI pooling — Slide 9 / Pagina 9 del PDF

I quattro massimi calcolati vengono copiati nella matrice di output finale di dimensione $2 \times 2$, illustrata nella Slide 9:
\[
\bm{Y}_{\text{pool}} = \begin{pmatrix}
0.85 & 0.84 \\
0.97 & 0.96
\end{pmatrix}
\]

\newpage

\section{YOLO: You Only Look Once}

Prima di introdurre il paradigma YOLO, si consideri come termine di paragone la classica architettura a due stadi \textbf{R-CNN} (\emph{Regions with CNN features}), il cui flusso operativo è sintetizzato nella Slide 11.

% [IMMAGINE DEL FILE DI SLIDE CON NOME "ROI E YOLO"] Schema della pipeline classica di R-CNN: estrazione delle proposte, warping, calcolo delle feature ed SVM — Slide 11 / Pagina 11 del PDF

La procedura di R-CNN prevede quattro passaggi sequenziali:
\begin{enumerate}
    \item Ricezione dell'immagine di input.
    \item Estrazione di circa 2000 proposte di regione (\emph{region proposals}).
    \item Riscalamento (\emph{warping}) di ciascuna regione proposta a una dimensione standard e calcolo dei relativi vettori di feature tramite una rete convoluzionale (CNN).
    \item Classificazione di ciascuna regione tramite modelli dedicati (ad esempio, per verificare la presenza di classi quali \emph{aeroplane}, \emph{person}, \emph{tvmonitor}, ecc.).
\end{enumerate}

Questo metodo risulta intrinsecamente lento, poiché richiede l'elaborazione separata di migliaia di sotto-regioni. L'approccio \textbf{YOLO} propone invece una filosofia differente: l'intera immagine viene analizzata in un'unica operazione di calcolo.

% [IMMAGINE DEL FILE DI SLIDE CON NOME "ROI E YOLO"] Schema generale della pipeline di YOLO (suddivisione in griglia, predizione dei box e classificazione) — Slide 10 / Pagina 10 del PDF

Le caratteristiche distintive di YOLO includono:
\begin{itemize}
    \item \textbf{Elevata velocità di elaborazione (\emph{Very fast})}: La versione veloce della rete (\emph{Fast YOLO}) è in grado di operare a una frequenza di 155 fotogrammi per secondo (FPS) sul medesimo benchmark di Faster R-CNN.
    \item \textbf{Rilevatore a colpo singolo (\emph{Single Shot Detector - SSD})}: Il sistema considera l'intera immagine in un solo passaggio (\emph{single shot}) per effettuare l'intera rilevazione degli oggetti.
    \item \textbf{Superiore capacità di generalizzazione}: YOLO mostra ottime capacità di generalizzazione quando applicato a domini di immagini differenti rispetto a quelli impiegati nella fase di addestramento.
\end{itemize}

\newpage

\subsection{Procedura Di Rilevazione}

Al contrario dei sistemi basati su proposte di regione, YOLO elabora l'immagine intera producendo simultaneamente le etichette di classificazione, le coordinate delle scatole di delimitazione e le relative probabilità di confidenza.

% [IMMAGINE DEL FILE DI SLIDE CON NOME "ROI E YOLO"] Rappresentazione concettuale del funzionamento di YOLO che riceve un'immagine stradale ed emette le predizioni finali — Slide 13 / Pagina 13 del PDF

La procedura logica di rilevazione si sviluppa in tre fasi:

\begin{enumerate}
    \item Suddivisione dell'immagine in una griglia di dimensioni $S \times S$ celle (ad esempio, una griglia $3 \times 3$ come riportato nella Slide 14).
    \item Ciascuna cella della griglia predice un numero $B$ di scatole di delimitazione (nell'esempio, $B = 2$).
    \item Vengono mantenute e mostrate solo le bounding box per le quali il valore di confidenza risulta superiore a una soglia definita dall'utente.
\end{enumerate}

% [IMMAGINE DEL FILE DI SLIDE CON NOME "ROI E YOLO"] Fasi dettagliate della procedura di predizione di YOLO tramite griglia ed eliminazione dei box a bassa confidenza — Slide 14 / Pagina 14 del PDF

\subsection{Definizione Geometrica E Struttura Dei Dati}

Per ciascuna scatola di delimitazione predetta, YOLO stima i seguenti parametri:

\begin{itemize}
    \item La probabilità $C$ che la scatola contenga un oggetto, denominata \textbf{objectness}:
    \[
    C = P(\text{Object})
    \]
    \item Le coordinate spaziali del centro della bounding box, denotate con $(x, y)$ e calcolate in modo relativo rispetto ai confini della cella corrente.
    \item La larghezza $w$ e l'altezza $h$ della bounding box, espresse come percentuale riferita alle dimensioni della cella stessa.
    \item La probabilità condizionale che la cella contenga un oggetto appartenente a una delle $K$ classi possibili, nell'ipotesi che la cella contenga effettivamente un oggetto.
\end{itemize}

% [IMMAGINE DEL FILE DI SLIDE CON NOME "ROI E YOLO"] Schema esplicativo dei parametri geometrici della bounding box e del vettore delle predizioni — Slide 15 / Pagina 15 del PDF

Il vettore di output associato a ogni singola scatola di delimitazione ha una lunghezza pari a $5 + K$. Considerando l'intera griglia di predizione sull'immagine completa, la rete produce complessivamente un volume di stime pari a:
\[
\boxed{S \times S \times (5 \times B + K)}
\]

\newpage

\subsection{Lo Strato Di Output E L'Architettura Di Rete}

Si consideri uno scenario concreto caratterizzato dai seguenti parametri operativi: dimensione della griglia $S = 7$, numero di scatole per cella $B = 2$ e numero di classi target $K = 20$. In questo scenario specifico, lo strato finale di YOLO restituisce un tensore tridimensionale di dimensioni:
\[
\bm{Y} \in \mathbb{R}^{7 \times 7 \times (5 \times 2 + 20)} \implies \bm{Y} \in \mathbb{R}^{7 \times 7 \times 30}
\]

% [IMMAGINE DEL FILE DI SLIDE CON NOME "ROI E YOLO"] Struttura del tensore tridimensionale di output di YOLO di dimensione 7x7x30 — Slide 16 / Pagina 16 del PDF

Per ciascuna delle $7 \times 7$ celle della griglia, il vettore di output di lunghezza 30 risulta così strutturato:
\begin{itemize}
    \item Primi 5 elementi: parametri geometrici e confidenza della prima bounding box, ovvero $(c_1, x_1, y_1, w_1, h_1)$ con il relativo \emph{score 1}.
    \item Successivi 5 elementi: parametri geometrici e confidenza della seconda bounding box, ovvero $(c_2, x_2, y_2, w_2, h_2)$ con il relativo \emph{score 2}.
    \item Restanti 20 elementi: probabilità condizionali riferite alle 20 classi oggetto di rilevazione (ad esempio: \emph{aeroplane, bicycle, bird, boat, bottle, bus, car, cat, chair, cow, diningtable, dog, horse, motorbike, person, pottedplant, sheep, sofa, train, tvmonitor}).
\end{itemize}

L'architettura di rete utilizzata per elaborare l'immagine e generare tale output si compone di tre moduli:

\begin{enumerate}
    \item Un'ossatura di rete (\textbf{backbone}) basata su un modello \textbf{ResNet50}, accoppiata a un'operazione di convoluzione $1 \times 1$ finalizzata al controllo e alla riduzione della dimensione spaziale delle mappe delle feature.
    \item Uno strato completamente connesso (\textbf{Fully connected layer}).
    \item Un'operazione di riorganizzazione dei dati (\textbf{Reshape}) per convertire il vettore unidimensionale di output nel tensore tridimensionale finale di dimensioni $(S, S, 5 \times B + K)$.
\end{enumerate}

% [IMMAGINE DEL FILE DI SLIDE CON NOME "ROI E YOLO"] Flusso architetturale della rete YOLO basata su backbone ResNet50 e blocco fully connected — Slide 17 / Pagina 17 del PDF

Il flusso architetturale della rete ResNet50 mostrato in figura si sviluppa secondo la sequenza:
\[
\text{Input} \to \text{Zero Padding} \to \text{CONV} \to \text{Batch Norm} \to \text{ReLU} \to \text{Max Pool} \to \text{Stage 1 a 5} \to \text{Avg Pool} \to \text{Flattening} \to \text{FC} \to \text{Output}
\]
Ciascuno stadio di elaborazione (Stage da 2 a 5) è strutturato alternando blocchi di tipo convoluzionale (\emph{Conv Block}) e blocchi di tipo identità (\emph{ID Block}).

\newpage

\subsection{La Funzione Di Perdita Multi-Task}

La funzione di perdita (\emph{Loss Function}) di YOLO è formulata per ottimizzare in modo simultaneo la precisione di localizzazione delle scatole di delimitazione, la stima di presenza dell'oggetto e l'errore di classificazione semantica. Essa si basa su un approccio a errore quadratico medio pesato ed è definita dalla seguente equazione complessa:

\begin{align*}
\mathcal{L} &= \underbrace{\lambda_{\text{coord}} \sum_{i=0}^{S^2} \sum_{j=0}^{B} \mathbb{1}_{ij}^{\text{obj}} \left[ (x_i - \hat{x}_i)^2 + (y_i - \hat{y}_i)^2 \right]}_{\text{(1) Errore sulle coordinate del centro}} \\
&+ \underbrace{\lambda_{\text{coord}} \sum_{i=0}^{S^2} \sum_{j=0}^{B} \mathbb{1}_{ij}^{\text{obj}} \left[ \left(\sqrt{w_i} - \sqrt{\hat{w}_i}\right)^2 + \left(\sqrt{h_i} - \sqrt{\hat{h}_i}\right)^2 \right]}_{\text{(2) Errore su larghezza e altezza della bounding box}} \\
&+ \underbrace{\sum_{i=0}^{S^2} \sum_{j=0}^{B} \mathbb{1}_{ij}^{\text{obj}} (C_i - \hat{C}_i)^2}_{\text{(3) Errore di confidenza per celle con oggetto}} \\
&+ \underbrace{\lambda_{\text{noobj}} \sum_{i=0}^{S^2} \sum_{j=0}^{B} \mathbb{1}_{ij}^{\text{noobj}} (C_i - \hat{C}_i)^2}_{\text{(4) Errore di confidenza per celle senza oggetto}} \\
&+ \underbrace{\sum_{i=0}^{S^2} \mathbb{1}_{i}^{\text{obj}} \sum_{c \in \text{classes}} (p_i(c) - \hat{p}_i(c))^2}_{\text{(5) Errore di classificazione}}
\end{align*}

Si analizzino singolarmente le cinque componenti della formula per comprenderne la motivazione teorica e il funzionamento:

\begin{enumerate}
    \item \textbf{Errore sulle coordinate del centro (1)}:
    Questo termine quantifica l'errore quadratico medio relativo allo scostamento tra la posizione reale del centro del box $(x_i, y_i)$ e quella stimata dalla rete $(\hat{x}_i, \hat{y}_i)$. La funzione indicatrice $\mathbb{1}_{ij}^{\text{obj}}$ assume valore pari a 1 se la scatola di delimitazione $j$ all'interno della cella $i$ è considerata responsabile per la predizione di quell'oggetto, altrimenti vale 0. L'iperparametro $\lambda_{\text{coord}}$ è un fattore di scala impostato a un valore elevato (tipicamente $\lambda_{\text{coord}} = 5$) per dare maggiore peso e priorità alla precisione geometrica della localizzazione rispetto alla pura classificazione delle classi.

    \item \textbf{Errore sulle dimensioni della bounding box (2)}:
    Questa componente calcola l'errore associato alle dimensioni della scatola di delimitazione (larghezza $w$ ed altezza $h$). L'impiego della radice quadrata $(\sqrt{w_i}, \sqrt{h_i})$ rispetto ai valori assoluti è introdotto per differenziare l'impatto degli errori di predizione a seconda della scala: in questo modo, uno scostamento di pochi pixel su una scatola di grandi dimensioni viene penalizzato meno rispetto al medesimo scostamento commesso su una scatola di piccole dimensioni.

    \item \textbf{Errore di confidenza per celle contenenti un oggetto (3)}:
    Questa sezione della perdita valuta l'accuratezza del punteggio di confidenza $\hat{C}_i$ stimato dalla rete rispetto al valore di confidenza reale $C_i$ (che è pari a 1 se l'oggetto è presente nella cella). La penalità viene applicata solo se la bounding box $j$ della cella $i$ è responsabile per l'oggetto ($\mathbb{1}_{ij}^{\text{obj}} = 1$).

    \item \textbf{Errore di confidenza per celle prive di oggetti (4)}:
    Questa componente penalizza la confidenza stimata $\hat{C}_i$ qualora la cella della griglia non contenga in realtà alcun oggetto ($C_i = 0$, con indicatore di assenza dell'oggetto $\mathbb{1}_{ij}^{\text{noobj}} = 1$). Poiché nella maggior parte delle immagini reali la maggior parte delle celle della griglia non contiene oggetti, il contributo di errore derivante dallo sfondo tenderebbe a sovrastare quello associato agli oggetti reali. Per scongiurare questo fenomeno di instabilità durante l'addestramento, si introduce il fattore di scala riduttivo $\lambda_{\text{noobj}}$ (tipicamente impostato a un valore basso, come $\lambda_{\text{noobj}} = 0.5$).

    \item \textbf{Errore di classificazione (5)}:
    L'ultimo termine misura l'errore quadratico medio associato alle probabilità condizionali di appartenenza alle classi $p_i(c)$ rispetto a quelle stimate $\hat{p}_i(c)$. Tale errore viene computato esclusivamente per le celle all'interno delle quali è accertata la presenza di un oggetto ($\mathbb{1}_{i}^{\text{obj}} = 1$).
\end{enumerate}

\end{document}
